% Los medios y la información responsable — Espagnol 1ère A — Cours signé OnBuch+
% Programme officiel MINESEC. Compiler avec : tectonic EA17-medios-informacion-responsable.tex
\newcommand{\DOCMATIERE}{Espagnol}
\newcommand{\DOCNIVEAU}{1ère A}
\newcommand{\DOCMODULE}{Módulo 5 : Les mass médias et la communication}
\newcommand{\DOCLECON}{EA17}
\newcommand{\DOCTITRE}{Los medios y la información responsable}
% OnBuch+ — préambule « cours signés » ENRICHI (compilé avec tectonic / XeLaTeX)
% Système visuel « Pop » d'OnBuch+ : fond crème (jamais blanc), bordures encre
% épaisses, ombres « dures » décalées, titres Archivo Black en capitales.
% Version MATHÉMATIQUES : graphes (pgfplots/tikz), boîtes pédagogiques
% (définition, propriété, méthode, attention, le savais-tu, exercice résolu,
% exercice, TP, bilan), page de garde et signature OnBuch+ ancrée en bas.
%
% Macros attendues dans main.tex AVANT \input{preamble} :
%   \DOCMATIERE  \DOCNIVEAU  \DOCMODULE  \DOCLECON (numéro)  \DOCTITRE
\documentclass[11pt]{article}

\usepackage{fontspec}
\usepackage[spanish,french]{babel} % français = langue principale ; l'espagnol via \esp{..} / espagnol
\usepackage[a4paper,margin=2cm,top=3.4cm,bottom=2.8cm,headheight=1.4cm,footskip=1.3cm]{geometry}
\usepackage{amsmath,amssymb,amsfonts}
\usepackage{mathtools} % \xmapsto, \coloneqq... souvent utilisés par les modèles
\usepackage{cancel} % simplifications barrées (bilans de forces)
% \abs{z} (module, valeur absolue) et \norm{u} (norme), utilisés par les modèles
\providecommand{\abs}{}\let\abs\relax\DeclarePairedDelimiter{\abs}{\lvert}{\rvert}
\providecommand{\norm}{}\let\norm\relax\DeclarePairedDelimiter{\norm}{\lVert}{\rVert}
\usepackage[table]{xcolor} % [table] : fournit aussi \rowcolors (alternance de lignes)
\usepackage{pagecolor}
\usepackage{tikz}
\usetikzlibrary{arrows.meta,positioning,calc,shapes.geometric,shapes.misc,shapes.arrows,decorations.pathmorphing,decorations.pathreplacing,decorations.markings,patterns,shadows,mindmap,trees,fit,backgrounds,matrix,intersections,angles,quotes}
\usepackage{pgfplots}
\usepackage[version=4]{mhchem} % équations nucléaires \ce{^{235}_{92}U}
\usepackage[european,siunitx]{circuitikz} % schémas électriques
\pgfplotsset{compat=1.18}
\usepgfplotslibrary{fillbetween,groupplots} % aires entre courbes (calcul intégral)
\usepackage{enumitem}
\usepackage{fancyhdr}
% [most] charge toutes les bibliothèques tcolorbox (listings, minted, raster,
% documentation...) et gonfle inutilement la mémoire TeX sur un document long
% (~300 boîtes) : on ne charge que ce qui est réellement utilisé.
\usepackage[skins,breakable]{tcolorbox}
\usepackage{graphicx}
\usepackage{colortbl}
\usepackage{booktabs}
\usepackage{tabularx}
\usepackage{diagbox} % case d'en-tête diagonale des tableaux à double entrée
% Colonnes extensibles centrée / gauche / droite (C, L, R), souvent utilisées par les modèles
\newcolumntype{C}{>{\centering\arraybackslash}X}
\newcolumntype{L}{>{\raggedright\arraybackslash}X}
\newcolumntype{R}{>{\raggedleft\arraybackslash}X}
\usepackage{array}
\usepackage{multicol}
\usepackage{siunitx}
% parse-numbers=false : accepte aussi \SI{2,5 \times 10^{-3}}{...}
\sisetup{locale=FR,output-decimal-marker={,},group-minimum-digits=4,parse-numbers=false}
\DeclareSIUnit\ohm{\ensuremath{\symup{\Omega}}} % U+2126 absent de la police
\DeclareSIUnit\division{div} % réglages d'oscilloscope (ms/div, V/div)
\DeclareSIUnit\tr{tr} % tours (vitesse de rotation, ex. \SI{1500}{\tr\per\minute})
\DeclareSIUnit\molar{M} % concentration molaire (ex. \SI{3}{\molar}, \SI{1}{\micro\molar})
\DeclareSIUnit\an{an} % année (le modèle confond parfois avec \year, un compteur TeX) — ex. \SI{5}{\kilo\meter\per\an}
\DeclareSIUnit\FCFA{FCFA} % franc CFA (ex. \SI{500}{\FCFA\per\kilo\gram})
\DeclareSIUnit\degreeBrix{\textdegree Brix} % teneur en sucre d'un sirop/jus (réfractométrie)
\DeclareSIUnit\month{mois} % le modèle utilise parfois \month, un compteur TeX, comme unité de durée
\usepackage{qrcode}
% \degree : commande fréquemment utilisée par les modèles pour un angle ou une
% température en dehors de \SI{}{} (non fournie par les paquets chargés).
\providecommand{\degree}{^{\circ}}
% \qed (fin de démonstration), utilisé par les modèles sans amsthm
\providecommand{\qed}{\hfill$\square$}
% \barre : double barre des valeurs interdites dans un tableau de variations (tkz-tab)
\providecommand{\barre}{\ensuremath{\|}}
% environnement proof (amsthm non chargé) utilisé par les modèles
\ifdefined\proof\else\newenvironment{proof}[1][Démonstration]{\par\noindent\textit{#1.}\ }{\qed\par}\fi
\usepackage{lastpage}
\usepackage{hyperref}
\hypersetup{hidelinks}

% ============================================================
% Palette « Pop » OnBuch+ — mêmes hex que la classe P (plus_ui.dart)
% ============================================================
\definecolor{poporange}{HTML}{F59321}
\definecolor{popdark}{HTML}{E07A0C}
\definecolor{popcream}{HTML}{FFF6E8}
\definecolor{popink}{HTML}{1C1714}
\definecolor{poppurple}{HTML}{7A5AE0}
\definecolor{popgreen}{HTML}{2E9E6A}
\definecolor{poppink}{HTML}{E0457B}
\definecolor{popblue}{HTML}{2F7DE1}
\definecolor{popgold}{HTML}{C98A12}
\definecolor{popmuted}{HTML}{8A7F76}
\definecolor{popline}{HTML}{EADFD2}
\definecolor{popgray}{HTML}{8A7F76}
\colorlet{popgrey}{popgray}
% Alias tolérants (le modèle nomme parfois les couleurs différemment)
\colorlet{popwhite}{white}
\colorlet{popblack}{popink}
\colorlet{popred}{poppink}
\colorlet{popyellow}{popgold}
% Teintes claires (fonds de tableaux/encadrés) — le modèle nomme parfois
% les couleurs avec un suffixe L (ex. popblueL) sans les avoir définies.
\colorlet{poporangeL}{poporange!20}
\colorlet{popdarkL}{popdark!20}
\colorlet{poppurpleL}{poppurple!20}
\colorlet{popgreenL}{popgreen!20}
\colorlet{poppinkL}{poppink!20}
\colorlet{popblueL}{popblue!20}
\colorlet{popgoldL}{popgold!20}
\colorlet{popredL}{popred!20}
\colorlet{popgrayL}{popgray!20}
% Teintes claires pour fonds de tableaux / zones
\colorlet{popblueL}{popblue!10!white}
\colorlet{poporangeL}{poporange!14!white}
\colorlet{popgreenL}{popgreen!12!white}
\colorlet{poppurpleL}{poppurple!12!white}
\colorlet{poppinkL}{poppink!12!white}
\colorlet{popgoldL}{popgold!14!white}

\pagecolor{popcream}

% ============================================================
% Typographie
% ============================================================
\defaultfontfeatures{Ligatures=TeX}
\setmainfont{PlusJakartaSans-Medium.ttf}[Path=../../../fonts/,BoldFont=PlusJakartaSans-Bold.ttf]
% Maths en sans-serif (Fira Math) avec chiffres et lettres droites de Plus
% Jakarta, pour que les formules se fondent dans le texte.
\usepackage[math-style=ISO,bold-style=ISO]{unicode-math}
\setmathfont{FiraMath-Regular.otf}[Path=../../../fonts/]
\setmathfont{PlusJakartaSans-Medium.ttf}[Path=../../../fonts/,range={up/num,up/{Latin,latin},\mathup}]
\usepackage{icomma} % virgule décimale sans espace en mode math
% Caractères mathématiques Unicode absents des polices de texte (Plus Jakarta,
% Archivo Black) : rendus par leur équivalent mathématique, en texte comme en
% formule — sinon ils s'affichent en carré vide (ex. « Arithmétique dans ℤ »).
\usepackage{newunicodechar}
\newunicodechar{ℝ}{\ensuremath{\mathbb{R}}}
\newunicodechar{ℤ}{\ensuremath{\mathbb{Z}}}
\newunicodechar{ℂ}{\ensuremath{\mathbb{C}}}
\newunicodechar{ℕ}{\ensuremath{\mathbb{N}}}
\newunicodechar{ℚ}{\ensuremath{\mathbb{Q}}}
\newunicodechar{𝔻}{\ensuremath{\mathbb{D}}}
\newunicodechar{α}{\ensuremath{\alpha}}
\newunicodechar{β}{\ensuremath{\beta}}
\newunicodechar{γ}{\ensuremath{\gamma}}
\newunicodechar{δ}{\ensuremath{\delta}}
\newunicodechar{ε}{\ensuremath{\varepsilon}}
\newunicodechar{θ}{\ensuremath{\theta}}
\newunicodechar{λ}{\ensuremath{\lambda}}
\newunicodechar{μ}{\ensuremath{\mu}}
\newunicodechar{σ}{\ensuremath{\sigma}}
\newunicodechar{τ}{\ensuremath{\tau}}
\newunicodechar{φ}{\ensuremath{\varphi}}
\newunicodechar{ψ}{\ensuremath{\psi}}
\newunicodechar{ω}{\ensuremath{\omega}}
\newunicodechar{Σ}{\ensuremath{\Sigma}}
% majuscules produites par \MakeUppercase dans les titres (α-aminés -> Α)
\newunicodechar{Α}{\ensuremath{\alpha}}
\newunicodechar{Β}{\ensuremath{\beta}}
\newunicodechar{Γ}{\ensuremath{\Gamma}}
\newunicodechar{Δ}{\ensuremath{\Delta}}
\newunicodechar{Ε}{\ensuremath{\varepsilon}}
\newunicodechar{Θ}{\ensuremath{\Theta}}
\newunicodechar{Λ}{\ensuremath{\Lambda}}
\newunicodechar{Μ}{\ensuremath{\mu}}
\newunicodechar{Τ}{\ensuremath{\tau}}
\newunicodechar{Φ}{\ensuremath{\Phi}}
\newunicodechar{Ψ}{\ensuremath{\Psi}}
\newunicodechar{Ω}{\ensuremath{\Omega}}
\newunicodechar{↦}{\ensuremath{\mapsto}}
\newunicodechar{∈}{\ensuremath{\in}}
\newunicodechar{∉}{\ensuremath{\notin}}
\newunicodechar{∧}{\ensuremath{\wedge}}
\newunicodechar{⇔}{\ensuremath{\Leftrightarrow}}
\newunicodechar{⟺}{\ensuremath{\Longleftrightarrow}}
\newunicodechar{⇒}{\ensuremath{\Rightarrow}}
\newunicodechar{⟹}{\ensuremath{\Longrightarrow}}
\newunicodechar{∘}{\ensuremath{\circ}}
\newunicodechar{∪}{\ensuremath{\cup}}
\newunicodechar{∩}{\ensuremath{\cap}}
\newunicodechar{≡}{\ensuremath{\equiv}}
\newunicodechar{⊂}{\ensuremath{\subset}}
\newunicodechar{⊥}{\ensuremath{\perp}}
\newunicodechar{∃}{\ensuremath{\exists}}
\newunicodechar{∀}{\ensuremath{\forall}}
\newunicodechar{∛}{\ensuremath{\sqrt[3]{\,}}}
\newunicodechar{∜}{\ensuremath{\sqrt[4]{\,}}}
\newunicodechar{⃗}{\ensuremath{{}^{\to}}}
\newunicodechar{ⁿ}{\ifmmode^{n}\else\textsuperscript{\ensuremath{n}}\fi}
\newunicodechar{ˣ}{\ifmmode^{x}\else\textsuperscript{\ensuremath{x}}\fi}
\newunicodechar{ᵇ}{\ifmmode^{b}\else\textsuperscript{\ensuremath{b}}\fi}
\newunicodechar{ʳ}{\ifmmode^{r}\else\textsuperscript{\ensuremath{r}}\fi}
\newunicodechar{ᵘ}{\ifmmode^{u}\else\textsuperscript{\ensuremath{u}}\fi}
\newunicodechar{ʸ}{\ifmmode^{y}\else\textsuperscript{\ensuremath{y}}\fi}
\newunicodechar{ᵃ}{\ifmmode^{a}\else\textsuperscript{\ensuremath{a}}\fi}
\newunicodechar{ᵉ}{\ifmmode^{e}\else\textsuperscript{\ensuremath{e}}\fi}
\newunicodechar{ᵏ}{\ifmmode^{k}\else\textsuperscript{\ensuremath{k}}\fi}
\newunicodechar{ᵖ}{\ifmmode^{p}\else\textsuperscript{\ensuremath{p}}\fi}
\newunicodechar{ᴺ}{\ifmmode^{N}\else\textsuperscript{\ensuremath{N}}\fi}
\newunicodechar{ᶜ}{\ifmmode^{c}\else\textsuperscript{\ensuremath{c}}\fi}
\newunicodechar{ʰ}{\ifmmode^{h}\else\textsuperscript{\ensuremath{h}}\fi}
\newunicodechar{ᵠ}{\ifmmode^{q}\else\textsuperscript{\ensuremath{q}}\fi}
\newunicodechar{ₙ}{\ifmmode_{n}\else\textsubscript{\ensuremath{n}}\fi}
\newunicodechar{ₐ}{\ifmmode_{a}\else\textsubscript{\ensuremath{a}}\fi}
\newunicodechar{ᵢ}{\ifmmode_{i}\else\textsubscript{\ensuremath{i}}\fi}
\newunicodechar{ᵣ}{\ifmmode_{r}\else\textsubscript{\ensuremath{r}}\fi}
\newunicodechar{ₖ}{\ifmmode_{k}\else\textsubscript{\ensuremath{k}}\fi}
\newunicodechar{ᵦ}{\ifmmode_{\beta}\else\textsubscript{\ensuremath{\beta}}\fi}
\newunicodechar{ₓ}{\ifmmode_{x}\else\textsubscript{\ensuremath{x}}\fi}
\newfontfamily\popdisplay{ArchivoBlack-Regular.ttf}[Path=../../../fonts/]
\newfontfamily\poplogo{SpaceGrotesk-Bold.ttf}[Path=../../../fonts/]
\newfontfamily\popsemi{PlusJakartaSans-SemiBold.ttf}[Path=../../../fonts/]
\newfontfamily\popxbold{PlusJakartaSans-ExtraBold.ttf}[Path=../../../fonts/]
\newfontfamily\popxboldls{PlusJakartaSans-ExtraBold.ttf}[Path=../../../fonts/,LetterSpace=3.0]

\setlist[enumerate]{leftmargin=1.7em,itemsep=5pt,topsep=4pt}
\setlist[itemize]{leftmargin=1.7em,itemsep=4pt,topsep=4pt,label={\color{poporange}\textbullet}}
\setlength{\parindent}{0pt}
\setlength{\parskip}{5pt}
\renewcommand{\arraystretch}{1.25}

% pgfplots : style Pop par défaut
\pgfplotsset{
  every axis/.append style={axis line style={popink,line width=1pt},tick style={popink},
    grid style={popline,line width=0.5pt},label style={font=\small},tick label style={font=\footnotesize}},
  popcurve/.style={poporange,line width=1.8pt,no markers,smooth},
}
% Même style disponible dans un \draw TikZ simple (hors pgfplots)
\tikzset{popcurve/.style={poporange,line width=1.8pt}}
\tikzset{popgrid/.style={popline,very thin}}
\colorlet{popcurve}{poporange} % le nom du style est parfois utilisé comme couleur

% ============================================================
% Pill / squiggle
% ============================================================
\newcommand{\poppill}[3]{%
  \tikz[baseline=-0.55ex]\node[draw=popink,line width=1.2pt,rounded corners=2.6pt,%
    fill=#2,inner xsep=6.5pt,inner ysep=3pt]{%
    \popxboldls\fontsize{7}{7}\selectfont\color{#3}\MakeUppercase{#1}};%
}
\newcommand{\popsquiggle}[1]{%
  \tikz[baseline=-0.4ex]{
    \draw[#1,line width=2.6pt,line cap=round]
      plot [smooth] coordinates {(0,0.09) (0.34,0.01) (0.68,0.21) (1.02,0.01) (1.36,0.10)};
  }%
}
% Mot-clé surligné (marqueur orange sous le texte)
\newcommand{\cle}[1]{{\popxbold\color{popdark}#1}}

% ============================================================
% En-tête / pied
% ============================================================
\pagestyle{fancy}
\fancyhf{}
\renewcommand{\headrulewidth}{0pt}
\renewcommand{\footrulewidth}{0pt}
\lhead{\raisebox{-0.15ex}{\poplogo\fontsize{13}{13}\selectfont\color{popink}OnBuch\color{poporange}+}}
\rhead{\poppill{\DOCMATIERE\ \textbullet\ \DOCNIVEAU\ \textbullet\ Leçon \DOCLECON}{white}{popink}}
\lfoot{\popxbold\fontsize{7.6}{7.6}\selectfont\color{popmuted}\MakeUppercase{Cours signé OnBuch+}\ \textbullet\ {\popsemi\DOCTITRE}}
\rfoot{\tikz[baseline=-0.6ex]\node[rounded corners=8.5pt,fill=popink,minimum height=17pt,minimum width=17pt,inner xsep=5pt,inner sep=0]{\popxbold\fontsize{8}{8}\selectfont\color{popcream}\thepage\,/\,\pageref*{LastPage}};}

% ============================================================
% Titres
% ============================================================
\newcommand{\chaptitre}[2]{%
  \begin{tcolorbox}[enhanced,colback=poporange,colframe=popink,boxrule=2.6pt,arc=11pt,
      drop shadow=popink,left=15pt,right=15pt,top=12pt,bottom=11pt,before skip=3pt,after skip=18pt]
    \poppill{#2}{popink}{popcream}\par\vspace{9pt}
    {\raggedright\hyphenpenalty=10000\exhyphenpenalty=10000\popdisplay\fontsize{22}{24}\selectfont\color{popcream}\MakeUppercase{#1}\par}
  \end{tcolorbox}
}
% Section numérotée : pastille orange avec le numéro + titre Archivo
\newcounter{popsec}
\newcommand{\coursec}[1]{%
  \stepcounter{popsec}\setcounter{popsub}{0}%
  \par\vspace{16pt}\noindent
  \begin{minipage}{\linewidth}
  \tikz[baseline=-0.7ex]\node[draw=popink,line width=1.6pt,fill=poporange,rounded corners=4pt,minimum size=22pt,inner sep=0]{\popdisplay\fontsize{12}{12}\selectfont\color{popcream}\thepopsec};%
  \hspace{8pt}{\popdisplay\fontsize{15}{17}\selectfont\color{popink}\MakeUppercase{#1}}\par
  \vspace{4pt}\noindent{\color{poporange}\rule{36pt}{3.6pt}}
  \end{minipage}\par\nopagebreak\vspace{8pt}
}
\newcounter{popsub}
\newcommand{\courssub}[1]{%
  \stepcounter{popsub}%
  \par\vspace{9pt}\noindent{\popxbold\fontsize{11.5}{13}\selectfont\color{popink}{\color{poporange}\thepopsec.\thepopsub}\hspace{6pt}#1}\par\nopagebreak\vspace{4pt}
}

% ============================================================
% Boîtes pédagogiques — même recette : fond blanc, bordure épaisse colorée,
% ombre dure encre, badge plein. Titre optionnel [#1].
% ============================================================
\tcbset{popbase/.style={breakable,enhanced,colback=white,boxrule=2.3pt,arc=9pt,
  shadow={3pt}{-3pt}{0pt}{popink},left=14pt,right=14pt,top=11pt,bottom=11pt,before skip=11pt,after skip=15pt,
  fontupper=\popsemi\fontsize{10.6}{14.5}\selectfont\color{popink},
  fontlower=\popsemi\fontsize{10.6}{14.5}\selectfont\color{popink}}}
% Coche dessinée (le glyphe ✓ n'existe pas dans les polices du document)
\newcommand{\popcheck}[1]{\tikz[baseline=-0.5ex]\draw[#1,line width=1.6pt,line cap=round,line join=round](-0.13,0.0)--(-0.04,-0.09)--(0.13,0.1);}
\providecommand{\coche}{\popcheck{popgreen}}
\providecommand{\resultat}[1]{\par\smallskip\noindent\fcolorbox{popgreen}{popgreenL}{\parbox{\dimexpr\linewidth-2\fboxsep-2\fboxrule\relax}{\textbf{Résultat.} #1}}\par\smallskip}
\newcommand{\popboxhead}[3]{\poppill{#1}{#2}{white}\ifx\relax#3\relax\else\hspace{6pt}{\popxbold\fontsize{10.6}{12}\selectfont\color{popink}#3}\fi\par\vspace{7pt}}

\newtcolorbox{aretenir}[1][]{popbase,colframe=popgreen,before upper={\popboxhead{À retenir}{popgreen}{#1}}}
% Texte en espagnol (document de lecture, dialogue, poème, extrait) : cadre
% d'étude. Un titre optionnel (source, auteur) s'affiche dans l'en-tête.
\newtcolorbox{textebox}[1][]{popbase,colframe=popblue,before upper={\popboxhead{Texto}{popblue}{#1}\begin{otherlanguage}{spanish}},after upper={\end{otherlanguage}}}
% Marques ✓ / ✗ (forme correcte / fautive) souvent utilisées par les modèles
\providecommand{\cross}{\textcolor{poppink}{\ensuremath{\times}}}
\providecommand{\xmark}{\cross}\providecommand{\crossmark}{\cross}\providecommand{\cmark}{\ensuremath{\checkmark}}
% Ligne de réponse / justification à compléter (inventée par les modèles)
\providecommand{\justification}[1]{\par\noindent\textit{Justification~:} \dotfill\par}
% \textipa (package tipa non chargé) : la police ne couvre pas l'API -> simple passage
\providecommand{\textipa}[1]{#1}
% Soulignement (ulem non chargé) : \uline, \ul
\providecommand{\uline}[1]{\underline{#1}}\providecommand{\ul}[1]{\underline{#1}}
% Ordinaux français écrits en macro par les modèles : 1\ière, 3\ième
\newcommand{\ière}{\textsuperscript{re}}\newcommand{\ième}{\textsuperscript{e}}
% \exemple (inventée par les modèles) : étiquette « Exemple : »
\providecommand{\exemple}{\textbf{Exemple~:}\ }
% Mot ou phrase en espagnol à l'intérieur d'un paragraphe français
\newcommand{\esp}[1]{\ifmmode\text{\textit{#1}}\else\begingroup\selectlanguage{spanish}\itshape #1\endgroup\fi}
\newtcolorbox{exemplebox}[1][]{popbase,colframe=poporange,before upper={\popboxhead{Exemple}{poporange}{#1}}}
\newtcolorbox{definition}[1][]{popbase,colframe=popblue,before upper={\popboxhead{Définition}{popblue}{#1}}}
\newtcolorbox{propriete}[1][]{popbase,colframe=poppurple,before upper={\popboxhead{Propriété}{poppurple}{#1}}}
\newtcolorbox{methode}[1][]{popbase,colframe=popgold,before upper={\popboxhead{Méthode}{popgold}{#1}}}
\newtcolorbox{attention}[1][]{popbase,colframe=poppink,before upper={\popboxhead{Attention}{poppink}{#1}}}
\newtcolorbox{experience}[1][]{popbase,colframe=popblue,colback=popblueL,before upper={\popboxhead{Expérience}{popblue}{#1}}}
% « Le savais-tu ? » : carte violette pleine (contexte, culture, Cameroun)
\newtcolorbox{savaistu}[1][]{popbase,colback=poppurple,colframe=popink,
  fontupper=\popsemi\fontsize{10.4}{14.2}\selectfont\color{white},
  before upper={\poppill{Le savais-tu\,?}{white}{poppurple}\ifx\relax#1\relax\else\hspace{6pt}{\popxbold\color{white}#1}\fi\par\vspace{7pt}}}
% Exercice résolu : énoncé en haut, \tcblower puis la solution sur fond crème
\newcounter{popexo}\newcounter{popexr}
\newtcolorbox{exoresolu}[1][]{popbase,colframe=poporange,colbacklower=popcream,
  segmentation style={popink,line width=1.2pt,dashed},
  before upper={\stepcounter{popexr}\popboxhead{Exercice résolu \thepopexr}{poporange}{#1}},
  before lower={\poppill{Solution}{popink}{popcream}\par\vspace{6pt}}}
% Exercice (non résolu en place) — la correction est dans la partie Corrigés
\newtcolorbox{exercice}[1][]{popbase,colframe=popink,
  before upper={\stepcounter{popexo}\popboxhead{Exercice \thepopexo}{popink}{#1}}}
% Corrigé
\newtcolorbox{corrige}[1][]{popbase,colframe=popgreen,colback=popgreenL,
  before upper={\popboxhead{Corrigé}{popgreen}{#1}}}
% Objectifs / prérequis / bilan
\newtcolorbox{objectifs}{popbase,colframe=popink,colback=poporangeL,
  before upper={\popboxhead{Objectifs de la leçon}{popink}{}}}
\newtcolorbox{prerequis}{popbase,colframe=popink,colback=popblueL,
  before upper={\popboxhead{Prérequis}{popblue}{}}}
\newtcolorbox{bilan}{popbase,colframe=popink,colback=white,boxrule=2.8pt,
  before upper={\popboxhead{Fiche bilan}{poporange}{L'essentiel à connaître pour l'examen}}}
% Figure encadrée avec légende
\newtcolorbox{popfigure}[1][]{enhanced,breakable=false,colback=white,colframe=popline,boxrule=1.4pt,arc=8pt,
  left=10pt,right=10pt,top=10pt,bottom=8pt,before skip=10pt,after skip=12pt,halign=center,
  lower separated=false,halign lower=center,
  fontlower=\popsemi\fontsize{9}{11.5}\selectfont\color{popmuted},
  #1}
\newcommand{\legende}[1]{\par\vspace{6pt}{\popsemi\fontsize{9}{11.5}\selectfont\color{popmuted}#1}}

% ============================================================
% Signature OnBuch+
% ============================================================
\newcommand{\poplogoinline}[1]{{\poplogo\fontsize{#1}{#1}\selectfont\color{popink}OnBuch\color{poporange}+}}
\newcommand{\signatureonbuch}{%
  \begin{tcolorbox}[enhanced,colback=popink,colframe=popink,boxrule=2.2pt,arc=10pt,
      drop shadow=poporange,left=14pt,right=14pt,top=10pt,bottom=10pt,before skip=0pt,after skip=0pt]
    \tikz[baseline=-0.7ex]\node[circle,fill=popgreen,draw=popcream,line width=1.3pt,minimum size=22pt,inner sep=0]{\popcheck{white}};%
    \hspace{9pt}\begin{minipage}[c]{0.62\linewidth}
      {\popxbold\fontsize{12}{13}\selectfont\color{popcream}Signé\ \textbullet\ {\poplogo OnBuch\color{poporange}+}}\par\vspace{2pt}
      {\raggedright\popsemi\fontsize{8}{10.5}\selectfont\color{popline}\MakeUppercase{Équipe pédagogique OnBuch+}\\\MakeUppercase{Conforme au programme officiel MINESEC}\par}
    \end{minipage}\hfill
    {\poplogo\fontsize{22}{22}\selectfont\color{popcream}OnBuch\color{poporange}+}
  \end{tcolorbox}%
}

% ============================================================
% Page de garde
% #1 titre · #2 matière · #3 niveau · #4 module · #5 n° leçon · #6 durée ·
% #7 description / objectifs (texte ou itemize)
% ============================================================
\newcommand{\pagegarde}[7]{%
  \thispagestyle{empty}
  \newgeometry{a4paper,margin=1.7cm,top=1.6cm,bottom=1.4cm}%
  \begin{tikzpicture}[remember picture,overlay]
    \fill[poporange!18] ([xshift=-3.2cm,yshift=-3.6cm]current page.north east) circle (4.6cm);
    \fill[poppurple!14] ([xshift=2.4cm,yshift=7.6cm]current page.south west) circle (3.8cm);
  \end{tikzpicture}%
  {\poplogo\fontsize{30}{30}\selectfont\color{popink}OnBuch\color{poporange}+}\hfill
  \raisebox{0.6ex}{\poppill{Cours signé \textbullet\ Premium}{popink}{popcream}}\par
  \vspace{1pt}\popsquiggle{poppurple}\par
  \vspace{8pt}
  \begin{tcolorbox}[enhanced,colback=poporange,colframe=popink,boxrule=2.8pt,arc=13pt,
      drop shadow=popink,left=18pt,right=18pt,top=12pt,bottom=13pt,after skip=10pt]
    \poppill{#2 \textbullet\ #3}{popink}{popcream}\hspace{5pt}\poppill{#4}{white}{popink}\par\vspace{8pt}
    {\popdisplay\fontsize{14}{14}\selectfont\color{popink}LEÇON #5}\par\vspace{4pt}
    {\raggedright\hyphenpenalty=10000\exhyphenpenalty=10000\popdisplay\fontsize{25}{27}\selectfont\color{popcream}\MakeUppercase{#1}\par}
    \vspace{8pt}
    {\popxbold\fontsize{9.5}{9.5}\selectfont\color{popink}\MakeUppercase{Durée conseillée : #6}}
  \end{tcolorbox}
  \begin{tcolorbox}[enhanced,colback=white,colframe=popink,boxrule=2.2pt,arc=10pt,
      drop shadow=popink,left=16pt,right=16pt,top=10pt,bottom=10pt,after skip=8pt]
    \poppill{Dans cette leçon}{poppurple}{white}\par\vspace{5pt}
    \popsemi\fontsize{9.3}{12.6}\selectfont\color{popink}#7
  \end{tcolorbox}
  \noindent\begin{minipage}[t]{0.487\linewidth}
    \begin{tcolorbox}[enhanced,colback=popgreen,colframe=popink,boxrule=2.2pt,arc=10pt,
        drop shadow=popink,left=11pt,right=11pt,top=10pt,bottom=10pt,height=3.1cm,valign=center]
      \tcbox[colback=white,colframe=popink,boxrule=1.3pt,arc=5pt,boxsep=0pt,nobeforeafter,
        left=3pt,right=3pt,top=3pt,bottom=3pt,valign=center]{\qrcode[height=1.9cm]{https://play.google.com/store/apps/details?id=com.luvvix.onbuch&hl=fr}}%
      \hspace{8pt}\begin{minipage}[c]{0.5\linewidth}
        {\popdisplay\fontsize{11}{12.5}\selectfont\color{white}\MakeUppercase{Télécharge OnBuch}}\par\vspace{4pt}
        {\raggedright\popsemi\fontsize{8.4}{11}\selectfont\color{white}Gratuit \textbullet\ Google Play\\Tuteur IA, annales, concours, résultats\par}
      \end{minipage}
    \end{tcolorbox}
  \end{minipage}\hfill
  \begin{minipage}[t]{0.487\linewidth}
    \begin{tcolorbox}[enhanced,colback=poppurple,colframe=popink,boxrule=2.2pt,arc=10pt,
        drop shadow=popink,left=11pt,right=11pt,top=10pt,bottom=10pt,height=3.1cm,valign=center]
      \tikz[baseline=-0.7ex]\node[circle,fill=white,draw=popink,line width=1.3pt,minimum size=1.6cm,inner sep=0]{\popdisplay\fontsize{24}{24}\selectfont\color{poppurple}+};%
      \hspace{8pt}\begin{minipage}[c]{0.6\linewidth}
        {\popdisplay\fontsize{11}{12.5}\selectfont\color{white}\MakeUppercase{Passe à OnBuch+}}\par\vspace{4pt}
        {\raggedright\popsemi\fontsize{8.4}{11}\selectfont\color{white}Tous les cours signés, Tuteur IA illimité, fiches PDF — onglet OnBuch+ de l'app\par}
      \end{minipage}
    \end{tcolorbox}
  \end{minipage}
  \vspace{8pt}
  \popcontenu
  \vfill
  \signatureonbuch
  \restoregeometry
  \newpage
}

% Rangée « Ce cours contient » (page de garde)
\newcommand{\poptile}[3]{% #1 couleur #2 gros texte #3 légende
  \begin{tcolorbox}[enhanced,colback=#1,colframe=popink,boxrule=2pt,arc=9pt,shadow={2.5pt}{-2.5pt}{0pt}{popink},
      left=6pt,right=6pt,top=9pt,bottom=9pt,halign=center,valign=center,height=1.75cm,width=\linewidth,nobeforeafter]
    {\popdisplay\fontsize{12.5}{13}\selectfont\color{white}\MakeUppercase{#2}}\par\vspace{3pt}
    {\centering\popsemi\fontsize{7.8}{9.5}\selectfont\color{white}#3\par}
  \end{tcolorbox}}
\newcommand{\popcontenu}{%
  \par\noindent{\popxboldls\fontsize{7.5}{7.5}\selectfont\color{popmuted}CE COURS SIGNÉ CONTIENT}\par\vspace{7pt}
  \noindent\begin{minipage}[t]{0.235\linewidth}\poptile{popblue}{Cours}{complet, illustré, pas à pas}\end{minipage}\hfill
  \begin{minipage}[t]{0.235\linewidth}\poptile{poporange}{Méthodes}{et exercices résolus}\end{minipage}\hfill
  \begin{minipage}[t]{0.235\linewidth}\poptile{poppink}{Exercices}{type examen + corrigés}\end{minipage}\hfill
  \begin{minipage}[t]{0.235\linewidth}\poptile{popgreen}{Fiche bilan}{l'essentiel à retenir}\end{minipage}\par}

% Sommaire
\newtcolorbox{sommaire}{enhanced,colback=white,colframe=popink,boxrule=2.2pt,arc=10pt,
  drop shadow=popink,left=18pt,right=18pt,top=13pt,bottom=13pt,before skip=6pt,after skip=20pt,
  before upper={\poppill{Sommaire}{poppurple}{white}\par\vspace{9pt}\popsemi\fontsize{10.6}{14.5}\selectfont\color{popink}}}

% Signature de fin de cours — poussée tout en bas de la dernière page
\newcommand{\findecours}{%
  \par\vfill\nopagebreak
  \begin{center}\popsquiggle{poporange}\end{center}\vspace{4pt}
  \signatureonbuch
}
\AtBeginDocument{\def\llbracket{\mathopen{[\mkern-3.5mu[}}\def\rrbracket{\mathclose{]\mkern-3.5mu]}}} % intervalles d'entiers (glyphe ⟦ absent de la police)
% Glyphes absents de Fira Math : redéfinis à partir de glyphes disponibles
\AtBeginDocument{%
  \def\setminus{\mathbin{\backslash}}\let\smallsetminus\setminus
  \def\vdots{\mathord{\vbox{\baselineskip3.2pt\lineskiplimit0pt\kern4pt\hbox{.}\hbox{.}\hbox{.}}}}%
  \def\ddots{\mathinner{\mkern1mu\raise6.5pt\hbox{.}\mkern2mu\raise3.6pt\hbox{.}\mkern2mu\raise0.7pt\hbox{.}\mkern1mu}}%
  \def\triangle{\mathord{\scalebox{0.9}{$\symup{\Delta}$}}}%
  \def\nabla{\mathord{\raisebox{\depth}{\scalebox{1}[-1]{$\symup{\Delta}$}}}}%
  \def\checkmark{\popcheck{popdark}}%
  \def\star{\mathbin{\ast}}\def\bigstar{\mathord{\ast}}%
  \def\diamond{\mathbin{\scalebox{0.6}{$\Diamond$}}}%
  \def\bigcup{\mathop{\mathchoice{\raisebox{-0.3em}{\scalebox{1.7}{$\cup$}}}{\scalebox{1.25}{$\cup$}}{\cup}{\cup}}\displaylimits}%
  \def\bigcap{\mathop{\mathchoice{\raisebox{-0.3em}{\scalebox{1.7}{$\cap$}}}{\scalebox{1.25}{$\cap$}}{\cap}{\cap}}\displaylimits}%
  \def\lesseqqgtr{\mathrel{\lessgtr}}\def\bigoplus{\mathop{\oplus}\displaylimits}%
}
\providecommand{\ssi}{si et seulement si} % abréviation fréquente
\AtBeginDocument{\providecommand{\wideparen}{\overparen}} % arc de cercle

\begin{document}

\pagegarde{Los medios y la información responsable}{Espagnol}{1ère A}{Módulo 5 : Les mass médias et la communication}{EA17}{2 h}{Cette leçon propose la lecture et le commentaire d'un article de presse sur les médias et les fausses informations. Elle fournit le lexique du journalisme et des réseaux sociaux, la méthode pour distinguer faits et opinions, et les outils pour exprimer son opinion en espagnol.
\vspace{4pt}
\begin{multicols}{2}\small
\begin{itemize}
\item À la fin de cette leçon, je sais lire et comprendre un article de presse simple en espagnol
\item À la fin de cette leçon, je sais utiliser le lexique des médias : prensa, noticia, periodista, titular, redes sociales, bulo, fuente, hecho, opinión
\item À la fin de cette leçon, je sais distinguer un fait d'une opinion dans un texte
\item À la fin de cette leçon, je sais repérer les indices d'une fausse information (bulo)
\item À la fin de cette leçon, je sais exprimer mon opinion avec creo que, en mi opinión, me parece que
\item À la fin de cette leçon, je sais rédiger un commentaire structuré d'un article de presse
\end{itemize}\end{multicols}}

\begin{sommaire}
\begin{multicols}{2}
\begin{enumerate}
\item Texte support : « ¿Verdad o bulo? »
\item Compréhension du texte
\item Lexique thématique : los medios de comunicación
\item Distinguer hecho y opinión — l'esprit critique
\item Méthode du commentaire de texte et commentaire modèle
\item Production guidée : escribir y opinar
\item Activité d'intégration
\item Méthodes et exercices résolus
\item Exercices
\item Corrigés des exercices
\item Fiche bilan
\end{enumerate}
\end{multicols}
\end{sommaire}

\begin{objectifs}
\begin{itemize}
\item À la fin de cette leçon, je sais lire et comprendre un article de presse simple en espagnol
\item À la fin de cette leçon, je sais utiliser le lexique des médias : prensa, noticia, periodista, titular, redes sociales, bulo, fuente, hecho, opinión
\item À la fin de cette leçon, je sais distinguer un fait d'une opinion dans un texte
\item À la fin de cette leçon, je sais repérer les indices d'une fausse information (bulo)
\item À la fin de cette leçon, je sais exprimer mon opinion avec creo que, en mi opinión, me parece que
\item À la fin de cette leçon, je sais rédiger un commentaire structuré d'un article de presse
\item À la fin de cette leçon, je sais vérifier une information avant de la partager
\item À la fin de cette leçon, je sais rédiger un court article ou un message responsable sur les réseaux sociaux
\end{itemize}
\end{objectifs}

\begin{prerequis}
\begin{itemize}
\item Présent de l'indicatif des verbes réguliers et irréguliers fréquents (Seconde)
\item Lexique général de la communication et des technologies (Seconde)
\item Structures de base de la phrase espagnole : sujet, verbe, compléments
\item Compréhension écrite de textes simples (stratégies de lecture vues en Seconde)
\item Passé composé (pretérito perfecto) pour rapporter des événements récents
\end{itemize}
\end{prerequis}

\newpage

\coursec{Texte support : « ¿Verdad o bulo? »}

\courssub{Situation d'accroche : une photo qui trompe tout le monde}

Imagine la scène. Pendant les élections au Cameroun, une photo circule sur WhatsApp : elle montrerait des urnes remplies de bulletins \emph{avant} même l'ouverture du vote. En quelques heures, des milliers de personnes la partagent, l'indignation monte, les commentaires s'enflamment... mais l'image date en réalité d'un scrutin organisé cinq ans plus tôt, dans un autre pays. Personne n'a vérifié. Tout le monde a cru.

En Espagne aussi, et dans toute l'Amérique latine, les \esp{bulos} (fausses informations) inondent les réseaux sociaux : fausses alertes sanitaires, rumeurs sur des personnalités, photos truquées. Face à ce phénomène, une question s'impose à tout lecteur de presse et à tout utilisateur des réseaux : comment distinguer un \esp{hecho} (un fait) d'une \esp{opinión} ? Comment vérifier une \esp{fuente} avant de partager ?

\textbf{Problématique :} qu'est-ce qu'une information responsable, et quels outils le lecteur possède-t-il pour ne pas se laisser tromper ? C'est tout l'enjeu de cette leçon, que nous allons découvrir à travers un article de presse hispanophone.

\begin{savaistu}[Le mot \esp{bulo} en Espagne]
En Espagne, le mot \esp{bulo} est entré dans le langage officiel : les médias sérieux l'utilisent dans leurs titres, et même la police nationale et la Garde civile publient régulièrement des messages sur Twitter/X pour démentir les \esp{bulos} qui circulent (« \esp{Es un bulo} », « \esp{No lo compartas} » — « Ne le partage pas »). Des journalistes spécialisés, appelés \esp{verificadores}, passent leur journée à vérifier les informations douteuses.
\end{savaistu}

\courssub{Première lecture : découverte globale du texte}

\begin{methode}[Lire un article de presse en cinq étapes]
Avant de lire en détail, adopte la démarche professionnelle du lecteur de presse :
\begin{enumerate}
\item Lire le \esp{titular} (le gros titre) : il annonce le sujet et parfois le point de vue du journaliste.
\item Identifier les questions essentielles : \esp{¿quién?} (qui ?), \esp{¿qué?} (quoi ?), \esp{¿dónde?} (où ?), \esp{¿cuándo?} (quand ?).
\item Repérer les \esp{hechos} : les faits vérifiables (dates, chiffres, événements).
\item Repérer les \esp{opiniones} : les jugements, les sentiments, les verbes comme \esp{creer}, \esp{pensar}, \esp{parecer}.
\item Vérifier la \esp{fuente} : qui a écrit l'article ? D'où vient l'information ?
\end{enumerate}
\end{methode}

Lis maintenant le texte suivant \textbf{silencieusement}, une première fois, sans dictionnaire. Objectif : répondre à deux questions simples. De quoi parle ce texte ? Quel type de texte est-ce ?

\begin{textebox}[¿Verdad o bulo? La información responsable]
Ayer por la tarde, Juan, un alumno de quince años, leyó en las redes sociales una noticia sorprendente: « ¡Cierran todas las escuelas mañana! ». El titular venía con una foto de una escuela vacía y muchos comentarios de personas enfadadas.

Sin comprobar la fuente, Juan compartió el titular con sus amigos y con los grupos de su barrio. En una hora, cientos de personas creyeron la noticia. Algunos padres decidieron no mandar a sus hijos a la escuela. Pero era un bulo: ningún periódico serio hablaba de eso, y el Ministerio de Educación no había publicado nada.

Al día siguiente, su profesora de lengua habló del caso en clase. Explicó la diferencia entre un hecho y una opinión: un hecho se puede comprobar, una opinión expresa lo que alguien piensa o siente. También enseñó a los alumnos a verificar siempre la fuente antes de compartir: ¿quién publica la noticia? ¿Es un medio conocido? ¿Otros periódicos dicen lo mismo?

Juan comprendió su error y pidió perdón a sus amigos. Ahora piensa que la información responsable es un deber de todos. « Antes de compartir, hay que pensar », dice. En mi opinión, esta lección es muy importante para todos los jóvenes que usan las redes sociales cada día.
\end{textebox}

\textbf{Glossaire (mots difficiles) :}

\begin{center}
\begin{tabularx}{\linewidth}{|X|X|X|}
\hline
\rowcolor{popblueL} Español & Français & Remarque \\
\hline
la prensa & la presse & ensemble des journaux \\
\hline
el titular & le gros titre & aussi « le titulaire » (d'un poste) \\
\hline
el bulo & la fausse information, l'intox & « fake news » en français courant \\
\hline
la fuente & la source & aussi « la fontaine » \\
\hline
el hecho & le fait & à ne pas confondre avec \esp{echo} (de \esp{echar}, je jette) \\
\hline
la opinión & l'opinion, l'avis & féminin : \esp{una opinión} \\
\hline
comprobar & vérifier & verbe avec changement \esp{o $\to$ ue} : \esp{compruebo} \\
\hline
compartir & partager & très utilisé sur les réseaux sociaux \\
\hline
sorprendente & surprenant & de \esp{sorprender} (surprendre) \\
\hline
enfadado & en colère, fâché & synonyme : \esp{enojado} (Amérique) \\
\hline
mandar & envoyer & aussi « commander » \\
\hline
el deber & le devoir & attention au pluriel : \esp{los deberes} = les devoirs (scolaires) \\
\hline
\end{tabularx}
\end{center}

\begin{definition}[¿Qué es un \esp{bulo} ?]
Un \esp{bulo} est une fausse information créée volontairement pour tromper le public et qui se propage rapidement, surtout sur les réseaux sociaux. En français, on parle de « fake news », d'« intox » ou de « rumeur ». Le verbe associé est \esp{desmentir} (démentir, radical changeant \esp{e $\to$ i} : \esp{desmiento, desmintió}) : \esp{El gobierno desmintió el bulo.} « Le gouvernement a démenti la fausse information. »
\end{definition}

\courssub{Deuxième lecture : repérage du lexique des médias}

Relis le texte une deuxième fois et \textbf{souligne} tous les mots qui appartiennent au champ lexical des médias. Tu devrais trouver au moins dix mots : \esp{noticia}, \esp{redes sociales}, \esp{titular}, \esp{fuente}, \esp{bulo}, \esp{periódico}, \esp{hecho}, \esp{opinión}, \esp{verificar}, \esp{compartir}, \esp{comentarios}, \esp{medio}.

Voici ce lexique organisé en carte mentale :

\begin{popfigure}
\begin{tikzpicture}[mindmap, grow cyclic, every node/.style={concept, align=center, font=\small},
root concept/.append style={concept color=popblue, font=\bfseries},
level 1/.append style={level distance=3.4cm, sibling angle=90},
level 2/.append style={level distance=2.2cm, sibling angle=45}]
\node[root concept]{Los medios}
  child[concept color=poporange]{ node {prensa escrita}
    child { node[concept color=poporangeL, text=popdark] {periódico, titular, periodista} } }
  child[concept color=popgreen]{ node {televisión y radio}
    child { node[concept color=popgreenL, text=popdark] {las noticias, el telediario} } }
  child[concept color=poppurple]{ node {redes sociales}
    child { node[concept color=poppinkL, text=popdark] {compartir, comentarios, bulo} } }
  child[concept color=popgold]{ node {verificación}
    child { node[concept color=popgoldL, text=popdark] {fuente, hecho, comprobar} } };
\end{tikzpicture}
\legende{Carte mentale du champ lexical des médias : quatre grandes familles de mots.}
\end{popfigure}

\begin{attention}[\esp{Noticia} : singulier ou pluriel ?]
En espagnol, \esp{una noticia} signifie « une nouvelle, une information ». Mais attention : \esp{las noticias} au pluriel désigne « les informations » ou « le journal télévisé » :
\begin{itemize}
\item \esp{He leído una noticia interesante.} « J'ai lu une nouvelle intéressante. »
\item \esp{¿Viste las noticias ayer?} « Tu as regardé le journal (télévisé) hier ? »
\end{itemize}
Ne traduis jamais « une nouvelle » par \esp{una nueva} : \esp{nueva} est l'adjectif « nouvelle » ; le nom est \esp{noticia}. De même, \esp{actualidad} signifie « l'actualité », pas « actuellement » (qui se dit \esp{actualmente} — faux ami classique !).
\end{attention}

\courssub{Lecture à voix haute : la prononciation}

Écoute d'abord la lecture de l'enseignant, puis lis à ton tour un paragraphe. Quelques points de prononciation à soigner :

\begin{itemize}
\item \esp{bulo} se prononce « bou-lo » : le \esp{b} initial est sonore, comme en français.
\item \esp{noticia} se prononce « no-ti-zia » en Espagne (le \esp{c} devant \esp{i} = « z » espagnol) et « no-ti-sia » en Amérique latine et en Guinée équatoriale.
\item \esp{compartió} : l'accent tombe sur la dernière syllabe « com-par-TIO » ; l'accent écrit est obligatoire (passé simple, 3\textsuperscript{e} personne).
\item \esp{opinión} : « o-pi-NION », accent final marqué.
\item \esp{hecho} : le \esp{h} est muet, le \esp{ch} se prononce « tch » : « è-tcho ».
\item \esp{periódico} : « pè-RI-o-di-ko », quatre syllabes, accent sur le \esp{i}.
\end{itemize}

\begin{propriete}[L'accent écrit des noms en \esp{-ción}]
Les noms féminins en \esp{-ción} portent toujours un accent écrit sur le \esp{ó} : \esp{opinión}, \esp{información}, \esp{educación}, \esp{verificación}. Au pluriel, l'accent disparaît : \esp{opiniones}, \esp{informaciones}. Compare avec le français : « opinion », « information » — même famille de mots, mais l'espagnol marque l'accent tonique final par un accent aigu.
\end{propriete}

\courssub{Le parcours d'une information : du fait au lecteur}

Le texte de Juan montre comment une information naît, circule... et parfois se déforme. Observe ce schéma :

\begin{popfigure}
\begin{tikzpicture}[node distance=0.55cm,
caja/.style={draw=popblue, fill=popblueL, rounded corners=3pt, align=center, font=\small, minimum height=0.9cm, inner sep=4pt},
flecha/.style={-{Stealth[length=2.5mm]}, thick, poporange}]
\node[caja, align=center] (fuente){\textbf{Fuente}\\Origine du fait\\témoin, document};
\node[caja, right=of fuente, align=center] (periodista){\textbf{Periodista}\\Enquête, vérifie,\\recoupe les sources};
\node[caja, right=of periodista, align=center] (titular){\textbf{Titular}\\Le média publie\\l'information vérifiée};
\node[caja, right=of titular, align=center] (lector){\textbf{Lector}\\Lit, analyse,\\vérifie à son tour};
\node[caja, right=of lector, align=center] (redes){\textbf{Redes sociales}\\Partage responsable\\ou propagation du bulo};
\draw[flecha] (fuente) -- (periodista);
\draw[flecha] (periodista) -- (titular);
\draw[flecha] (titular) -- (lector);
\draw[flecha] (lector) -- (redes);
\end{tikzpicture}
\legende{Le parcours d'une information fiable. Si le lecteur partage sans vérifier (\esp{sin comprobar}), la chaîne sérieuse se brise et le \esp{bulo} se propage.}
\end{popfigure}

Dans le texte, Juan a sauté l'étape essentielle : il est passé directement du \esp{titular} au partage, sans vérifier la \esp{fuente}. C'est exactement ce qui s'est passé avec la photo des urnes au Cameroun évoquée dans l'accroche.

\courssub{Reformulation orale : à toi de parler}

\begin{experience}[Resumir el texto en dos o tres frases]
En binôme, reformulez oralement le texte en deux ou trois phrases, sans le regarder. Utilisez les structures suivantes :
\begin{itemize}
\item \esp{El texto habla de...} « Le texte parle de... »
\item \esp{Juan compartió... sin comprobar...} « Juan a partagé... sans vérifier... »
\item \esp{La profesora explicó que...} « La professeure a expliqué que... »
\item \esp{La lección del texto es que...} « La leçon du texte, c'est que... »
\end{itemize}
Modèle attendu : \esp{El texto habla de Juan, un alumno que compartió un bulo en las redes sociales sin comprobar la fuente. Su profesora le enseñó a verificar la información. Ahora Juan piensa que la información responsable es un deber de todos.}
\end{experience}

\begin{exoresolu}[Compréhension globale : vrai ou faux]
Énoncé — Lis les phrases suivantes et dis si elles sont \esp{verdadero} (vrai) ou \esp{falso} (faux) d'après le texte, en justifiant avec une phrase du texte :
\begin{enumerate}
\item \esp{Juan comprobó la fuente antes de compartir la noticia.}
\item \esp{Ningún periódico serio hablaba del cierre de las escuelas.}
\item \esp{La profesora explicó la diferencia entre un hecho y una opinión.}
\end{enumerate}
\tcblower
Solution détaillée :
\begin{enumerate}
\item \textbf{Falso.} Le texte dit : \esp{Sin comprobar la fuente, Juan compartió el titular con sus amigos.} Juan a partagé \emph{sans} vérifier. Piège : le mot \esp{sin} (sans) inverse le sens de la phrase.
\item \textbf{Verdadero.} Le texte dit : \esp{Pero era un bulo: ningún periódico serio hablaba de eso.} Aucun journal sérieux ne parlait de cette fermeture : c'était donc une fausse information.
\item \textbf{Verdadero.} Le texte dit : \esp{Explicó la diferencia entre un hecho y una opinión: un hecho se puede comprobar, una opinión expresa lo que alguien piensa o siente.}
\end{enumerate}
\end{exoresolu}

\begin{aretenir}[L'essentiel de cette étape]
\begin{itemize}
\item Un article de presse se lit en cinq étapes : \esp{titular} $\to$ qui/quoi/où/quand $\to$ \esp{hechos} $\to$ \esp{opiniones} $\to$ \esp{fuente}.
\item Un \esp{hecho} se peut vérifier ; une \esp{opinión} exprime ce que quelqu'un pense ou ressent.
\item Un \esp{bulo} est une fausse information créée pour tromper ; avant de \esp{compartir}, il faut \esp{comprobar la fuente}.
\item Le lexique clé : \esp{prensa}, \esp{noticia}, \esp{periodista}, \esp{titular}, \esp{redes sociales}, \esp{bulo}, \esp{fuente}, \esp{hecho}, \esp{opinión}.
\end{itemize}
\end{aretenir}

\courssub{Distinguer \esp{hecho} et \esp{opinión} : l'esprit critique}

C'est la compétence centrale de ce module. Un lecteur responsable doit pouvoir étiqueter chaque phrase d'un texte : \textbf{Fait} (objectif, vérifiable) ou \textbf{Opinion} (subjectif, discutable).

\begin{propriete}[Marqueurs linguistiques : \esp{Hecho} vs \esp{Opinión}]
\begin{center}
\begin{tabularx}{\linewidth}{|X|X|X|}
\hline
\rowcolor{popblueL} Critère & \textbf{Hecho (Fait)} & \textbf{Opinión (Opinion)} \\
\hline
\cle{Définition} & Réalité observable, mesurable, datable, vérifiable par des sources indépendantes. & Jugement de valeur, sentiment, interprétation, croyance personnelle. \\
\hline
\cle{Verbes fréquents} & \esp{ser/estar} (objectif), \esp{hay}, \esp{ocurrir}, \esp{publicar}, \esp{confirmar}, \esp{anunciar}, \esp{cifrar}. & \esp{creer}, \esp{pensar}, \esp{opinar}, \esp{considerar}, \esp{parecer}, \esp{sentir}, \esp{preferir}, \esp{deber} (obligation morale). \\
\hline
\cle{Adjectifs / Noms} & Chiffres, dates, noms propres, termes techniques, \esp{según...}, \esp{de acuerdo con...}. & \esp{importante}, \esp{grave}, \esp{mejor/peor}, \esp{injusto}, \esp{necesario}, \esp{imprescindible}, \esp{escándalo}, \esp{éxito}. \\
\hline
\cle{Connecteurs} & \esp{porque} (cause réelle), \esp{ya que}, \esp{puesto que}, \esp{según la fuente X}. & \esp{en mi opinión}, \esp{para mí}, \esp{a mi juicio}, \esp{desde mi punto de vista}, \esp{me parece que}. \\
\hline
\cle{Test de vérification} & « Peut-on le prouver par un document, une photo, un témoin, une loi ? » $\to$ \textbf{OUI}. & « Tout le monde est-il d'accord ? » $\to$ \textbf{NON} (c'est discutable). \\
\hline
\end{tabularx}
\end{center}
\end{propriete}

\begin{exemplebox}[Classification dans le texte « ¿Verdad o bulo? »]
\begin{itemize}
\item \esp{« Juan... leyó en las redes sociales una noticia »} $\to$ \textbf{Hecho} (action datée, vérifiable).
\item \esp{« El titular venía con una foto... »} $\to$ \textbf{Hecho} (description objective).
\item \esp{« Pero era un bulo »} $\to$ \textbf{Hecho} (constatation après vérification).
\item \esp{« Ningún periódico serio hablaba de eso »} $\to$ \textbf{Hecho} (vérifiable : consultation des kiosques/sites).
\item \esp{« Un hecho se puede comprobar »} $\to$ \textbf{Hecho} (définition).
\item \esp{« Una opinión expresa lo que alguien piensa o siente »} $\to$ \textbf{Hecho} (définition).
\item \esp{« En mi opinión, esta lección es muy importante »} $\to$ \textbf{Opinión} (marqueur \esp{En mi opinión}, adjectif subjectif \esp{importante}).
\item \esp{« La información responsable es un deber de todos »} $\to$ \textbf{Opinión} (jugement moral \esp{deber}, bien que présenté comme vérité partagée).
\end{itemize}
\end{exemplebox}

\begin{attention}[Piège : \esp{Creo que} + Indicatif vs \esp{No creo que} + Subjonctif]
C'est une règle d'or du niveau B1 pour exprimer l'opinion :
\begin{itemize}
\item \textbf{Affirmatif :} \esp{Creo que / Pienso que / Opino que / Me parece que} + \textbf{INDICATIF} (certitude du locuteur).
  \exemple \esp{Creo que el bulo es peligroso.} (Indicatif \esp{es})
\item \textbf{Négatif :} \esp{No creo que / No pienso que / No me parece que} + \textbf{SUBJONCTIF} (doute, subjectivité).
  \exemple \esp{No creo que el bulo \textbf{sea} peligroso.} (Subjonctif \esp{sea})
\item \textbf{Structures impersonnelles :} \esp{Es importante / necesario / urgente que} + \textbf{SUBJONCTIF}.
  \exemple \esp{Es importante que \textbf{verifiquemos} la fuente.}
\end{itemize}
\emph{Attention :} \esp{En mi opinión}, \esp{Para mí}, \esp{A mi juicio} appellent l'indicatif (ce sont des cadres énonciatifs, pas des verbes de doute).
\end{attention}

\begin{exercice}[Entraînement : Hecho u Opinión]
Classifie les phrases suivantes en \esp{Hecho (H)} ou \esp{Opinión (O)}. Justifie ton choix par le marqueur linguistique repéré.
\begin{enumerate}
\item \esp{El Ministerio de Educación publicó un comunicado ayer.}
\item \esp{Es urgente que los jóvenes aprendan a verificar fuentes.}
\item \esp{Según el periódico \textit{El País}, la inflación bajó al 3\%.}
\item \esp{Me parece que las redes sociales son un peligro para la democracia.}
\item \esp{El 60\% de los estudiantes comparten noticias sin leerlas.}
\item \esp{No creo que Juan haya actuado con mala fe.}
\end{enumerate}
\end{exercice}

\begin{corrige}[Exercice : Hecho u Opinión]
\begin{enumerate}
\item \textbf{H.} Verbe \esp{publicó} (passé simple, action accomplie), source identifiable (\esp{Ministerio}).
\item \textbf{O.} Structure impersonnelle \esp{Es urgente que} + \textbf{subjonctif} \esp{aprendan} (jugement de valeur).
\item \textbf{H.} Introducteur \esp{Según} + source précise (\esp{El País}), chiffre vérifiable (\esp{3\%}).
\item \textbf{O.} Marqueur \esp{Me parece que} + adjectif subjectif \esp{peligro}. Note : indicatif \esp{son} car phrase affirmative.
\item \textbf{H.} Statistique précise (\esp{60\%}), bien qu'il faille vérifier la source de ce chiffre.
\item \textbf{O.} \esp{No creo que} + \textbf{subjonctif} \esp{haya actuado} (doute du locuteur sur l'intention).
\end{enumerate}
\end{corrige}

\courssub{Méthode du commentaire de texte et commentaire modèle}

L'épreuve de commentaire de texte (examen de Première, Baccalauréat) suit une structure rigoureuse. Voici la méthode officielle en 4 étapes.

\begin{methode}[Commenter un texte de presse : 4 étapes]
\begin{enumerate}
\item \textbf{Introduction (1 paragraphe) :}
  \begin{itemize}
  \item \textit{Accroche} : contexte général (thème des médias, fake news, citoyenneté numérique).
  \item \textit{Présentation du texte} : titre (entre guillemets), auteur (si connu, sinon « artículo anónimo »), nature (article de presse / récit didactique), source (manuelle), date (si connue), thème principal.
  \item \textit{Problématique} : question ouverte à laquelle le texte répond (ex. « Comment ce texte illustre-t-il la nécessité de vérifier l'information avant de la partager ? »).
  \item \textit{Annonce du plan} : « Nous analyserons d'abord... puis... enfin... »
  \end{itemize}
\item \textbf{Développement (2 ou 3 parties, 1 paragraphe chacune) :}
  \begin{itemize}
  \item \textit{Partie 1 : Analyse du contenu — Faits vs Opinions.} Repérer les \esp{hechos} (intrigue : Juan, bulo, école) et les \esp{opiniones} (leçon de la prof, conclusion de Juan, commentaire final du narrateur). Citer le texte (guillemets, numéros de ligne).
  \item \textit{Partie 2 : Analyse de la structure et de la langue.} Type de discours (narratif $\to$ didactique), temps verbaux (passé \esp{pretérito indefinido/imperfecto} pour le récit, présent/généralités pour la leçon), connecteurs logiques (\esp{pero, también, ahora}), champ lexical (médias, éducation).
  \item \textit{Partie 3 (si place) : Portée citoyenne / Éducation aux médias.} Le texte comme outil pédagogique : modélisation d'une erreur, correction par l'école, message final \esp{deber de todos}.
  \end{itemize}
\item \textbf{Conclusion (1 paragraphe) :}
  \begin{itemize}
  \item \textit{Bilan} : réponse synthétique à la problématique.
  \item \textit{Ouverture} : lien avec l'actualité (Cameroun, Guinée équatoriale, Espagne), rôle de l'école, défi de l'IA (deepfakes), responsabilité individuelle.
  \end{itemize}
\item \textbf{Relecture :} orthographe, accents, accord participes, cohérence du français.
\end{enumerate}
\end{methode}

\begin{exemplebox}[Commentaire modèle : « ¿Verdad o bulo? »]
\textbf{Introduction}
Dans un contexte où les \esp{bulos} (fausses informations) circulent plus vite que les vérités sur WhatsApp, Facebook ou X, l'éducation aux médias devient une urgence citoyenne. Le texte « ¿Verdad o bulo? La información responsable », extrait de notre manuel de Première, est un court récit didactique anonyme qui met en scène Juan, un élève de 15 ans, victime et vecteur d'une fausse annonce de fermeture d'écoles. Comment ce texte illustre-t-il le passage de la crédulité à l'esprit critique ? Nous verrons d'abord comment l'opposition faits/opinions structure le récit, puis comment la langue sert la visée pédagogique, enfin la portée civique du message.

\textbf{Développement}
\textit{1. Une intrigue construite sur la confusion fait/opinion.}
Le texte suit une chronologie claire : l'erreur (l. 1-5), la propagation (l. 5-7), la vérité rétablie (l. 7-9), la correction pédagogique (l. 10-14), la prise de conscience (l. 15-17). Les \textbf{faits} dominent la première partie : verbes d'action au \esp{pretérito indefinido} (\esp{leyó, compartió, creyeron, decidieron, era, hablaba, había publicado}), indicateurs spatio-temporels (\esp{ayer por la tarde, en una hora, al día siguiente}). La phrase « \esp{Pero era un bulo: ningún periódico serio hablaba de eso} » (l. 8-9) est le pivot factuel : elle oppose la rumeur à la réalité vérifiable (absence de source officielle). La seconde partie bascule vers l'\textbf{opinion} : la professeure énonce une définition (\esp{un hecho se puede comprobar, una opinión expresa lo que alguien piensa o siente}), puis le narrateur conclut par un jugement de valeur explicite : « \esp{En mi opinión, esta lección es muy importante} » (l. 17). Juan lui-même intériorise une norme morale : « \esp{la información responsable es un deber de todos} » (l. 16).

\textit{2. Une langue au service de la modélisation citoyenne.}
Le vocabulaire technique (\esp{titular, fuente, bulo, verificar, periódico, Ministerio}) ancre le texte dans le champ lexical des médias. La syntaxe privilégie la clarté didactique : phrases courtes, coordination (\esp{y, pero}), subordination explicative (\esp{que} + indicatif). L'alternance des temps marque la progression : passé narratif pour l'anecdote, présent de vérité générale pour la leçon (\esp{es, piensa, dice, usan}), impératif implicite (\esp{hay que pensar}) pour la prescription.

\textbf{Conclusion}
Ce texte remplit parfaitement sa fonction : transformer un incident banal en leçon d'hygiène informationnelle. Il montre que l'esprit critique s'apprend (\esp{enseñó a los alumnos a verificar}) et que la responsabilité incombe à chaque maillon de la chaîne (lecteur, partageur). Ouverture : au Cameroun comme en Guinée équatoriale, la viralité des \esp{bulos} électoraux ou sanitaires menace la cohésion sociale. L'école a pour mission de former des \esp{verificadores} citoyens, capables de dire « \esp{Es un bulo, no lo compartas} » avant qu'il ne soit trop tard.
\end{exemplebox}

\courssub{Production guidée : \esp{escribir y opinar}}

À ton tour d'écrire ! Tu vas produire un court texte d'opinion (80-100 mots) pour alerter tes camarades sur un \esp{bulo} circulant dans ton établissement ou ton quartier.

\begin{propriete}[Exprimer son opinion : structures clés (niveau A2+/B1)]
\begin{center}
\begin{tabularx}{\linewidth}{|l|X|X|}
\hline
\rowcolor{popblueL} Structure & Régime & Exemple traduit \\
\hline
\esp{Creo que / Pienso que / Opino que} & \textbf{Indicatif} & \esp{Creo que \textbf{es} un bulo.} (Je crois que c'est un bulo.) \\
\hline
\esp{No creo que / No pienso que} & \textbf{Subjonctif} & \esp{No creo que \textbf{sea} verdad.} (Je ne crois pas que ce soit vrai.) \\
\hline
\esp{Me parece que} & \textbf{Indicatif} (souvent) & \esp{Me parece que \textbf{es} peligroso.} (Il me semble que c'est dangereux.) \\
\hline
\esp{No me parece que} & \textbf{Subjonctif} & \esp{No me parece que \textbf{haya} pruebas.} (Il ne me semble pas qu'il y ait des preuves.) \\
\hline
\esp{En mi opinión / Para mí / A mi juicio} & \textbf{Indicatif} & \esp{En mi opinión, \textbf{hay} que verificar.} (À mon avis, il faut vérifier.) \\
\hline
\esp{Es + adj (importante, necesario, urgente) que} & \textbf{Subjonctif} & \esp{Es necesario que \textbf{comprobemos} la fuente.} (Il faut que nous vérifiions.) \\
\hline
\end{tabularx}
\end{center}
\end{propriete}

\begin{experience}[Écrire une « Alerta Bulo » pour le journal du lycée]
\textbf{Consigne :} Un \esp{bulo} circule : « \esp{El examen de español del lunes está anulado} ». Tu as vérifié : c'est faux (source : professeur principal, emploi du temps officiel). Rédige un court message (style tweet / post Instagram / panneau d'affichage) pour le démentir.
\textbf{Contraintes :}
\begin{itemize}
\item Titre accrocheur (hashtag bienvenu).
\item Annonce le bulo clairement.
\item Apporte la preuve (source officielle).
\item Donne ton opinion / conseil en utilisant \textbf{au moins 3 structures différentes} du tableau ci-dessus.
\item Ton : calme, responsable, pédagogique. Pas d'insultes.
\item 80-100 mots.
\end{itemize}
\end{experience}

\begin{exemplebox}[Modèle de production attendue]
\textbf{\#AlertaBulo \#InfoResponsable}

\esp{¡Atención! Circula en los grupos de WhatsApp que \textbf{el examen de español del lunes está anulado}. \textbf{ES UN BULO}. He verificado con el profesor titular y el horario oficial: el examen se mantiene a las 8:00.}

\esp{\textbf{En mi opinión}, compartir esto sin comprobar crea pánico innecesario. \textbf{Creo que} cada uno tiene la responsabilidad de \textbf{verificar la fuente} antes de reenviar. \textbf{No creo que} sea gracioso jugar con la ansiedad de los compañeros. \textbf{Es urgente que} usemos las redes con cabeza. \textbf{¡Piensa antes de compartir!}}
\end{exemplebox}

\begin{exercice}[Production écrite notée : Article d'opinion]
\textbf{Sujet :} « \esp{Las redes sociales: ¿ventana al mundo o caja de bulos?} » (Les réseaux sociaux : fenêtre sur le monde ou boîte à fake news ?).
Rédige un article d'opinion structuré (150-180 mots) pour le journal numérique de ton lycée.
\textbf{Plan imposé :}
\begin{enumerate}
\item Titre original.
\item Introduction : accroche + thèse (ton avis général).
\item Argument 1 : Les avantages (information rapide, lien diaspora, mobilisation citoyenne). Exemple concret (Cameroun / Guinée équatoriale / Espagne).
\item Argument 2 : Les dangers (bulos, bulles de filtre, haine, deepfakes). Exemple concret.
\item Argument 3 : La solution (éducation, vérification, régulation ?). Ton avis personnel (\esp{En mi opinión...}, \esp{Creo que...}, \esp{Es necesario que...}).
\item Conclusion : phrase choc / appel à l'action.
\end{enumerate}
\end{exercice}

\begin{corrige}[Exercice : Article d'opinion — Proposition de correction]
\textbf{Titre : \#PiensaAntesDeCompartir: La doble cara de las redes}

\esp{¿Ventana al mundo o caja de bulos? Las redes sociales son hoy la plaza pública de la juventud camerunesa e hispana. \textbf{En mi opinión}, son una herramienta poderosa, pero su uso exige una ciudadanía digital madura.}

\esp{Por un lado, \textbf{creo que} democratizan la información. En Yaoundé o en Malabo, un estudiante accede en segundos a la prensa de Madrid, Buenos Aires o París. Permiten movilizaciones ciudadanas: \textbf{he visto} campañas contra la corrupción o por el clima nacer en Twitter/X y llegar a la televisión.}

\esp{Por otro lado, \textbf{me parece que} el algoritmo premia lo viral, no lo vero. Los \textbf{bulos} electorales, los montajes de vídeo (deepfakes) y el discurso de odio se propagan sin freno. \textbf{No creo que} la autorregulación de las plataformas baste: \textbf{es imprescindible que} la escuela enseñe a \textbf{verificar fuentes} desde la Primaria.}

\esp{En definitiva, la red no es buena ni mala: es un espejo de quien la usa. \textbf{Propongo} un pacto: « \textbf{Verifico, luego comparto} ». \textbf{¡Tu dedo tiene el poder: úsalo con cabeza!}}
\end{corrige}

\begin{aretenir}[Synthèse de la leçon EA17 — Los medios y la información responsable]
\begin{itemize}
\item \textbf{Lexique clé :} \cle{prensa}, \cle{noticia}, \cle{periodista}, \cle{titular}, \cle{redes sociales}, \cle{bulo}, \cle{fuente}, \cle{hecho}, \cle{opinión}, \cle{verificar}, \cle{compartir}, \cle{desmentir}.
\item \textbf{Compétence lecture :} Méthode en 5 étapes (Titular $\to$ Qui/Qué/Dónde/Cuándo $\to$ Hechos $\to$ Opiniones $\to$ Fuente).
\item \textbf{Pensée critique :} \esp{Hecho} = vérifiable (indicatif, sources, chiffres) ; \esp{Opinión} = subjectif (verbes de pensée, adjectifs de valeur, marqueurs \esp{en mi opinión}).
\item \textbf{Grammaire opinion :} \esp{Creo que} + Ind. / \esp{No creo que} + Subj. / \esp{En mi opinión} + Ind. / \esp{Es necesario que} + Subj.
\item \textbf{Écriture :} Commentaire de texte (Intro / Développement 2-3 parties / Conclusion + Ouverture) ; Article d'opinion structuré (Thèse / Arguments / Exemples / Conclusion).
\item \textbf{Attitude :} \esp{Antes de compartir, hay que pensar.} La información responsable es un \cle{deber de todos}.
\end{itemize}
\end{aretenir}

\coursec{Los medios y la información responsable}

\courssub{Texte support : ¿Verdad o bulo?}

\begin{textebox}[¿Verdad o bulo?]
Juan es un estudiante de dieciséis años que vive en Yaoundé. Como todos los jóvenes de su edad, pasa mucho tiempo en las redes sociales. Un lunes por la tarde, leyó en su teléfono una noticia sorprendente: el titular decía que un famoso jugador de fútbol venía a visitar su barrio el sábado siguiente. Juan se puso muy contento. Sin comprobar la fuente, compartió el titular con todos sus amigos y con su familia. En pocas horas, la noticia circuló por todo el barrio y mucha gente la creyó.

El sábado, decenas de personas esperaron al jugador delante del campo de fútbol. Pero nadie vino. La noticia era un bulo: la página que la publicó no era un medio serio, sino un sitio creado para hacer bromas. Juan se sintió avergonzado. Su profesora de español le explicó que antes de compartir una noticia hay que verificar la fuente, comparar varios medios y distinguir los hechos de las opiniones.

Desde aquel día, Juan aprendió una lección importante: la información responsable es un deber de todos. Ahora, cuando lee un titular sorprendente, siempre se pregunta: ¿es verdad o es un bulo?
\end{textebox}

\textbf{Glossaire :} \esp{el barrio} : le quartier ; \esp{compartir} : partager ; \esp{la fuente} : la source ; \esp{el bulo} : la fausse information, le canular, le \emph{fake news} ; \esp{avergonzado} : honteux ; \esp{la broma} : la blague ; \esp{el deber} : le devoir ; \esp{aquel} : ce...-là (démonstratif éloigné) ; \esp{la noticia} : la nouvelle, l'information ; \esp{el titular} : le titre (d'un article) ; \esp{verificar} : vérifier ; \esp{el medio (de comunicación)} : le média.

\begin{definition}[Lexique fondamental : los medios de comunicación]
\begin{itemize}
\item \esp{La prensa} : la presse (écrite). \esp{El periódico} : le journal. \esp{La revista} : le magazine.
\item \esp{La noticia} : l'information, la nouvelle (événement rapporté).
\item \esp{El/la periodista} : le/la journaliste.
\item \esp{El titular} : le titre, la une.
\item \esp{Las redes sociales} : les réseaux sociaux (ex. \esp{WhatsApp, Facebook, X, TikTok}).
\item \esp{La fuente} : la source (origine de l'info).
\item \esp{El bulo} : la rumeur, l'infox, le canular (\emph{fake news}). Syn. : \esp{la noticia falsa}, \esp{el bulo informativo}.
\item \esp{El hecho} : le fait (réalité vérifiable).
\item \esp{La opinión} : l'opinion (jugement subjectif).
\item \esp{Verificar / Comprobar} : vérifier, contrôler.
\item \esp{Contrastar} : recouper, croiser les sources.
\end{itemize}
\end{definition}

\courssub{Compréhension du texte}

Avant d'étudier la distinction entre \esp{hecho} et \esp{opinión}, il faut s'assurer que le texte est \cle{parfaitement compris}. Voici la méthode rigoureuse exigée au baccalauréat.

\begin{methode}[Répondre à une question de compréhension en trois étapes]
\begin{enumerate}
\item \textbf{Repérer} : retrouve dans le texte la ligne ou le passage qui contient la réponse. Souligne-le.
\item \textbf{Citer la preuve} : recopie les mots exacts du texte entre guillemets espagnols (\esp{« ... »}). Exemple : \esp{« Sin comprobar la fuente, compartió el titular »}.
\item \textbf{Reformuler} : rédige ta réponse avec tes propres mots, en phrase complète, en reprenant les termes de la question. Exemple : \esp{Juan compartió el titular sin comprobar la fuente.}
\end{enumerate}
Une bonne réponse = \textbf{preuve citée} + \textbf{reformulation personnelle}. Ni l'une ni l'autre ne suffit seule.
\end{methode}

\courssub{Questions de compréhension globale}

\begin{itemize}
\item \esp{¿De qué habla el texto?} — De quoi parle le texte ?
      \begin{exemplebox}[Réponse modèle]
      \esp{El texto habla de un joven camerunés, Juan, que comparte una noticia falsa en las redes sociales sin verificarla.} (Le texte parle d'un jeune Camerounais, Juan, qui partage une fausse nouvelle sur les réseaux sociaux sans la vérifier.)
      \end{exemplebox}
\item \esp{¿Quién es el protagonista?} — Qui est le protagoniste ?
      \begin{exemplebox}[Réponse modèle]
      \esp{El protagonista es Juan, un estudiante de dieciséis años que vive en Yaoundé.}
      \end{exemplebox}
\item \esp{¿Dónde y cuándo ocurre la historia?} — Où et quand se passe l'histoire ?
      \begin{exemplebox}[Réponse modèle]
      \esp{La historia ocurre en Yaoundé, un lunes por la tarde y el sábado siguiente.}
      \end{exemplebox}
\end{itemize}

\courssub{Questions de détail}

La réponse se trouve presque toujours dans une seule phrase. Applique la méthode en 3 étapes.

\begin{exoresolu}[Exemples résolus]
Énoncé : \esp{¿Qué noticia leyó Juan?}
\tcblower
Solution détaillée :
\begin{enumerate}
\item \textbf{Repérer} : 1er paragraphe, phrase 3.
\item \textbf{Citer} : \esp{« el titular decía que un famoso jugador de fútbol venía a visitar su barrio el sábado siguiente »}.
\item \textbf{Reformuler} : \esp{Juan leyó que un famoso jugador de fútbol iba a visitar su barrio el sábado.}
\end{enumerate}
\end{exoresolu}

\begin{exoresolu}[Exemples résolus]
Énoncé : \esp{¿Qué hizo con la noticia?}
\tcblower
Solution détaillée :
\begin{enumerate}
\item \textbf{Repérer} : 1er paragraphe, phrase 5.
\item \textbf{Citer} : \esp{« compartió el titular con todos sus amigos y con su familia »}.
\item \textbf{Reformuler} : \esp{Compartió el titular con sus amigos y su familia.}
\end{enumerate}
\end{exoresolu}

\begin{exoresolu}[Exemples résolus]
Énoncé : \esp{¿Era verdad la noticia?}
\tcblower
Solution détaillée :
\begin{enumerate}
\item \textbf{Repérer} : 2e paragraphe, phrase 3.
\item \textbf{Citer} : \esp{« La noticia era un bulo »}.
\item \textbf{Reformuler} : \esp{No, no era verdad. Era un bulo.}
\end{enumerate}
\end{exoresolu}

\begin{exoresolu}[Exemples résolus]
Énoncé : \esp{¿Quién le explicó el error a Juan?}
\tcblower
Solution détaillée :
\begin{enumerate}
\item \textbf{Repérer} : 2e paragraphe, phrase 5.
\item \textbf{Citer} : \esp{« Su profesora de español le explicó que... »}.
\item \textbf{Reformuler} : \esp{Su profesora de español le explicó el error.}
\end{enumerate}
\end{exoresolu}

\begin{attention}[¿Qué hizo? (Prétérit) vs ¿Qué hacía? (Imparfait) — Piège classique]
Les francophones confondent ces deux questions car le français utilise souvent le même temps (passé composé ou imparfait) pour les deux.

\begin{center}
\begin{tabularx}{\linewidth}{|X|X|X|}
\hline
\rowcolor{popblueL} Question & Temps verbal & Sens / Réponse attendue \\
\hline
\esp{¿Qué \textbf{hizo} Juan?} & Prétérit simple (\emph{pretérito indefinido}) & Action \textbf{ponctuelle, achevée, datée}. \\
& \esp{hizo} (verbe \esp{hacer}) & Réponse : \esp{Compartió la noticia.} (Il a partagé \textbf{une fois}.) \\
\hline
\esp{¿Qué \textbf{hacía} Juan?} & Imparfait (\emph{pretérito imperfecto}) & Action \textbf{habituelle} ou \textbf{en cours} (décor). \\
& \esp{hacía} (verbe \esp{hacer}) & Réponse : \esp{Pasaba mucho tiempo en las redes sociales.} (Il \textbf{passait} son temps \textbf{généralement}.) \\
\hline
\end{tabularx}
\end{center}

\textbf{Règle :} Avant de répondre, identifie le verbe de la question.
\begin{itemize}
\item \esp{hizo, leyó, compartió, esperó, vino, sintió, aprendió} $\rightarrow$ Prétérit $\rightarrow$ Action précise, une fois.
\item \esp{pasaba, era, decía, vivía, leía} $\rightarrow$ Imparfait $\rightarrow$ Habitude, description, action en cours.
\end{itemize}
\end{attention}

\courssub{Questions sur le lexique (synonymes et champ lexical)}

L'examinateur demande de retrouver un mot précis dans le texte. Cherche un \cle{synonyme} ou un mot du même \cle{champ lexical}.

\begin{exoresolu}[Question lexicale type Bac]
Énoncé : \esp{Encuentra en el texto una palabra que signifique « información falsa » y otra que signifique « comprobar ».}
\tcblower
Solution détaillée :
\begin{enumerate}
\item Je cherche le champ lexical de l'information fausse. Je repère : \esp{« La noticia era un \textbf{bulo} »}. $\rightarrow$ \esp{bulo}.
\item Je cherche le sens « vérifier ». Je repère : \esp{« hay que \textbf{verificar} la fuente »} (synonyme de \esp{comprobar} présent dans \esp{« Sin \textbf{comprobar} la fuente »}). $\rightarrow$ \esp{verificar} (ou \esp{comprobar}).
\item \textbf{Rédaction finale :} \esp{Una palabra que significa « información falsa » es « bulo »: « La noticia era un bulo ». Una palabra que significa « comprobar » es « verificar »: « hay que verificar la fuente ».}
\end{enumerate}
\end{exoresolu}

\begin{exercice}[Entraînement lexique]
\begin{enumerate}
\item Trouve dans le texte le synonyme de « quartier » : \esp{\_\_\_\_\_\_\_\_\_\_}.
\item Trouve le mot qui signifie « honteux » : \esp{\_\_\_\_\_\_\_\_\_\_}.
\item Trouve le mot pour « le titre d'un article » : \esp{\_\_\_\_\_\_\_\_\_\_}.
\item Trouve le mot pour « l'origine d'une information » : \esp{\_\_\_\_\_\_\_\_\_\_}.
\item Trouve un synonyme de « vérifier » (autre que \esp{comprobar}) : \esp{\_\_\_\_\_\_\_\_\_\_}.
\end{enumerate}
\end{exercice}

\begin{corrige}[Exercice « Entraînement lexique »]
\begin{enumerate}
\item \esp{el barrio} (l. 2)
\item \esp{avergonzado} (l. 10)
\item \esp{el titular} (l. 4)
\item \esp{la fuente} (l. 6, 12)
\item \esp{verificar} (l. 12)
\end{enumerate}
\end{corrige}

\courssub{Questions d'interprétation (dépasser la lettre du texte)}

La réponse n'est pas écrite mot pour mot ; il faut la \cle{déduire} à partir d'indices textuels.

\begin{exoresolu}[Interprétation résolue]
Énoncé : \esp{¿Por qué la gente creyó la noticia?}
\tcblower
Solution détaillée :
\begin{enumerate}
\item \textbf{Indices textuels :} \esp{« Juan se puso muy contento »} (crédibilité de l'émetteur), \esp{« compartió... con todos sus amigos y su familia »} (proximité, confiance), \esp{« la noticia circuló por todo el barrio y mucha gente la creyó »} (effet viral, absence de vérification).
\item \textbf{Inférence :} Les gens ont cru la nouvelle parce qu'elle venait d'un ami de confiance (Juan), qu'elle était agréable (visite d'une star), et que personne n'a pris le temps de vérifier la source.
\item \textbf{Réponse rédigée :} \esp{La gente creyó la noticia porque Juan, una persona de confianza, la compartió con entusiasmo; además, la noticia era atractiva y nadie verificó la fuente.}
\end{enumerate}
\end{exoresolu}

\begin{exoresolu}[Interprétation résolue]
Énoncé : \esp{¿Qué aprendió Juan?}
\tcblower
Solution détaillée :
\begin{enumerate}
\item \textbf{Preuve textuelle :} \esp{« la información responsable es un deber de todos »} et \esp{« hay que verificar la fuente, comparar varios medios y distinguir los hechos de las opiniones »}.
\item \textbf{Réponse rédigée :} \esp{Juan aprendió que la información responsable es un deber de todos y que antes de compartir una noticia hay que verificar la fuente, comparar varios medios y distinguir hechos de opiniones.}
\end{enumerate}
\end{exoresolu}

\begin{propriete}[Vrai ou Faux : la justification est obligatoire]
Au baccalauréat, une question \esp{Verdadero o falso} ne rapporte des points que si elle est \cle{justifiée par une citation exacte du texte}. Répondre seulement \esp{Falso} ou \esp{Verdadero} donne \textbf{0 point}.

\textbf{Modèle de réponse :}
\esp{Falso. El texto dice: « Sin comprobar la fuente, compartió el titular ». Por lo tanto, Juan no verificó la fuente antes de compartir.}
\end{propriete}

\begin{exoresolu}[Vrai ou Faux justifié]
Énoncé : \esp{Verdadero o falso. Justifica con una cita: « Juan verificó la fuente antes de compartir la noticia ».}
\tcblower
Solution détaillée :
\begin{enumerate}
\item Je cherche le passage sur le partage : \esp{« Sin comprobar la fuente, compartió el titular »}.
\item C'est le contraire de l'affirmation.
\item \textbf{Réponse :} \esp{Falso. Juan no verificó la fuente. El texto dice: « Sin comprobar la fuente, compartió el titular ».}
\end{enumerate}
\end{exoresolu}

\courssub{Distinguer \esp{hecho} et \esp{opinión} : l'esprit critique}

C'est la compétence clé du module. Un \esp{hecho} (fait) est une information \cle{vérifiable} (on peut la prouver vraie ou fausse). Une \esp{opinión} est un \cle{jugement personnel} (subjectif, ni vrai ni faux).

\begin{propriete}[Critères de distinction]
\begin{itemize}
\item \textbf{Hecho (Fait) :} Énoncé constatable, mesurable, datable, vérifiable par des sources externes. Verbes fréquents : \esp{ocurrir, suceder, registrar, confirmar, mostrar, ser} (identité, classification). \textbf{Test :} « Peut-on le prouver par des données ? » $\rightarrow$ Oui = Hecho.
\item \textbf{Opinión (Opinion) :} Jugement de valeur, sentiment, appréciation, hypothèse. Verbes/structures : \esp{creer que, pensar que, parecer, ser} + adjectif subjectif (\esp{bueno, malo, peligroso, sorprendente, serio, importante}), \esp{debería, es necesario que} + subjonctif. \textbf{Test :} « Peut-on être d'accord ou pas ? » $\rightarrow$ Oui = Opinión.
\end{itemize}
\end{propriete}

\begin{center}
\begin{tabularx}{\linewidth}{|X|X|X|}
\hline
\rowcolor{popblueL} Phrase extraite du texte & Classification & Justification \\
\hline
\esp{Juan leyó una noticia.} & \textbf{Hecho} & Action datée (un lunes), constatable. \\
\hline
\esp{Compartió el titular con sus amigos.} & \textbf{Hecho} & Action précise, vérifiable (historique téléphone). \\
\hline
\esp{Nadie vino el sábado.} & \textbf{Hecho} & Fait observable, constatable par témoins. \\
\hline
\esp{La noticia era sorprendente.} & \textbf{Opinión} & \esp{Sorprendente} = jugement subjectif (pour Juan). \\
\hline
\esp{La información responsable es un deber de todos.} & \textbf{Opinión} & Jugement moral / valeur (déontologie). \\
\hline
\esp{El sitio no era un medio serio.} & \textbf{Opinión} & \esp{Serio} = appréciation qualitative. \\
\hline
\end{tabularx}
\end{center}

\begin{popfigure}
\begin{tikzpicture}[node distance=0.3cm]
\node[draw=popblue, fill=popblueL, rounded corners, minimum width=6.5cm, minimum height=1cm, font=\bfseries\large, text=popdark] (h) at (0,0.5) {HECHOS (Vérifiables)};
\node[draw=poporange, fill=poporangeL, rounded corners, minimum width=6.5cm, minimum height=1cm, font=\bfseries\large, text=popdark] (o) at (7.5,0.5) {OPINIONES (Subjectives)};

\node[draw=popline, fill=white, rounded corners, minimum width=6.5cm, minimum height=4.5cm, align=left, anchor=north west, text=popdark] (hb) at (-3.25,-0.2) {
\textbullet\ \esp{Juan leyó una noticia.}\\
\textbullet\ \esp{Compartió el titular.}\\
\textbullet\ \esp{Nadie vino el sábado.}\\
\textbullet\ \esp{La página publicó la noticia.}\\
\textbullet\ \esp{Juan se sintió avergonzado.}\\
\textbullet\ \esp{La profesora explicó la regla.}
};

\node[draw=popline, fill=white, rounded corners, minimum width=6.5cm, minimum height=4.5cm, align=left, anchor=north west, text=popdark] (ob) at (4.25,-0.2) {
\textbullet\ \esp{La noticia era sorprendente.}\\
\textbullet\ \esp{El sitio no era serio.}\\
\textbullet\ \esp{La info responsable es un deber.}\\
\textbullet\ \esp{Es bueno compartir noticias.}\\
\textbullet\ \esp{Las redes son peligrosas.}\\
\textbullet\ \esp{Juan fue imprudente.}
};
\end{tikzpicture}
\legende{Classification des énoncés : à gauche les faits (constatables), à droite les opinions (discutables).}
\end{popfigure}

\begin{attention}[Piège majeur : Un fait n'est pas forcément \emph{vrai} — \esp{Hecho} $\neq$ \esp{Verdad}]
Dans le texte, le \esp{titular} affirmait : \esp{« un famoso jugador venía a visitar el barrio »}.
\begin{itemize}
\item C'est un \textbf{hecho} (un énoncé factuel : « X viendra »). On \textbf{peut} le vérifier (il vient / il ne vient pas).
\item C'est un \textbf{bulo} : ce fait est \textbf{faux}.
\item Ce n'est \textbf{pas une opinion} (ce n'est pas un jugement comme « c'est bien » ou « c'est dangereux »).
\end{itemize}
\cle{Ne confonds jamais :} \esp{Hecho} = énoncé vérifiable (vrai \textbf{ou} faux). \esp{Opinión} = jugement (ni vrai, ni faux). \esp{Bulo} = \textbf{faux fait} présenté comme vrai.
\end{attention}

\courssub{Méthode du commentaire de texte (Bac / Évaluation)}

Le commentaire de texte en 1ère A suit une structure codifiée. Il ne s'agit pas de résumer, mais d'\textbf{analyser}.

\begin{methode}[Commentaire de texte en 4 étapes (20--25 lignes visées)]
\begin{enumerate}
\item \textbf{Introduction (3--4 lignes) :}
      \begin{itemize}
      \item \textbf{Présenter} : Nature du document (article de presse, récit, témoignage), titre (\esp{« ¿Verdad o bulo? »}), auteur (anonyme / manuel), thème (les fake news / responsabilité numérique), contexte (Cameroun, lycée, réseaux sociaux).
      \item \textbf{Problématiser} : Quelle question le texte pose-t-il ? Ex. : \esp{« ¿Cómo evitar la propagación de bulos en redes sociales? »} / « Comment développer une information responsable ? »
      \item \textbf{Annoncer le plan} : \esp{« En primer lugar...; en segundo lugar...; finalmente... »}.
      \end{itemize}
\item \textbf{Développement (2--3 paragraphes argumentés) :}
      \begin{itemize}
      \item \textbf{Axe 1 : Le mécanisme du bulo (Récit).} Analyse de la chaîne : réception $\rightarrow$ émotion (\esp{contento}) $\rightarrow$ partage impulsif (\esp{sin comprobar}) $\rightarrow$ viralité (\esp{circuló por todo el barrio}). Citer le texte.
      \item \textbf{Axe 2 : La leçon critique (Analyse).} Distinction \esp{hecho / opinión}. Rôle de l'école (\esp{profesora}). Outils : \esp{verificar la fuente, contrastar medios}. Vocabulaire clé : \esp{fuente, bulo, medio serio}.
      \item \textbf{Axe 3 (si place) : Portée citoyenne.} \esp{Deber de todos}. Changement de comportement (\esp{ahora... siempre se pregunta}).
      \end{itemize}
\item \textbf{Conclusion (2--3 lignes) :}
      \begin{itemize}
      \item Résumer la thèse du texte : la vigilance est un devoir citoyen.
      \item Ouverture personnelle / contextuelle : lien avec le Cameroun (ex. rumeurs sur la CAN, le cacao, les résultats du Bac), rôle des \esp{influencers}, éducation aux médias (EMI).
      \end{itemize}
\item \textbf{Relecture :} Orthographe (accents, \esp{ñ}, \esp{¿¡}), concordance des temps, connecteurs logiques (\esp{por tanto, sin embargo, además, en conclusión}).
\end{enumerate}
\end{methode}

\begin{exemplebox}[Commentaire modèle (niveau B1 / Bac)]
\textbf{Introducción}
\esp{El texto « ¿Verdad o bulo? » es un relato didáctico que narra la experiencia de Juan, un adolescente de Yaoundé, víctima de una noticia falsa en redes sociales. El tema central es la propagación de los bulos y la necesidad de una información responsable. La pregunta que guía el análisis es: ¿qué mecanismos llevan a creer y compartir una mentira, y cómo formar un espíritu crítico?}

\textbf{Desarrollo}
\esp{En primer lugar, el texto muestra la mecánica del bulo: la noticia « sorprendente » genera una emoción inmediata (Juan « se puso muy contento ») que anula el sentido crítico. La acción « compartir » se vuelve reflejo (« sin comprobar la fuente »), amplificada por la viralidad (« circuló por todo el barrio »). Es un hecho constatable: « decenas de personas esperaron » un jugador que nuncavino.}

\esp{En segundo lugar, la reacción pedagógica de la profesora introduce la solución: la verificación. El texto opone « medio serio » y « sitio para bromas », « hecho » y « opinión ». La distinción es clave: el bulo es un hecho falso, no una opinión. La lección final (« la información responsable es un deber de todos ») transforma el error en norma ética.}

\esp{Finalmente, el cambio de actitud de Juan (« ahora... siempre se pregunta ») ilustra la internalisation de la règle.}

\textbf{Conclusión}
\esp{En conclusión, el texto demuestra que la credulidad digital nace de la emoción y la confianza mal placée. La escuela tiene un papel vital para enseñar a « verificar, contrastar, distinguir ». En Camerún, como en todo el mundo hispanohablante, esta competencia es urgente frente a los bulos electorales, sanitarios o deportivos.}
\end{exemplebox}

\courssub{Outils linguistiques : Exprimer son opinion}

Pour l'oral et l'écrit (commentaire, production), tu dois maîtriser les structures d'opinion. Au niveau A2+/B1, on utilise principalement l'\textbf{indicatif} (certitude/opinion personnelle affirmative) et le \textbf{subjonctif} (doute, négation, impersonnels).

\begin{propriete}[Structures pour exprimer l'opinion]
\begin{center}
\begin{tabularx}{\linewidth}{|X|X|X|}
\hline
\rowcolor{popblueL} Structure & Mode & Exemple traduit \\
\hline
\esp{Creo que} + indicatif & Indicatif & \esp{Creo que es un bulo.} (Je pense que c'est un canular.) \\
\hline
\esp{No creo que} + subjonctif & Subjonctif & \esp{No creo que sea verdad.} (Je ne pense pas que ce soit vrai.) \\
\hline
\esp{Me parece que} + indicatif & Indicatif & \esp{Me parece que la fuente no es fiable.} (Il me semble que la source n'est pas fiable.) \\
\hline
\esp{En mi opinión,} + indicatif & Indicatif & \esp{En mi opinión, hay que verificar.} (À mon avis, il faut vérifier.) \\
\hline
\esp{Desde mi punto de vista,} + indicatif & Indicatif & \esp{Desde mi punto de vista, es un deber.} (De mon point de vue, c'est un devoir.) \\
\hline
\esp{Es importante / necesario / fundamental que} + subjonctif & Subjonctif & \esp{Es fundamental que verifiquemos la fuente.} (Il est fondamental que nous vérifiions la source.) \\
\hline
\esp{Opino que} + indicatif & Indicatif & \esp{Opino que las redes sociales son útiles si se usan bien.} \\
\hline
\end{tabularx}
\end{center}
\cle{Rappel :} Après \esp{creo que, me parece que, en mi opinión, opino que} (affirmatif) $\rightarrow$ \textbf{Indicatif}. Après \esp{no creo que, dudo que, es importante que} $\rightarrow$ \textbf{Subjonctif}.
\end{propriete}

\begin{exercice}[Mets au bon mode (Indicatif ou Subjonctif)]
\begin{enumerate}
\item \esp{Creo que la noticia \_\_\_\_\_\_\_\_ (ser) falsa.}
\item \esp{No creo que el sitio \_\_\_\_\_\_\_\_ (ser) serio.}
\item \esp{Me parece que Juan \_\_\_\_\_\_\_\_ (compartir) el bulo sin querer.}
\item \esp{Es necesario que los jóvenes \_\_\_\_\_\_\_\_ (aprender) a verificar.}
\item \esp{En mi opinión, los medios \_\_\_\_\_\_\_\_ (tener) mucha responsabilidad.}
\end{enumerate}
\end{exercice}

\begin{corrige}[Exercice « Mode »]
\begin{enumerate}
\item \esp{es} (Indicatif : \esp{Creo que} + certitude)
\item \esp{sea} (Subjonctif : \esp{No creo que} + doute)
\item \esp{compartió} (Indicatif : \esp{Me parece que} + fait constatable passé) \textit{ou} \esp{comparta} (Subjonctif si nuance de doute/hypothèse, mais ici fait narré $\rightarrow$ Indicatif).
\item \esp{aprendan} (Subjonctif : \esp{Es necesario que} + impersonnel de nécessité)
\item \esp{tienen} (Indicatif : \esp{En mi opinión} + opinion affirmative)
\end{enumerate}
\end{corrige}

\courssub{Production guidée : Escribir y opinar}

\begin{experience}[Atelier « Chasse aux bulos » — Production écrite / Orale]
\textbf{Consigne :} En binôme. Choisissez un thème d'actualité camerounaise ou sportive (ex. : transfert de footballeur, résultat d'examen, prix du cacao, rumeur politique).
\begin{enumerate}
\item \textbf{Invente} un \esp{bulo} court (2 phrases) sur ce thème. Ex. : \esp{« URGENTE: La FECAFOOT annonce que Samuel Eto'o devient sélectionneur du Real Madrid mañana. »}
\item \textbf{Rédige} un \esp{post de démenti} (style « Fact-checking ») pour alerter tes abonnés. Structure imposée :
      \begin{itemize}
      \item \esp{ALERTA BULO:} [Titre accrocheur]
      \item \esp{Hecho:} [Ce qui est vrai / ce qui s'est passé vraiment]
      \item \esp{Fuente oficial:} [Où as-tu vérifié ? Ex. \esp{Web oficial de la FECAFOOT, cuenta verificada de X}]
      \item \esp{Consejo:} \esp{Antes de compartir, verifica y contrasta. La información responsable es un deber.}
      \end{itemize}
\item \textbf{Présente} ton travail à la classe. Vote pour le démenti le plus clair et le plus complet.
\end{enumerate}
\textbf{Objectif :} Réutiliser \esp{hecho, bulo, fuente, verificar, medio serio, deber, creo que / en mi opinión}. Travailler l'impératif (\esp{verifica, contrasta, no compartas}) et la structure argumentative.
\end{experience}

\begin{exercice}[Production écrite individuelle — Type Bac]
\textbf{Sujet :} \esp{« Escribe un artículo corto (80--100 palabras) para el periódico de tu instituto con el título: "No a los bulos: mi compromiso". Utiliza el vocabulario de la lección y al menos tres structures d'opinion. »}

\textbf{Pistes :}
\begin{itemize}
\item Intro : Définir le bulo, danger pour la société.
\item Corps : Ton expérience (ou celle de Juan), les 3 règles d'or (\esp{verificar, contrastar, distinguir}).
\item Conclusion : Appel à la responsabilité (\esp{deber de todos}), ton engagement personnel (\esp{Me comprometo a...}).
\end{itemize}
\end{exercice}

\begin{savaistu}[Cameroun \& Guinée équatoriale : Médias et infox]
\begin{itemize}
\item \textbf{Guinée équatoriale :} Seul pays africain hispanophone officiel. Principaux journaux : \esp{« La Gaceta de Guinea Ecuatorial »} (officiel), \esp{« Diario Rombe »} (opposition, en ligne). La radio d'État (\esp{Radio Nacional}) et \esp{TVGE} dominent. L'accès à Internet est surveillé ; les \esp{bulos} politiques circulent beaucoup sur \esp{WhatsApp} (groupes familiaux, \esp{« cadenas »}).
\item \textbf{Cameroun :} Pays bilingue (FR/EN), mais l'espagnol est LV2 majeure. Presse francophone/anglophone forte (\esp{Mutations, Cameroon Tribune, The Guardian Post}). \textbf{Fake news locales :} rumeurs sur la santé (Ebola, Covid, vaccins), résultats du Bac/Bepc (listes falsifiées sur Facebook), transferts de \esp{Les Lions Indomptables}, prix du cacao/café. L'initiative \esp{\#StopIntox} (MINPOSTEL) lutte contre la désinformation.
\item \textbf{Étymologie :} \esp{Bulo} vient du catalan \esp{bolo} (boule, mensonge gros comme une boule) ou de l'argot \esp{bulo} = « mentira gorda ». En Espagne, on dit aussi \esp{bulos} pour \emph{fake news}. L'Académie royale (\esp{RAE}) l'accepte depuis 2017.
\end{itemize}
\end{savaistu}

\begin{aretenir}[L'essentiel de la leçon EA17]
\begin{itemize}
\item \textbf{Lexique clé :} \esp{prensa, noticia, periodista, titular, redes sociales, bulo, fuente, hecho, opinión, verificar, contrastar, medio serio}.
\item \textbf{Compréhension :} 3 étapes (Repérer, Citer, Reformuler). Justifier \textbf{toujours} les Vrai/Faux.
\item \textbf{Grammaire :} \esp{¿Qué hizo?} (Prétérit = action ponctuelle) $\neq$ \esp{¿Qué hacía?} (Imparfait = habitude/décor).
\item \textbf{Esprit critique :} \esp{Hecho} = vérifiable (vrai ou faux). \esp{Opinión} = subjectif. \esp{Bulo} = \textbf{faux fait}.
\item \textbf{Commentaire :} Introduction (présenter, problématiser, plan), Développement (2 axes + citations), Conclusion (bilan + ouverture).
\item \textbf{Opinion :} \esp{Creo que / Me parece que / En mi opinión} + \textbf{Indicatif} (affirmatif) ; \esp{No creo que / Es importante que} + \textbf{Subjonctif}.
\item \textbf{Production :} Écrire un démenti structuré (\esp{ALERTA BULO, Hecho, Fuente, Consejo}) ou un article d'opinion.
\end{itemize}
\end{aretenir}

\begin{exercice}[Bilan complet — Entraînement évaluation]
\begin{enumerate}
\item \textbf{Compréhension :} Réponds en phrases complètes : \esp{¿Dónde vive Juan? ¿Qué leyó un lunes? ¿Qué hizo con el titular? ¿Quién le explicó su error?}
\item \textbf{Vrai / Faux justifié :}
      \begin{enumerate}
      \item \esp{« La página que publicó la noticia era un medio serio. »}
      \item \esp{« Juan se sintió orgulloso después del sábado. »}
      \end{enumerate}
\item \textbf{Lexique :} Trouve dans le texte : un synonyme de « comprobar », le mot pour « quartier », le mot pour « honteux ».
\item \textbf{Hecho / Opinión :} Classe avec justification :
      \begin{enumerate}
      \item \esp{Juan pasa mucho tiempo en las redes sociales.}
      \item \esp{Las redes sociales son peligrosas.}
      \item \esp{Decenas de personas esperaron delante del campo.}
      \end{enumerate}
\item \textbf{Grammaire (Opinion) :} Complète avec le mode correct :
      \esp{Es fundamental que los estudiantes \_\_\_\_\_\_\_\_ (aprender) a detectar bulos.}
\item \textbf{Expression écrite :} Rédige 5 lignes : \esp{« En mi opinión, para evitar los bulos hay que... »} (utilise \esp{verificar, contrastar, fuente, medio serio}).
\end{enumerate}
\end{exercice}

\begin{corrige}[Exercice « Bilan complet »]
\begin{enumerate}
\item \esp{Juan vive en Yaoundé. Un lunes leyó una noticia sorprendente sobre un jugador de fútbol. Compartió el titular con sus amigos y su familia sin comprobar la fuente. Su profesora de español le explicó su error.}
\item a) \esp{Falso. El texto dice: « la página que la publicó no era un medio serio, sino un sitio creado para hacer bromas ».} b) \esp{Falso. El texto dice: « Juan se sintió avergonzado ».}
\item « Comprobar » : \esp{verificar} (l. 12) ; « Quartier » : \esp{el barrio} (l. 2) ; « Honteux » : \esp{avergonzado} (l. 10).
\item \begin{enumerate}
      \item \esp{Hecho} (habitude constatable : \esp{pasa}).
      \item \esp{Opinión} (jugement de valeur : \esp{peligrosas}).
      \item \esp{Hecho} (fait observable, daté, quantifié : \esp{decenas, esperaron}).
      \end{enumerate}
\item \esp{aprendan} (Subjonctif présent : \esp{Es fundamental que} + nécessité).
\item \textbf{Production type :} \esp{En mi opinión, para evitar los bulos hay que verificar siempre la fuente antes de compartir. Es fundamental contrastar la información en varios medios serios. No creo que sea bueno reenviar cadenas sin leer. Me comprometo a distinguir hechos de opiniones. La información responsable es un deber de todos.}
\end{enumerate}
\end{corrige}

\coursec{Texte support : « ¿Verdad o bulo? »}

Avant d'étudier le vocabulaire et les outils d'analyse, lisons un article de presse typique. Ce texte original met en scène une situation camerounaise familière : une rumeur sur les réseaux sociaux, sa propagation et le travail de vérification d'un journaliste.

\begin{textebox}[Article original : ¿Verdad o bulo?]
\esp{
¿Verdad o bulo? La prensa ante el reto de la desinformación

En Yaoundé, el pasado martes, un mensaje se hizo viral en WhatsApp y Facebook: « ¡Cierre total del mercado Mfoundi por orden del gobernador! ». El titular alarmaba a los comerciantes y compradores. Miles de personas compartieron la noticia sin comprobar la fuente.

La periodista Marie Ndom, de la redacción de \emph{Le Jour}, decidió verificar el hecho. Llamó a la fuente oficial: la delegación del comercio. La respuesta fue clara: « Es un bulo. El mercado abre con normalidad ». El gobernador desmintió el rumor en un comunicado oficial.

Este caso muestra un problema real: en las redes sociales, la opinión circula más rápido que el hecho. Un bulo genera pánico y pérdidas económicas. El periodista tiene la responsabilidad de contrastar las fuentes antes de publicar. El ciudadano, también: no creer todo, no compartir sin verificar. La verdad es un derecho, la información responsable, un deber.
}
\end{textebox}

\begin{definition}[Traduction et glossaire]
{\em « Vérité ou canular ? La presse face au défi de la désinformation »

À Yaoundé, mardi dernier, un message est devenu viral sur WhatsApp et Facebook : « Fermeture totale du marché Mfoundi sur ordre du gouverneur ! » Le gros titre alarmait les commerçants et les acheteurs. Des milliers de personnes ont partagé la nouvelle sans vérifier la source.

La journaliste Marie Ndom, de la rédaction de \emph{Le Jour}, a décidé de vérifier le fait. Elle a appelé la source officielle : la délégation du commerce. La réponse fut claire : « C'est un canular. Le marché ouvre normalement. » Le gouverneur a démenti la rumeur dans un communiqué officiel.

Ce cas montre un problème réel : sur les réseaux sociaux, l'opinion circule plus vite que le fait. Un canular génère panique et pertes économiques. Le journaliste a la responsabilité de confronter les sources avant de publier. Le citoyen aussi : ne pas tout croire, ne pas partager sans vérifier. La vérité est un droit, l'information responsable, un devoir.}

\textbf{Glossaire :} \esp{el pasado martes} (mardi dernier) ; \esp{se hizo viral} (est devenu viral) ; \esp{alarmaba} (alarmait, imparfait) ; \esp{comprobar} (vérifier) ; \esp{la delegación del comercio} (la délégation du commerce) ; \esp{con normalidad} (normalement) ; \esp{un comunicado oficial} (un communiqué officiel) ; \esp{genera} (génère, présent) ; \esp{pérdidas económicas} (pertes économiques) ; \esp{contrastar} (confronter, croiser) ; \esp{un derecho / un deber} (un droit / un devoir).
\end{definition}

\coursec{Compréhension du texte}

\begin{exercice}[Questions de compréhension]
Répondez en français (ou en espagnol si vous le pouvez) aux questions suivantes :
\begin{enumerate}
\item Quel est le thème général de l'article ?
\item Quelle fausse information a circulé mardi dernier à Yaoundé ?
\item Sur quelles plateformes le message s'est-il propagé ?
\item Qui est Marie Ndom et que fait-elle ?
\item Quelle démarche suit-elle pour vérifier l'information ?
\item Quelle est la réponse de la source officielle ?
\item Selon l'article, qui a la responsabilité de vérifier l'information ? Citez deux acteurs.
\item Relevez dans le texte trois verbes d'action liés au travail du journaliste.
\end{enumerate}
\end{exercice}

\begin{corrige}[Exercice « Compréhension du texte »]
1. Thème : la désinformation (les \esp{bulos}) sur les réseaux sociaux et le rôle de la presse pour la vérifier. 2. « Fermeture totale du marché Mfoundi sur ordre du gouverneur. » 3. WhatsApp et Facebook. 4. Marie Ndom est journaliste à la rédaction de \emph{Le Jour} ; elle vérifie le fait (\esp{verificar el hecho}). 5. Elle appelle la source officielle (la délégation du commerce) → \esp{contrastar las fuentes}. 6. « C'est un bulo. Le marché ouvre normalement. » Le gouverneur dément par communiqué. 7. Le journaliste (\esp{el periodista}) et le citoyen (\esp{el ciudadano}). 8. \esp{verificar}, \esp{llamar} (appeler), \esp{contrastar}, \esp{publicar}, \esp{desmentir} (le gouverneur).
\end{corrige}

\coursec{Lexique thématique : los medios de comunicación}

Dans la section précédente, nous avons lu et compris le texte \esp{¿Verdad o bulo?}. Vous avez rencontré des mots nouveaux : \esp{prensa}, \esp{titular}, \esp{bulo}, \esp{fuente}... C'est le moment de les organiser, de les définir et de les mémoriser. Un bon lexique ne s'apprend pas mot par mot : il s'apprend par \cle{champs lexicaux}, c'est-à-dire par familles de mots qui se rapportent à un même thème. Nous allons donc classer le vocabulaire des médias en cinq champs, puis nous entraîner à l'utiliser.

\courssub{Vue d'ensemble : les cinq champs lexicaux}

Avant d'entrer dans le détail, observons la carte mentale ci-dessous. Elle résume toute la leçon de vocabulaire : au centre, le thème \esp{los medios de comunicación} ; autour, cinq branches colorées, une par champ lexical.

\begin{popfigure}
\begin{tikzpicture}[mindmap, grow cyclic, every node/.style={concept, align=center, font=\small},
  root concept/.append style={concept color=popblue, text=white, font=\small\bfseries},
  level 1/.append style={level distance=3.4cm, sibling angle=72},
  level 1 concept/.append style={font=\small\bfseries}]
\node[root concept, align=center]{Los medios\\de comunicación}
  child[concept color=poporange]{ node[align=center]{La prensa\\y el\\periodismo} }
  child[concept color=popgreen]{ node[align=center]{La\\información} }
  child[concept color=poppurple]{ node[align=center]{Las redes\\sociales} }
  child[concept color=poppink]{ node[align=center]{La desinfor-\\mación} }
  child[concept color=popgold]{ node[align=center]{Verbos\\útiles} };
\end{tikzpicture}
\legende{Carte mentale : les cinq champs lexicaux des médias.}
\end{popfigure}

Retenez cette image mentale : quand vous chercherez un mot, demandez-vous d'abord à quelle branche il appartient.

\courssub{Champ 1 — La presse et le journalisme}

Commençons par l'observation. Lisez ces phrases tirées de la presse hispanophone :

\esp{El periodista escribió un artículo para la redacción del periódico.} « Le journaliste a écrit un article pour la rédaction du journal. »

\esp{El titular de hoy habla del campeonato de fútbol.} « Le gros titre d'aujourd'hui parle du championnat de football. »

\begin{definition}[La presse et le journalisme]
\begin{itemize}
\item \esp{la prensa} : la presse (l'ensemble des journaux et des journalistes). \emph{Attention : toujours au singulier avec l'article défini.} \esp{La prensa informa a los ciudadanos.} « La presse informe les citoyens. »
\item \esp{el periódico} : le journal (quotidien). \esp{Compro el periódico cada mañana.} « J'achète le journal chaque matin. »
\item \esp{el diario} : synonyme courant de \esp{periódico}. \esp{El diario sale a las seis.} « Le journal paraît à six heures. »
\item \esp{la revista} : le magazine, la revue. \esp{Mi hermana lee una revista de moda.} « Ma sœur lit un magazine de mode. »
\item \esp{el artículo} : l'article. \esp{Este artículo es muy interesante.} « Cet article est très intéressant. »
\item \esp{el titular} : le gros titre, le titre principal. \esp{El titular anuncia una buena noticia.} « Le titre annonce une bonne nouvelle. »
\item \esp{el/la periodista} : le/la journaliste.
\item \esp{la redacción} : la rédaction (le lieu et l'équipe). \esp{Trabaja en la redacción de un periódico de Madrid.} « Il travaille dans la rédaction d'un journal de Madrid. »
\end{itemize}
\end{definition}

\begin{attention}[Le genre de « periodista »]
Les noms en \esp{-ista} ont \cle{un seul genre grammatical} : on dit \esp{el periodista} pour l'homme et \esp{la periodista} pour la femme. Ne dites jamais \esp{un periodisto} ! C'est la même règle que pour \esp{el/la artista}, \esp{el/la dentista}, \esp{el/la turista}. Seul l'article change : \esp{María es una periodista famosa.} « María est une journaliste célèbre. »
\end{attention}

\begin{attention}[Faux amis]
\esp{La redacción} ne signifie pas seulement « la rédaction » au sens d'équipe : c'est aussi le texte que l'élève rédige en classe (une « dissertation »). De même, \esp{el titular} est d'abord un \emph{nom} (le gros titre), alors que le français « titulaire » est un adjectif. Ne traduisez pas « le titulaire du poste » par \esp{el titular del puesto} sans vérifier le contexte : en espagnol, \esp{el titular del puesto} existe, mais dans le vocabulaire des médias, \esp{el titular} désigne le titre d'un article.
\end{attention}

\courssub{Champ 2 — L'information : hecho y opinión}

C'est le cœur de la leçon EA17 : pour analyser l'information, il faut distinguer le \cle{fait} de l'\cle{opinion}.

\begin{definition}[L'information]
\begin{itemize}
\item \esp{la noticia} : la nouvelle, l'information. \esp{He escuchado una noticia importante en la radio.} « J'ai entendu une nouvelle importante à la radio. »
\item \esp{el hecho} : le fait (ce qui est réel, vérifiable). \esp{Es un hecho: el partido terminó 2-0.} « C'est un fait : le match s'est terminé 2-0. »
\item \esp{la opinión} : l'opinion (ce que quelqu'un pense). \esp{En mi opinión, fue el mejor partido del año.} « À mon avis, ce fut le meilleur match de l'année. »
\item \esp{la fuente} : la source (d'où vient l'information). \esp{El periodista no revela su fuente.} « Le journaliste ne révèle pas sa source. »
\item \esp{la verdad} : la vérité. \esp{Siempre dice la verdad.} « Il dit toujours la vérité. »
\item \esp{la mentira} : le mensonge. \esp{Esa noticia es una mentira.} « Cette nouvelle est un mensonge. »
\end{itemize}
\end{definition}

\begin{propriete}[Hecho ou opinión ? Le test simple]
Pour savoir si une phrase exprime un fait ou une opinion, posez-vous la question : \emph{peut-on la vérifier ?}
\begin{itemize}
\item \textbf{Hecho} : vérifiable, objectif, souvent avec des chiffres, des dates, des lieux. \esp{Yaoundé es la capital de Camerún.} « Yaoundé est la capitale du Cameroun. »
\item \textbf{Opinión} : personnelle, subjective, souvent introduite par \esp{creo que}, \esp{me parece que}, \esp{en mi opinión}, ou contenant un adjectif de jugement (\esp{bonito}, \esp{terrible}, \esp{mejor}). \esp{Creo que Yaoundé es una ciudad maravillosa.} « Je crois que Yaoundé est une ville merveilleuse. »
\end{itemize}
\end{propriete}

\begin{exemplebox}[Exemples traduits]
\esp{El periodista comprobó la fuente antes de publicar la noticia.} « Le journaliste a vérifié la source avant de publier la nouvelle. »

\esp{Es un hecho que el equipo ganó; es una opinión que jugó bien.} « C'est un fait que l'équipe a gagné ; c'est une opinion qu'elle a bien joué. »

\esp{La verdad siempre sale a la luz.} « La vérité finit toujours par éclater au grand jour. »
\end{exemplebox}

\courssub{Champ 3 — Les réseaux sociaux}

\begin{definition}[Las redes sociales]
\begin{itemize}
\item \esp{las redes sociales} : les réseaux sociaux. \esp{Los jóvenes pasan horas en las redes sociales.} « Les jeunes passent des heures sur les réseaux sociaux. »
\item \esp{el mensaje} : le message. \esp{Recibí un mensaje de mi primo de Douala.} « J'ai reçu un message de mon cousin de Douala. »
\item \esp{compartir} : partager. \esp{No compartas noticias sin verificarlas.} « Ne partage pas de nouvelles sans les vérifier. »
\item \esp{publicar} : publier. \esp{Ella publica fotos cada semana.} « Elle publie des photos chaque semaine. »
\item \esp{el/la seguidor(a)} : l'abonné, le « follower ». \esp{Tiene mil seguidores.} « Il a mille abonnés. »
\item \esp{viral} : viral (invariable en genre : \esp{un vídeo viral}, \esp{una foto viral}). \esp{El vídeo se hizo viral en un día.} « La vidéo est devenue virale en un jour. »
\end{itemize}
\end{definition}

\begin{attention}[Prononciation et orthographe]
\esp{Las redes sociales} se prononce « las ré-dess so-sia-less » : le \esp{d} entre voyelles est doux, presque comme un « d » anglais léger ([ð]). Notez aussi que \esp{compartir} est un verbe régulier en \esp{-ir} : \esp{yo comparto, tú compartes, él comparte}. Ne confondez pas avec le français « comparer » : \esp{compartir} signifie \emph{partager}, jamais « comparer » (qui se dit \esp{comparar}).
\end{attention}

\courssub{Champ 4 — La désinformation et la vérification}

\begin{definition}[Desinformación y verificación]
\begin{itemize}
\item \esp{el bulo} : la fausse information, le canular. \esp{La noticia del cierre del mercado era un bulo.} « La nouvelle de la fermeture du marché était une fausse information. »
\item \esp{la fake news} : anglicisme admis, synonyme de \esp{bulo} (souvent au pluriel : \esp{las fake news}).
\item \esp{comprobar} : vérifier (que quelque chose est vrai). \esp{Comprueba la fecha antes de creerlo.} « Vérifie la date avant d'y croire. »
\item \esp{verificar} : vérifier (synonyme plus formel). \esp{Hay que verificar las fuentes.} « Il faut vérifier les sources. »
\item \esp{contrastar} : confronter, croiser (les sources). \esp{Contrasté la noticia con tres periódicos.} « J'ai confronté la nouvelle avec trois journaux. »
\item \esp{desmentir} : démentir. \esp{El ministro desmintió el bulo.} « Le ministre a démenti la fausse information. »
\end{itemize}
\end{definition}

\begin{attention}[« Una noticia falsa » ou « un bulo » ?]
Dire \esp{una noticia falsa} est grammaticalement correct, et tout hispanophone vous comprendra. Mais dans la presse espagnole et latino-américaine, le mot le plus naturel est \esp{un bulo}. L'anglicisme \esp{fake news} est également admis et très répandu, surtout chez les jeunes. Dans une rédaction, préférez donc : \esp{El rumor era un bulo.} « La rumeur était une fausse information. »
\end{attention}

\begin{savaistu}[Le verbe « desmentir » est irrégulier]
\esp{Desmentir} suit la même conjugaison que \esp{mentir} : l'alternance \cle{e $\rightarrow$ ie} au présent de l'indicatif et aux trois personnes du singulier + 3e pluriel du subjonctif présent. 

\begin{center}
\begin{tabular}{|l|l|l|}\hline
\rowcolor{popblueL} \textbf{Personne} & \textbf{Présent indicatif} & \textbf{Passé simple (indefinido)} \\ \hline
yo & desmiento & desmentí \\ \hline
tú & desmientes & desmentiste \\ \hline
él/ella/usted & desmiente & \textbf{desmintió} \\ \hline
nosotros/as & desmentimos & desmentimos \\ \hline
vosotros/as & desmentís & desmentisteis \\ \hline
ellos/ellas/ustedes & desmienten & desmintieron \\ \hline
\end{tabular}
\end{center}

Exemple : \esp{Yo desmiento esa noticia.} « Je démens cette nouvelle. » Au passé simple, attention à la 3e personne : \esp{él desmintió} (e $\rightarrow$ i). Le nom correspondant est \esp{el desmentido} : \esp{El periódico publicó un desmentido.} « Le journal a publié un démenti. »
\end{savaistu}

\courssub{Champ 5 — Les verbes utiles}

\begin{definition}[Verbos útiles]
\begin{itemize}
\item \esp{informar} : informer. \esp{La radio informa cada hora.} « La radio informe chaque heure. » Avec \esp{de/sobre} : \esp{El artículo informa sobre el clima.} « L'article informe sur le climat. »
\item \esp{leer} : lire. \esp{Leo el periódico los domingos.} « Je lis le journal le dimanche. »
\item \esp{creer} : croire. \esp{No creo todo lo que leo.} « Je ne crois pas tout ce que je lis. »
\item \esp{pensar} (e $\rightarrow$ ie) : penser. \esp{Pienso que la prensa es importante.} « Je pense que la presse est importante. »
\item \esp{opinar} : donner son avis. \esp{Todos opinan sobre el partido.} « Tout le monde donne son avis sur le match. »
\item \esp{difundir} : diffuser, répandre. \esp{No difundas bulos.} « Ne diffuse pas de fausses informations. »
\end{itemize}
\end{definition}

\begin{attention}[Piège de traduction : « informer »]
Le français « s'informer » se traduit par le verbe pronominal \esp{informarse} : \esp{Me informo en Internet.} « Je m'informe sur Internet. » N'oubliez pas le pronom réfléchi ! Et ne confondez pas \esp{informar} (informer) avec \esp{formar} (former, constituer) : \esp{El profesor forma a los alumnos} « Le professeur forme les élèves ».
\end{attention}

\courssub{Tableau récapitulatif du lexique}

\begin{textebox}[Vocabulario esencial : español — français]
\begin{center}
\begin{tabularx}{\linewidth}{|l|X|X|}\hline
\rowcolor{popblueL} \textbf{Español} & \textbf{Français} & \textbf{Exemple} \\ \hline
la prensa & la presse & La prensa informa. \\ \hline
el titular & le gros titre & El titular es sorprendente. \\ \hline
el/la periodista & le/la journaliste & La periodista escribe bien. \\ \hline
la fuente & la source & La fuente es fiable. \\ \hline
el hecho & le fait & Es un hecho verificable. \\ \hline
la opinión & l'opinion & Es solo una opinión. \\ \hline
el bulo & la fausse information & Era un bulo. \\ \hline
comprobar & vérifier & Comprueba la noticia. \\ \hline
compartir & partager & No compartas bulos. \\ \hline
desmentir & démentir & Desmiento el rumor. \\ \hline
\end{tabularx}
\end{center}
\end{textebox}

\textbf{Prononciation et genre : points de vigilance.}\par

\begin{itemize}
\item \esp{el titular} : masculin, se prononce « el ti-tou-lar » (le \esp{u} se prononce [u] après \esp{t}).
\item \esp{la fuente} : féminin, « la fouèn-té » ; le \esp{ue} est une diphtongue [we] en une seule syllabe.
\item \esp{el bulo} : masculin, « él bou-lo » ; le \esp{b} initial est doux [β].
\item \esp{la noticia} : féminin, accent sur le \textbf{i} (\esp{noticia}) ; en Espagne, le \esp{c} devant \esp{i} se prononce [θ] (« no-ti-thia »), en Amérique latine et en Guinée équatoriale, [s] (« no-ti-sia »). Les deux sont correctes.
\item \esp{el/la periodista} : accent sur le \textbf{i} (\esp{periodista}) : pro-paroxytone terminé par voyelle $\rightarrow$ accent écrit obligatoire.
\end{itemize}

\courssub{Exercices de fixation}

\begin{exoresolu}[Association mot / définition]
Énoncé. Reliez chaque mot espagnol à sa définition française : 1. \esp{el bulo} ; 2. \esp{la fuente} ; 3. \esp{el titular} ; 4. \esp{el seguidor} ; 5. \esp{la redacción}. Définitions : a. le gros titre d'un article ; b. l'origine d'une information ; c. une fausse information ; d. l'équipe qui prépare un journal ; e. un abonné sur les réseaux sociaux.
\tcblower
Solution détaillée. 1-c : \esp{el bulo} est la fausse information (le mot préféré des médias espagnols). 2-b : \esp{la fuente} est la source, l'origine de l'information. 3-a : \esp{el titular} est le gros titre. 4-e : \esp{el seguidor} est l'abonné (le « follower »). 5-d : \esp{la redacción} est la rédaction, l'équipe du journal. Astuce de mémorisation : \esp{fuente} vient du latin \emph{fons} (la fontaine) : la source est « la fontaine » d'où coule l'information.
\end{exoresolu}

\begin{exercice}[Phrases à compléter]
Complétez avec le mot qui convient : \esp{bulo}, \esp{fuente}, \esp{titular}, \esp{opinión}, \esp{compartir}, \esp{periodista}.
\begin{enumerate}
\item La \ldots\ldots escribió un artículo sobre el mercado de Yaoundé.
\item Antes de \ldots\ldots una noticia, hay que verificarla.
\item El \ldots\ldots del periódico dice: «El equipo nacional gana».
\item Esa información es falsa: es un \ldots\ldots.
\item El periodista protege el nombre de su \ldots\ldots.
\item «Me parece que el partido fue aburrido» es una \ldots\ldots.
\end{enumerate}
\end{exercice}

\begin{corrige}[Exercice « Phrases à compléter »]
1. \esp{periodista} (la journaliste, avec l'article \esp{la} $\rightarrow$ \esp{la periodista}). 2. \esp{compartir} (infinitif après la préposition \esp{antes de}). 3. \esp{titular} (le gros titre annonce le résultat). 4. \esp{bulo} (une information fausse). 5. \esp{fuente} (le journaliste protège sa source). 6. \esp{opinión} (la phrase commence par \esp{me parece que}, marqueur de subjectivité).
\end{corrige}

\begin{exercice}[Classement : ¿hecho u opinión?]
Classez ces phrases dans deux colonnes : \esp{hecho} ou \esp{opinión}.
\begin{enumerate}
\item \esp{El partido terminó con el resultado de 2-0.}
\item \esp{Creo que fue el mejor partido del campeonato.}
\item \esp{La redacción del periódico está en Madrid.}
\item \esp{En mi opinión, las redes sociales son peligrosas.}
\item \esp{El vídeo tiene un millón de visitas.}
\end{enumerate}
\end{exercice}

\begin{corrige}[Exercice « Classement »]
Hechos : phrases 1, 3 et 5 (vérifiables : un score, un lieu, un nombre de vues). Opiniones : phrases 2 et 4 (introduites par \esp{creo que} et \esp{en mi opinión}, avec des jugements de valeur : \esp{mejor}, \esp{peligrosas}).
\end{corrige}

\begin{aretenir}[L'essentiel du lexique]
\begin{itemize}
\item Cinq champs : la presse (\esp{prensa, titular, periodista}), l'information (\esp{noticia, hecho, opinión, fuente}), les réseaux (\esp{redes sociales, compartir, viral}), la désinformation (\esp{bulo, comprobar, desmentir}), les verbes (\esp{informar, creer, opinar, difundir}).
\item \esp{el/la periodista} : un seul mot pour les deux sexes (accent sur le \textbf{i}).
\item \esp{la noticia} : accent sur le \textbf{i}.
\item \esp{un bulo} est le mot naturel pour « une fausse information ».
\item \esp{desmentir} : irrégulier, \esp{yo desmiento} / \esp{él desmintió}.
\item Le test de l'esprit critique : vérifiable = \esp{hecho} ; personnel = \esp{opinión}.
\end{itemize}
\end{aretenir}

\coursec{Distinguer hecho y opinión — L'esprit critique en action}

Le lexique est acquis ; place maintenant à la compétence visée par le programme : \cle{analyser l'information} pour développer l'esprit critique. Dans un monde saturé d'images et de messages courts (TikTok, X, WhatsApp), la frontière entre fait et opinion s'efface. Votre mission : la rétablir.

\courssub{Les marqueurs linguistiques}

\begin{propriete}[Repérer l'opinion : les signaux grammaticaux]
L'opinion en espagnol se signale souvent par des structures subjectives :
\begin{itemize}
\item \textbf{Verbes d'opinion + indicatif} (certitude personnelle) : \esp{creer que}, \esp{pensar que}, \esp{opinar que}, \esp{considerar que}. \esp{Creo que el bulo es peligroso.}
\item \textbf{Verbes d'opinion + subjonctif} (doute, évaluation, sentiment) : \esp{me parece que} + \textbf{subjonctif}, \esp{es bueno/malo que} + \textbf{subjonctif}, \esp{es importante que} + \textbf{subjonctif}. \esp{Me parece que las redes sociales sean un riesgo.} (Subjonctif \esp{sean} car évaluation subjective).
\item \textbf{Adverbes et locutions} : \esp{en mi opinión}, \esp{a mi juicio}, \esp{según yo}, \esp{probablemente}, \esp{seguramente}.
\item \textbf{Adjectifs de jugement} : \esp{mejor, peor, terrible, maravilloso, injusto, necesario}.
\end{itemize}
\end{propriete}

\begin{propriete}[Repérer le fait : les signaux objectifs]
Le fait s'exprime généralement à l'\textbf{indicatif} (présent, passé composé, imparfait, futur) sans marqueur de subjectivité :
\begin{itemize}
\item Verbes d'action ou d'état constaté : \esp{ocurrir, suceder, registrar, confirmar, anunciar, terminar, ganar, perder}.
\item Données chiffrées, dates, lieux précis : \esp{El 12 de octubre}, \esp{el 75\%}, \esp{en Douala}.
\item Verbes de communication neutres : \esp{informar que}, \esp{anunciar que}, \esp{confirmar que} + \textbf{indicatif}. \esp{La prensa informó que el mercado abrió.}
\end{itemize}
\end{propriete}

\begin{exemplebox}[Analyse contrastée : même sujet, deux traitements]
\begin{tabularx}{\linewidth}{|X|X|}\hline
\rowcolor{popblueL} \textbf{HECHO (Article factuel)} & \textbf{OPINIÓN (Éditorial / Tweet)} \\ \hline
\esp{El mercado Mfoundi abrió sus puertas a las 6h.} & \esp{Es una vergüenza que el mercado abra tan temprano.} \\ \hline
\esp{La delegación confirmó la normalidad.} & \esp{Me parece inaceptable la gestión del gobernador.} \\ \hline
\esp{El bulo circuló en 500 grupos de WhatsApp.} & \esp{Creo que WhatsApp debería prohibir esos mensajes.} \\ \hline
\end{tabularx}
\end{exemplebox}

\courssub{La vérification en 3 étapes (Méthode « CCC »)} 

\begin{methode}[Comment vérifier une information : Contraster, Comprobar, Confirmar]
\begin{enumerate}
\item \textbf{Contrastar (Croiser) :} Cherchez la même nouvelle dans \textbf{au moins 3 sources fiables} (presse écrite reconnue, site officiel, agence de presse). \esp{Contrastar la noticia con El País, RFI, y la web del gobierno.}
\item \textbf{Comprobar (Vérifier les détails) :} Regardez la \textbf{date}, l'\textbf{auteur}, l'\textbf{URL} (site officiel ? \esp{.gob.es}, \esp{.gov.cm}), les \textbf{images} (recherche inversée Google/TinEye). \esp{Comprobar la fecha: la foto es de 2018, no de hoy.}
\item \textbf{Confirmar (Valider) :} Attendez le \textbf{démenti ou la confirmation officielle}. Ne partagez qu'après. \esp{El ministro confirmó / desmintió la versión.}
\end{enumerate}
\end{methode}

\begin{experience}[Atelier « Détectives de l'info »]
\textbf{Consigne :} En groupes de 4. Chaque groupe reçoit 3 captures d'écran (imprimées ou projetées) : une vraie nouvelle (ex. résultat CAN), un bulo classique (ex. « Coca-Cola ferme son usine à Douala »), une opinion déguisée en fait (ex. « Le nouveau ministre est incompétent, le pays va mal »).
\textbf{Tâche :} 1. Classer : \esp{Hecho / Bulo / Opinión}. 2. Justifier avec les marqueurs vus ci-dessus. 3. Appliquer la méthode CCC. 4. Présenter à la classe en espagnol : \esp{« Es un bulo porque... La fuente no es oficial... No hay datos verificables... »}.
\end{experience}

\coursec{Méthode du commentaire de texte et commentaire modèle}

Le commentaire de texte est une épreuve majeure du Baccalauréat (épreuve écrite). Il ne s'agit pas de résumer, mais d'\textbf{analyser} : comment le texte construit-il son sens ? Voici la méthode en 4 étapes, appliquée à notre texte \esp{¿Verdad o bulo?}.

\begin{methode}[Étapes du commentaire de texte (Type « Texte de presse »)]
\begin{enumerate}
\item \textbf{Lecture active \& Contextualisation} (10 min) : Identifier le \textbf{genre} (article d'opinion / reportage / éditorial), la \textbf{source} (journal, auteur, date), le \textbf{thème} central, le \textbf{ton} (alarmiste, pédagogique, ironique...). Noter les mots-clés récurrents.
\item \textbf{Structure \& Organisation} (10 min) : Découper en \textbf{paragraphes logiques} (introduction, développement, conclusion). Repérer les \textbf{connecteurs} (\esp{por tanto, sin embargo, además, en conclusión}). Tracer le schéma : \esp{Planteamiento $\rightarrow$ Caso concreto $\rightarrow$ Análisis $\rightarrow$ Llamamiento}.
\item \textbf{Analyse du contenu : Hechos vs Opiniones} (15 min) : Surligner les \textbf{faits vérifiables} (couleur verte) et les \textbf{opinions/jugements} (couleur orange). Identifier la \textbf{thèse} de l'auteur. Relever le \textbf{vocabulaire thématique} (champs lexicaux) et les \textbf{figures de style} (métaphores, répétitions).
\item \textbf{Rédaction structurée} (20 min) : 
\begin{itemize}
    \item \textbf{Introduction} : Accroche $\rightarrow$ Présentation du texte (titre, auteur, source, date, thème) $\rightarrow$ Problématique $\rightarrow$ Annonce du plan.
    \item \textbf{Développement} (2-3 parties) : Chaque partie = une idée force + citation courte + analyse (lexique, grammaire, argumentation). \emph{Ex: « Le texte dénonce la viralité des bulos... »}.
    \item \textbf{Conclusion} : Bilan synthétique $\rightarrow$ Réponse à la problématique $\rightarrow$ Ouverture (lien avec l'actualité camerounaise, rôle du citoyen).
\end{itemize}
\item \textbf{Relecture} (5 min) : Orthographe, accents, connecteurs, temps verbaux, respect de la consigne.
\end{enumerate}
\end{methode}

\begin{exemplebox}[Commentaire modèle (extrait — Introduction + Partie 1)]
\textbf{Introducción}
En la era digital, la frontera entre información y desinformación se ha vuelto difusa. El artículo « ¿Verdad o bulo? », publicado en un contexto camerunés, aborda este desafío a través de un caso concreto: la viralización de un falso cierre de mercado en Yaoundé. \textbf{Problématique :} ¿Cómo puede la prensa, apoyada en la verificación, combatir la propagación de bulos en redes sociales? \textbf{Plan:} 1. La anatomía de un bulo viral; 2. El método periodístico como antídoto; 3. La responsabilidad compartida del ciudadano.

\textbf{Desarrollo 1: La anatomía de un bulo viral}
El texto expone la mecánica de la desinformación: un \esp{mensaje} alarmista (\esp{« ¡Cierre total...! »}) circula en \esp{WhatsApp y Facebook} sin \esp{fuente} verificable. El vocabulario del \esp{Champ 3} (\esp{viral, compartir, mensaje}) y del \esp{Champ 4} (\esp{bulo, pánico}) muestra la velocidad y el impacto emocional. La estructura \esp{« Miles de personas compartieron... sin comprobar »} (indicativo de hecho + gerundio de manera) denuncia la irresponsabilidad colectiva. Es un \esp{hecho} constatable: la opinión (el miedo) viaja más rápido que el \esp{hecho} (la realidad del mercado abierto).
\end{exemplebox}

\coursec{Production guidada : escribir y opinar}

\courssub{Expresión de la opinión : estructuras clave}

Pour l'oral (exposé, débat) et l'écrit (article, lettre, commentaire), le programme exige l'usage correct des marqueurs d'opinion.

\begin{propriete}[Structures pour exprimer son avis (Niveau A2+/B1)]
\begin{center}
\begin{tabularx}{\linewidth}{|l|X|X|}\hline
\rowcolor{popblueL} \textbf{Estructura} & \textbf{Registro / Nuance} & \textbf{Ejemplo} \\ \hline
\esp{Creo que} + \textbf{Indicatif} & Opinion personnelle, certitude & \esp{Creo que la prensa es libre.} \\ \hline
\esp{No creo que} + \textbf{Subjonctif} & Doute / Négation de l'opinion & \esp{No creo que sea verdad.} \\ \hline
\esp{Me parece que} + \textbf{Indicatif/Subjonctif} & Impression / Jugement subjectif & \esp{Me parece que es un bulo.} (Ind.) / \esp{Me parece bien que verifiquen.} (Subj.) \\ \hline
\esp{En mi opinión,} + \textbf{Indicatif} & Formel, structuré & \esp{En mi opinión, el periodista debe verificar.} \\ \hline
\esp{A mi juicio,} + \textbf{Indicatif} & Plus formel, écrit & \esp{A mi juicio, la ley es necesaria.} \\ \hline
\esp{Desde mi punto de vista,} & Très formel, essai & \esp{Desde mi punto de vista, la educación mediática es clave.} \\ \hline
\esp{Es evidente que} + \textbf{Indicatif} & Certitude partagée (fait social) & \esp{Es evidente que las redes cambian la información.} \\ \hline
\end{tabularx}
\end{center}
\end{propriete}

\begin{attention}[Piège majeur : \esp{Creo que} vs \esp{No creo que}]
\esp{Creo que} (affirmatif) $\rightarrow$ \textbf{Indicatif}. \esp{Creo que es verdad.}
\esp{No creo que} (négatif) $\rightarrow$ \textbf{Subjonctif}. \esp{No creo que sea verdad.}
\emph{Raison :} La négation introduit le doute/la subjectivité $\rightarrow$ subjonctif. C'est une règle d'or du BAC.
\end{attention}

\courssub{Production écrite : « Mi artículo de verificación »}

\begin{exercice}[Rédaction guidée : Écrire un article de fact-checking (120-150 mots)]
\textbf{Consigne :} Vous êtes journaliste pour le journal scolaire \esp{« La Voz del Liceo »}. Écrivez un court article pour démentir une rumeur circulant au lycée (ex. : « Le bac est annulé », « La cantine sert de la viande de chien », « Le prof d'espagnol part demain »).
\textbf{Contraintes :}
\begin{enumerate}
\item \textbf{Titular} accrocheur (max 8 mots).
\item \textbf{Entrada (Lead) :} Qui ? Quoi ? Où ? Quand ? Source officielle.
\item \textbf{Cuerpo :} Expliquer l'origine du bulo, la vérification (méthode CCC), la réponse officielle.
\item \textbf{Conclusión :} Appel à la responsabilité (\esp{No compartan sin verificar}).
\item \textbf{Vocabulaire imposé :} \esp{bulo, fuente, verificar, periodista, redacción, opinión, hecho, compartir, desmentir, verdad}.
\item \textbf{Grammaire :} Au moins 2 opinions avec \esp{creo que / me parece que / en mi opinión} + indicatif/subjonctif correct.
\end{enumerate}
\end{exercice}

\begin{exemplebox}[Modèle de production (Guide)]
\esp{Titular: « Bulo: No hay suspensión de clases el lunes »}

\esp{Entrada:} \esp{En la redacción de « La Voz del Liceo » hemos recibido decenas de consultas sobre un supuesto comunicado del MINESEC que anuncia la suspensión de clases el próximo lunes.}

\esp{Cuerpo:} \esp{El mensaje, viral en WhatsApp, carece de firma oficial y de número de decreto. \textbf{He comprobado} la fuente en la web del Ministerio: no existe tal comunicado. El secretario general \textbf{ha desmentido} el rumor esta mañana. \textbf{En mi opinión,} es un bulo clásico de fin de trimestre para crear confusión. \textbf{Creo que} los estudiantes deben actuar con responsabilidad.}

\esp{Conclusión:} \esp{La verdad es un derecho. \textbf{No compartan} informaciones sin \textbf{verificar}. El periodismo serio contrasta; el bulo destruye.}
\end{exemplebox}

\courssub{Production orale : Débat « ¿Regulación o libertad?»}

\begin{experience}[Débat structuré : « Les réseaux sociaux doivent-ils être régulés par l'État ? »]
\textbf{Préparation (10 min) :} En binômes. L'un défend la régulation (lutte contre bulos, haine, protection mineurs), l'autre la liberté totale (censure, droit à l'info, autorégulation). Chaque élève prépare 3 arguments avec \esp{creo que, me parece que, en mi opinión, es evidente que}.
\textbf{Déroulement (15 min) :} 
\begin{enumerate}
\item Tour 1 : Argument 1 (Pro) vs Argument 1 (Contra). Réplique courte.
\item Tour 2 : Argument 2. Réplique.
\item Tour 3 : Argument 3. Réplique.
\end{enumerate}
\textbf{Synthèse (5 min) :} La classe vote. Le professeur note au tableau les meilleures structures d'opinion entendues. Correction collective des erreurs de subjonctif (\esp{no creo que + subj.}).
\end{experience}

\begin{aretenir}[Bilan de la leçon EA17 — Compétences validées]
\begin{itemize}
\item \textbf{Lexique :} Maîtrise des 5 champs (presse, info, réseaux, désinfo, verbes) + accents clés (\esp{periodista, noticia}).
\item \textbf{Grammaire :} Distinction \esp{hecho} (indicatif) / \esp{opinión} (subjonctif après négation/doute) ; \esp{desmentir} (irrégulier).
\item \textbf{Compréhension :} Lire un article de presse, identifier thèse, structure, marqueurs fait/opinion.
\item \textbf{Esprit critique :} Appliquer la méthode CCC (Contrastar, Comprobar, Confirmar) face à un message viral.
\item \textbf{Expression écrite :} Rédiger un article de fact-checking structuré avec vocabulaire imposé.
\item \textbf{Expression orale :} Débattre en utilisant les marqueurs d'opinion corrects (indicatif/subjonctif).
\end{itemize}
\end{aretenir}

\coursec{Distinguer hecho y opinión — l'esprit critique}

Dans la section précédente, nous avons découvert le vocabulaire des médias : \esp{prensa}, \esp{noticia}, \esp{periodista}, \esp{titular}, \esp{redes sociales}, \esp{bulo}, \esp{fuente}, \esp{hecho}, \esp{opinión}. Deux de ces mots sont la clé de toute lecture intelligente de l'information : \esp{hecho} (le fait) et \esp{opinión} (l'opinion). Un bon lecteur de presse, à Yaoundé comme à Madrid, doit savoir les distinguer en une fraction de seconde. C'est exactement la compétence que nous allons construire maintenant, étape par étape : observer, définir, repérer les indices, puis s'entraîner sur un texte authentique.

\courssub{Observer d'abord : deux phrases, deux réalités}

Lisons attentivement ces deux phrases, qui parlent du même sujet :

\esp{El examen de español es el lunes a las ocho.} « L'examen d'espagnol a lieu lundi à huit heures. »

\esp{El examen de español es muy difícil.} « L'examen d'espagnol est très difficile. »

Quelle différence voyez-vous ? La première phrase peut être \cle{vérifiée} : il suffit de regarder l'emploi du temps affiché au tableau du lycée. Elle est vraie ou fausse, point final. La seconde phrase ne peut pas être vérifiée de la même manière : « très difficile » dépend de la personne qui parle. Pour un élève qui travaille beaucoup, l'examen sera facile ; pour un autre, il sera terrible. C'est un \cle{jugement personnel}.

\begin{definition}[Hecho y opinión]
Un \cle{hecho} (fait) est une information \textbf{vérifiable et objective} : on peut la confirmer par des preuves, des documents, des chiffres, des témoins. \esp{El examen es el lunes.}

Une \cle{opinión} (opinion) est un \textbf{jugement personnel} : elle exprime ce que quelqu'un pense, ressent ou apprécie. Elle ne peut pas être prouvée vraie ou fausse. \esp{El examen es muy difícil.}

Attention au vocabulaire : \esp{hecho} est aussi le participe passé du verbe \esp{hacer} (faire) : \esp{el trabajo está hecho} « le travail est fait ». Dans le vocabulaire des médias, c'est un nom masculin : \esp{el hecho}, \esp{los hechos}.
\end{definition}

\begin{exemplebox}[Hecho o opinión : entraînement rapide]
\begin{itemize}
\item \esp{El partido terminó 2-0.} « Le match s'est terminé 2-0. » $\rightarrow$ \textbf{hecho} (le score se vérifie).
\item \esp{Fue el mejor partido del año.} « Ce fut le meilleur match de l'année. » $\rightarrow$ \textbf{opinión} (jugement).
\item \esp{La reunión empieza a las diez.} « La réunion commence à dix heures. » $\rightarrow$ \textbf{hecho}.
\item \esp{Las reuniones son aburridas.} « Les réunions sont ennuyeuses. » $\rightarrow$ \textbf{opinión}.
\end{itemize}
\end{exemplebox}

\courssub{Les indices linguistiques du fait}

Un \esp{hecho} se reconnaît à des signes précis dans la phrase. Apprends à les repérer comme un détective.

\begin{propriete}[Indices d'un fait]
Une phrase factuelle contient généralement :
\begin{itemize}
\item des \textbf{chiffres et des données} : \esp{2-0}, \esp{500 alumnos}, \esp{el 60 \%};
\item des \textbf{dates et des heures} : \esp{el lunes}, \esp{a las ocho}, \esp{en 2024};
\item des \textbf{lieux précis} : \esp{en Yaoundé}, \esp{en el estadio};
\item des \textbf{verbes objectifs} qui rapportent sans juger : \esp{decir} (dire), \esp{anunciar} (annoncer), \esp{ocurrir} (se produire), \esp{informar} (informer), \esp{confirmar} (confirmer);
\item la marque de source \esp{según} (selon) : \esp{Según el Ministerio, las clases empiezan el lunes.} « Selon le Ministère, les cours commencent lundi. »
\end{itemize}
Exemple : \esp{El ministro anunció ayer en Douala la apertura de un nuevo hospital.} « Le ministre a annoncé hier à Douala l'ouverture d'un nouvel hôpital. » (verbe objectif + date + lieu : fait.)
\end{propriete}

\courssub{Les indices linguistiques de l'opinion}

L'opinion, elle aussi, laisse des traces très visibles dans la langue.

\begin{propriete}[Indices d'une opinion]
Une phrase d'opinion contient généralement :
\begin{itemize}
\item des \textbf{expressions personnelles} : \esp{creo que} (je crois que), \esp{pienso que} (je pense que), \esp{en mi opinión} (à mon avis), \esp{me parece que} (il me semble que), \esp{para mí} (pour moi);
\item des \textbf{adjectifs évaluatifs} : \esp{bueno} (bon), \esp{malo} (mauvais), \esp{terrible} (terrible), \esp{sorprendente} (surprenant), \esp{increíble} (incroyable), \esp{estupendo} (formidable);
\item des \textbf{verbes de jugement} : \esp{gustar} (plaire), \esp{encantar} (adorer), \esp{odiar} (détester).
\end{itemize}
Exemple : \esp{En mi opinión, las redes sociales son útiles.} « À mon avis, les réseaux sociaux sont utiles. »
\end{propriete}

\begin{center}
\begin{tabularx}{\linewidth}{|X|X|X|}
\hline
\rowcolor{popblueL} Hecho (fait) & Opinión (opinion) & Remarque \\
\hline
\esp{El examen es el lunes.} & \esp{El examen es difícil.} & date vérifiable / adjectif évaluatif \\
\hline
\esp{El partido terminó 2-0.} & \esp{Fue un partido terrible.} & chiffre objectif / jugement \\
\hline
\esp{La ministra anunció la medida.} & \esp{Creo que la medida es buena.} & verbe objectif / \esp{creo que} \\
\hline
\esp{Llueve en Yaoundé.} & \esp{Me parece que llueve demasiado.} & observation / appréciation \\
\hline
\end{tabularx}
\end{center}

\begin{popfigure}
\begin{tikzpicture}[node distance=0.4cm]
\node[draw=popblue, fill=popblueL, rounded corners, align=center, minimum width=6.2cm, font=\small\bfseries] (h) {HECHO (fait)};
\node[draw=poppink, fill=poppinkL, rounded corners, align=center, minimum width=6.2cm, font=\small\bfseries, right=0.8cm of h] (o) {OPINIÓN (opinion)};
\node[draw=popline, fill=popcream, rounded corners, align=center, font=\small, below=0.35cm of h, minimum width=6.2cm] (h1) {verificable, objetivo\\datos, cifras, fechas, lugares\\decir, anunciar, ocurrir};
\node[draw=popline, fill=popcream, rounded corners, align=center, font=\small, below=0.35cm of o, minimum width=6.2cm] (o1) {juicio personal, subjetivo\\creo que, me parece que\\adjetivos: bueno, terrible};
\end{tikzpicture}
\legende{Les deux visages de l'information : le fait se vérifie, l'opinion se discute.}
\end{popfigure}

\courssub{Exprimer son opinion : creo que + indicatif}

Pour donner ton avis à l'écrit comme à l'oral, l'espagnol utilise des verbes d'opinion suivis de \esp{que}.

\begin{propriete}[Verbe d'opinion affirmatif + indicatif]
Après \esp{creo que}, \esp{pienso que}, \esp{me parece que}, \esp{opino que} à la forme \textbf{affirmative}, le verbe de la subordonnée est à l'\textbf{indicatif}, comme en français.

\esp{Creo que la noticia es falsa.} « Je crois que la nouvelle est fausse. »

\esp{Pienso que los periodistas trabajan mucho.} « Je pense que les journalistes travaillent beaucoup. »

\esp{Me parece que este titular exagera.} « Il me semble que ce titre exagère. »
\end{propriete}

\begin{propriete}[Complément utile au Bac : no creo que + subjonctif]
Après \esp{no creo que}, \esp{no pienso que}, \esp{no me parece que} à la forme \textbf{négative}, l'espagnol \textbf{exige le subjonctif}, car le locuteur exprime un doute.

\esp{No creo que sea verdad.} « Je ne crois pas que ce soit vrai. »

\esp{No creo que este titular sea verdad.} « Je ne crois pas que ce titre soit vrai. »

\esp{No me parece que la fuente sea fiable.} « Il ne me semble pas que la source soit fiable. »

Comparaison : le français, lui, hésite librement entre indicatif et subjonctif (« je ne crois pas qu'il vient / qu'il vienne ») ; l'espagnol n'a pas le choix : \textbf{subjonctif obligatoire}.

Rappel de formation : subjonctif présent de \esp{ser} : \esp{sea, seas, sea, seamos, seáis, sean}.
\end{propriete}

\begin{attention}[Le piège de « me parece que »]
\esp{Me parece que} ne s'accorde \textbf{jamais} avec le sujet de la subordonnée : le verbe \esp{parecer} reste à la troisième personne du singulier, car son sujet est toute la proposition introduite par \esp{que}.

Correct : \esp{Me parece que las noticias son falsas.} « Il me semble que les nouvelles sont fausses. »

Incorrect : \esp{Me parecen que las noticias son falsas.}

Astuce : traduis mentalement par « il me semble que » : le « il » français t'empêche d'accorder. Même logique que \esp{me gusta el fútbol} / \esp{me gustan los deportes} : ici, le sujet de \esp{parecer} est la phrase entière, donc singulier.
\end{attention}

\begin{attention}[Faux amis du lexique de l'information]
\begin{itemize}
\item \esp{la noticia} = la nouvelle, l'information (une seule) ; \esp{las noticias} = les informations, le journal télévisé. Ne dis pas \esp{una noticia} pour « des nouvelles » en général.
\item \esp{el bulo} = la fausse information, le canular (mot courant en Espagne) ; on dit aussi \esp{noticia falsa} ou, par anglicisme, \esp{fake news}.
\item \esp{la fuente} = la source (d'une information), mais aussi « la fontaine » ; ne confonds pas avec \esp{la salsa} (la sauce !).
\item \esp{actualmente} = actuellement, en ce moment (et non « actuellement » au sens de « en fait ») ; \esp{de hecho} = en fait, en réalité.
\end{itemize}
\end{attention}

\begin{exoresolu}[Hecho, opinión et mode du verbe]
Énoncé : (a) Classe chaque phrase : \esp{hecho} ou \esp{opinión}. (b) Traduis la phrase 4 en français.

1. \esp{El nuevo mercado abrió en marzo.} 2. \esp{El nuevo mercado es muy bonito.} 3. \esp{Creo que la prensa es libre.} 4. \esp{No creo que este bulo sea nuevo.}
\tcblower
Solution détaillée :

1. \textbf{Hecho} : date précise (\esp{en marzo}) et verbe objectif (\esp{abrió}) ; on peut vérifier l'ouverture du marché.

2. \textbf{Opinión} : l'adjectif évaluatif \esp{bonito} exprime un goût personnel.

3. \textbf{Opinión} : la marque \esp{creo que} signale un jugement. Le verbe \esp{es} est à l'indicatif car \esp{creo que} est affirmatif.

4. \textbf{Opinión} (négative) : \esp{no creo que} exige le subjonctif, d'où \esp{sea} (et non \esp{es}). Traduction : « Je ne crois pas que cette fausse information soit nouvelle. »
\end{exoresolu}

\courssub{Texte d'application : Un bulo en el grupo de la clase}

\begin{textebox}[Un bulo en el grupo de la clase]
El martes por la noche, un mensaje apareció en el grupo de WhatsApp de la clase de Première : « ¡Urgente! ¡No habrá clases mañana en todo el país! ¡Comparte este mensaje con todos tus contactos! ».

Muchos alumnos se alegraron y compartieron el mensaje enseguida. Pero Amina, una alumna seria, decidió pensar antes de actuar. Primero, miró la fuente : el mensaje no venía del Ministerio ni del director del liceo, sino de un número desconocido. Después, buscó la noticia en otros medios : ningún periódico, ninguna radio y ninguna cadena de televisión hablaba de eso. Luego, miró la fecha : el mensaje era una foto vieja de otro año. Finalmente, Amina escribió en el grupo : « No creo que este mensaje sea verdad. En mi opinión, es un bulo. Esperemos la confirmación oficial ».

A la mañana siguiente, todos los alumnos fueron al liceo : las clases eran normales. El director felicitó a Amina delante de toda la clase : « Verificar antes de compartir es la primera lección del ciudadano moderno ».

\emph{Glossaire :} \esp{por la noche} : le soir ; \esp{enseguida} : aussitôt ; \esp{desconocido} : inconnu ; \esp{viejo} : vieux ; \esp{a la mañana siguiente} : le lendemain matin ; \esp{felicitar} : féliciter ; \esp{verificar} : vérifier.
\end{textebox}

Questions de compréhension :

\begin{enumerate}
\item ¿Qué anuncia el mensaje de WhatsApp ? « Qu'annonce le message de WhatsApp ? »
\item ¿Cuáles son las cuatro verificaciones de Amina ? « Quelles sont les quatre vérifications d'Amina ? »
\item ¿Por qué el mensaje es un bulo ? Cita dos pruebas del texto. « Pourquoi le message est-il un \esp{bulo} ? Cite deux preuves du texte. »
\item ¿Qué frase usa Amina para dar su opinión ? « Quelle phrase Amina utilise-t-elle pour donner son opinion ? »
\end{enumerate}

\begin{corrige}[Questions de compréhension]
1. El mensaje anuncia que no habrá clases al día siguiente en todo el país y pide compartir el mensaje. « Le message annonce qu'il n'y aura pas cours le lendemain dans tout le pays et demande de le partager. »

2. Amina verifica cuatro cosas : la fuente (un número desconocido), la confirmación en otros medios (ninguno habla de eso), la fecha (una foto vieja de otro año) y espera la confirmación oficial.

3. El mensaje es un bulo porque no viene de una fuente oficial (Ministerio o director) y porque ningún periódico, radio o televisión lo confirma ; además, la foto es vieja.

4. Amina usa : \esp{No creo que este mensaje sea verdad. En mi opinión, es un bulo.} (verbe d'opinion négatif + subjonctif \esp{sea}, puis \esp{en mi opinión}).
\end{corrige}

\courssub{La méthode anti-bulo en quatre étapes}

Amina a appliqué, sans le savoir, la méthode professionnelle des journalistes. La voici, à apprendre par cœur.

\begin{methode}[¿Es un bulo? Cuatro preguntas antes de compartir]
\begin{enumerate}
\item \textbf{¿Quién publica?} Vérifie la \esp{fuente} : est-ce un média connu, un ministère, un compte officiel ? Un numéro inconnu ou un compte anonyme est suspect.
\item \textbf{¿Otros medios lo confirman?} Cherche la même \esp{noticia} dans deux ou trois médias sérieux. Si personne d'autre n'en parle, méfie-toi.
\item \textbf{¿La fecha es reciente?} Regarde la date : beaucoup de \esp{bulos} sont de vieilles nouvelles recyclées.
\item \textbf{¿La imagen es auténtica?} Une photo peut être retouchée ou prise dans un autre pays. Cherche l'image sur Internet pour trouver son origine.
\end{enumerate}
Conclusion : si une seule réponse est négative, \textbf{no compartas} (ne partage pas).
\end{methode}

\begin{popfigure}
\begin{tikzpicture}[node distance=0.55cm, font=\small]
\node[draw=popdark, fill=popcream, rounded corners, align=center, minimum width=4.5cm] (q0) {¿Es un bulo?};
\node[draw=popblue, fill=popblueL, rounded corners, align=center, minimum width=4.5cm, below=of q0] (q1) {¿Fuente fiable?};
\node[draw=popblue, fill=popblueL, rounded corners, align=center, minimum width=4.5cm, below=of q1] (q2) {¿Otros medios\\lo confirman?};
\node[draw=popgreen, fill=popgreenL, rounded corners, align=center, minimum width=3.6cm, below left=0.6cm and -1.2cm of q2] (ok) {COMPARTIR\\(con prudencia)};
\node[draw=poppink, fill=poppinkL, rounded corners, align=center, minimum width=3.6cm, below right=0.6cm and -1.2cm of q2] (no) {NO COMPARTIR\\es un bulo};
\draw[-{Stealth}, thick, popdark] (q0) -- (q1);
\draw[-{Stealth}, thick, popdark] (q1) -- node[right, font=\footnotesize]{sí} (q2);
\draw[-{Stealth}, thick, popgreen] (q2) -- node[left, font=\footnotesize]{sí} (ok);
\draw[-{Stealth}, thick, poppink] (q2) -- node[right, font=\footnotesize]{no} (no);
\draw[-{Stealth}, thick, poppink] (q1.east) -- ++(2.6,0) node[midway, above, font=\footnotesize]{no} |- (no.east);
\end{tikzpicture}
\legende{Organigramme de décision : deux questions suffisent souvent à bloquer un \esp{bulo}.}
\end{popfigure}

\courssub{Atelier : hecho u opinión}

\begin{exercice}[Trier les titres]
Classe ces huit phrases, extraites de titres de presse, dans deux colonnes : \esp{hecho} ou \esp{opinión}. Justifie chaque réponse avec un indice linguistique.

1. \esp{El presidente visitó ayer tres escuelas de Douala.} 2. \esp{Fue una visita histórica y maravillosa.} 3. \esp{El precio del cacao subió un diez por ciento.} 4. \esp{Creo que el cacao es el mejor producto del país.} 5. \esp{La selección ganó el partido por 3-1.} 6. \esp{Me parece que el árbitro fue terrible.} 7. \esp{El hospital abrirá sus puertas en enero.} 8. \esp{En mi opinión, es el hospital más moderno de África.}
\end{exercice}

\begin{corrige}[Exercice : trier les titres]
\textbf{Hechos} : 1 (verbe objectif \esp{visitó} + date \esp{ayer} + lieu) ; 3 (chiffre \esp{un diez por ciento}) ; 5 (score \esp{3-1} vérifiable) ; 7 (date \esp{en enero} + verbe factuel).

\textbf{Opiniones} : 2 (adjectifs évaluatifs \esp{histórica}, \esp{maravillosa}) ; 4 (\esp{creo que} + superlatif de goût) ; 6 (\esp{me parece que} + adjectif \esp{terrible}) ; 8 (\esp{en mi opinión} + superlatif).
\end{corrige}

\begin{experience}[Étude de cas : le message viral]
Ton camarade reçoit ce message WhatsApp : « ¡Increíble! Un estudio dice que beber agua del grifo de Yaoundé cura todas las enfermedades. ¡Comparte antes de que lo borren! »

En groupes de quatre : (1) appliquez les quatre questions de la méthode anti-bulo ; (2) décidez si vous partagez le message ; (3) rédigez une réponse en espagnol au groupe, avec \esp{creo que} ou \esp{no creo que}. Modèle possible : \esp{No creo que este mensaje sea verdad : ningún medio serio lo confirma. En mi opinión, no debemos compartirlo.}
\end{experience}

\begin{savaistu}[¿Sabías que...?]
Le mot espagnol \esp{bulo} vient de l'argot et désigne une rumeur fabriquée pour tromper. En Espagne, des journalistes spécialisés, les \esp{verificadores} ou \esp{fact-checkers}, passent leur journée à démentir les \esp{bulos} qui circulent sur les réseaux sociaux. En Guinée équatoriale, pays hispanophone voisin du Cameroun, la radio reste le média le plus écouté : là aussi, savoir distinguer \esp{hecho} et \esp{opinión} est une compétence de citoyen.
\end{savaistu}

\begin{aretenir}[L'essentiel de la section]
\begin{itemize}
\item Un \esp{hecho} se vérifie (chiffres, dates, lieux, verbes objectifs, \esp{según}) ; une \esp{opinión} se discute (\esp{creo que}, adjectifs évaluatifs).
\item \esp{Creo que} affirmatif $\rightarrow$ indicatif ; \esp{no creo que} $\rightarrow$ subjonctif (exigence du Bac) : \esp{No creo que sea verdad.}
\item \esp{Me parece que} reste au singulier, quel que soit le sujet de la subordonnée.
\item Avant de partager : fuente, confirmación, fecha, imagen. En cas de doute : \esp{no compartas}.
\end{itemize}
\end{aretenir}

\coursec{Méthode du commentaire de texte et commentaire modèle}

Dans la section précédente, nous avons appris à distinguer un \cle{hecho} (un fait, vérifiable) d'une \cle{opinión} (un jugement personnel). Cette compétence d'analyse est exactement celle que l'on te demandera au baccalauréat dans l'épreuve de \cle{comentario de texto} : lire un texte, le comprendre, puis en rendre compte de manière organisée et personnelle. Nous allons maintenant apprendre à construire ce commentaire, étape par étape, puis à lire ensemble un modèle complet rédigé sur notre texte support \esp{¿Verdad o bulo?}.

\courssub{Qu'est-ce qu'un commentaire de texte ?}

\begin{definition}[El comentario de texto]
Le \cle{comentario de texto} est un exercice écrit dans lequel l'élève \textbf{présente} un texte, en \textbf{analyse} le contenu et la langue, puis donne son \textbf{opinion personnelle} justifiée. Ce n'est ni une traduction, ni un résumé mot à mot : c'est un regard critique et organisé sur le texte.
\end{definition}

Observe d'abord ces trois questions, qui correspondent aux trois parties obligatoires du commentaire :

\begin{itemize}
\item \esp{¿Qué es el texto?} — Qu'est-ce que ce texte ? (type, thème, auteur)
\item \esp{¿Qué dice y cómo lo dice?} — Que dit-il et comment le dit-il ? (idées, lexique, faits et opinions)
\item \esp{¿Qué pienso yo?} — Qu'en pense-je, moi ? (opinion personnelle)
\end{itemize}

\begin{propriete}[La structure en trois parties]
Tout commentaire de texte au lycée suit le plan suivant :
\begin{enumerate}
\item \textbf{Introducción} : on présente le texte (type de document, thème principal, source ou auteur si connus).
\item \textbf{Desarrollo} : on analyse le texte (idées principales, lexique important, distinction hechos/opiniones), en citant de courtes phrases du texte.
\item \textbf{Conclusión} : on fait le bilan et on donne son opinion personnelle avec \esp{creo que}, \esp{en mi opinión}, \esp{me parece que}.
\end{enumerate}
\end{propriete}

\begin{popfigure}
\begin{center}
\begin{tikzpicture}[node distance=0.9cm]
\node[draw=popblue, fill=popblueL, rounded corners, minimum width=9cm, minimum height=1.4cm, align=center] (intro) {\textbf{1. INTRODUCCIÓN}\\ \esp{¿Qué es el texto?} (tipo, tema, autor)};
\node[draw=poporange, fill=poporangeL, rounded corners, minimum width=9cm, minimum height=1.4cm, align=center, below=of intro] (des) {\textbf{2. DESARROLLO}\\ \esp{Análisis:} ideas, léxico, \esp{hechos} y \esp{opiniones}};
\node[draw=popgreen, fill=popgreenL, rounded corners, minimum width=9cm, minimum height=1.4cm, align=center, below=of des] (conc) {\textbf{3. CONCLUSIÓN}\\ \esp{Mi opinión personal} + bilan};
\draw[-{Stealth}, very thick, popdark] (intro) -- (des);
\draw[-{Stealth}, very thick, popdark] (des) -- (conc);
\end{tikzpicture}
\end{center}
\legende{Les trois blocs du commentaire de texte : présenter, analyser, conclure.}
\end{popfigure}

\courssub{Les phrases-outils du commentaire}

Pour rédiger un bon commentaire, il faut connaître par cœur des \cle{frases útiles} (phrases-outils). Les voici, classées par partie, avec leur traduction française.

\begin{center}
\begin{tabularx}{\linewidth}{|X|X|X|}
\hline
\rowcolor{popblueL} \textbf{Partie} & \textbf{Español} & \textbf{Français} \\
\hline
Introducción & \esp{Este texto es un artículo de prensa que trata de...} & Ce texte est un article de presse qui traite de... \\
\hline
Introducción & \esp{El tema principal es la información responsable.} & Le thème principal est l'information responsable. \\
\hline
Introducción & \esp{El texto fue escrito por un periodista.} & Le texte a été écrit par un journaliste. \\
\hline
Desarrollo & \esp{El autor explica que...} & L'auteur explique que... \\
\hline
Desarrollo & \esp{En el segundo párrafo, el texto presenta...} & Dans le deuxième paragraphe, le texte présente... \\
\hline
Desarrollo & \esp{Podemos distinguir hechos y opiniones.} & Nous pouvons distinguer faits et opinions. \\
\hline
Desarrollo & \esp{Por ejemplo, el texto dice: «...».} & Par exemple, le texte dit : «...». \\
\hline
Conclusión & \esp{En conclusión, este texto nos enseña que...} & En conclusion, ce texte nous enseigne que... \\
\hline
Conclusión & \esp{En mi opinión, este texto es interesante porque...} & À mon avis, ce texte est intéressant parce que... \\
\hline
Conclusión & \esp{Me parece que es importante verificar la fuente.} & Il me semble qu'il est important de vérifier la source. \\
\hline
\end{tabularx}
\end{center}

\begin{exemplebox}[« Tratar de » : traiter de]
Le verbe \esp{tratar de} suivi d'un nom signifie « traiter de, parler de ». C'est le verbe roi de l'introduction du commentaire.

\esp{El texto trata de los medios de comunicación.} « Le texte traite des médias. »

\esp{El artículo trata de un bulo en las redes sociales.} « L'article traite d'une fausse information sur les réseaux sociaux. »

\esp{¿De qué trata el texto?} « De quoi traite le texte ? »

Attention à la préposition : on dit toujours \esp{tratar \textbf{de} algo}, jamais \esp{tratar algo} dans ce sens.
\end{exemplebox}

\begin{attention}[« Sin embargo » : le faux ami à ne jamais oublier]
\esp{Sin embargo} signifie « \textbf{cependant} » et non « sans embargo » ! C'est un des faux amis les plus fréquents pour les francophones.

\esp{La noticia era falsa; sin embargo, Juan la compartió.} « La nouvelle était fausse ; cependant, Juan l'a partagée. »

Autres connecteurs utiles à ne pas confondre : \esp{además} = « de plus » ; \esp{por eso} = « c'est pourquoi » ; \esp{en cambio} = « en revanche ».
\end{attention}

\begin{attention}[Le genre des mots des médias]
Les francophones se trompent souvent sur le genre des noms de ce thème. Retiens :

\esp{la noticia} (féminin) : \esp{una noticia falsa}, \esp{las noticias} — « une nouvelle fausse », « les informations ».

\esp{el medio} (masculin) : \esp{los medios de comunicación} — « les médias ».

\esp{la fuente} (féminin) : \esp{una fuente fiable} — « une source fiable ».

\esp{el bulo} (masculin) : \esp{un bulo} — « une fausse information, une intox ».

\esp{el titular} (masculin) : \esp{el titular del artículo} — « le titre (la manchette) de l'article ».
\end{attention}

\courssub{Méthode complète pour réussir le commentaire au Bac}

\begin{methode}[Réussir le commentaire de texte en 5 étapes]
\begin{enumerate}
\item \textbf{Lire le texte deux fois} : une première fois pour le sens général, une deuxième fois en soulignant les mots-clés (\esp{prensa, noticia, bulo, fuente...}).
\item \textbf{Présenter le texte} (2-3 lignes) : type de document + thème principal, avec \esp{Este texto es... que trata de...}.
\item \textbf{Analyser avec des citations} (5-7 lignes) : dégager 2 ou 3 idées principales, citer de courtes phrases du texte entre guillemets, distinguer \esp{hechos} et \esp{opiniones}, relever le lexique du thème.
\item \textbf{Conclure avec son opinion personnelle} (2-3 lignes) : bilan + \esp{En mi opinión...} + justification avec \esp{porque}.
\item \textbf{Relire} : vérifier l'accord des adjectifs, les accents (\esp{opinión, artículo, párrafo}) et la présence des trois parties.
\end{enumerate}
\end{methode}

\begin{aretenir}[Les trois interdits du commentaire]
\begin{itemize}
\item Ne \textbf{traduis pas} le texte : le commentaire analyse, il ne traduit pas.
\item N'\textbf{oublie jamais} l'opinion personnelle à la fin : sans \esp{en mi opinión}, le commentaire est incomplet.
\item Ne reste pas \textbf{vague} : cite toujours le texte (\esp{el texto dice: «...»}) pour prouver ce que tu affirmes.
\end{itemize}
\end{aretenir}

\courssub{Commentaire modèle sur « ¿Verdad o bulo? »}

Lisons maintenant ensemble un commentaire complet, rédigé selon la méthode, sur notre texte support. Lis-le d'abord en silence, puis nous l'analyserons.

\begin{textebox}[Comentario modelo — « ¿Verdad o bulo? »]
Este texto es un artículo de prensa sobre la información responsable. El tema principal es el problema de los bulos en las redes sociales. El autor cuenta la historia de Juan, un joven que compartió un bulo sin comprobar la fuente.

Primero, el texto presenta los hechos: Juan leyó una noticia falsa y la compartió con sus amigos. El artículo dice: «Juan no verificó la fuente antes de compartir». Además, el texto distingue claramente el hecho de la opinión: el hecho es que la noticia era falsa; la opinión es que los jóvenes deben ser más responsables. Sin embargo, la idea más importante es la lección final: verificar siempre la fuente antes de creer una noticia.

En conclusión, este texto nos enseña a leer la información con espíritu crítico. En mi opinión, este texto es muy útil porque nos ayuda a ser responsables con la información que compartimos. Me parece que todos los alumnos deberían leer artículos como este.
\end{textebox}

\textbf{Glossaire :} \esp{compartió} = a partagé (prétérit de \esp{compartir}) ; \esp{sin comprobar} = sans vérifier ; \esp{la fuente} = la source ; \esp{deberían} = devraient (conditionnel de \esp{deber}) ; \esp{espíritu crítico} = esprit critique ; \esp{enseña a} = apprend à.

\textbf{Questions de lecture :}
\begin{enumerate}
\item \esp{¿Qué tipo de texto es?} (Quel type de texte est-ce ?)
\item \esp{¿Cuál es el tema principal?} (Quel est le thème principal ?)
\item \esp{¿Qué hecho cita el comentario?} (Quel fait le commentaire cite-t-il ?)
\item \esp{¿Cuál es la opinión personal del autor del comentario?} (Quelle est l'opinion personnelle de l'auteur du commentaire ?)
\end{enumerate}

\courssub{Analyse du modèle : les connecteurs et la structure}

Repérons maintenant dans le modèle les \cle{conectores} qui donnent du rythme au texte et marquent la progression logique :

\begin{center}
\begin{tabularx}{\linewidth}{|X|X|X|}
\hline
\rowcolor{popblueL} \textbf{Conector} & \textbf{Français} & \textbf{Fonction dans le modèle} \\
\hline
\esp{Primero} & D'abord, premièrement & Annonce la première idée del desarrollo \\
\hline
\esp{Además} & De plus, en outre & Ajoute une idée (la distinction hecho/opinión) \\
\hline
\esp{Sin embargo} & Cependant & Introduit une nuance : l'idée la plus importante \\
\hline
\esp{En conclusión} & En conclusion & Ouvre la dernière partie \\
\hline
\esp{En mi opinión} & À mon avis & Introduit l'opinion personnelle \\
\hline
\esp{Porque} & Parce que & Justifie l'opinion personnelle \\
\hline
\end{tabularx}
\end{center}

Observe aussi la répartition : le premier paragraphe est l'\textbf{introducción} (type + thème), le deuxième est le \textbf{desarrollo} (faits, citation, distinction hecho/opinión), le troisième est la \textbf{conclusión} (bilan + opinion personnelle). Chaque partie remplit exactement sa fonction : c'est cela, un commentaire réussi.

\begin{propriete}[Opinion personnelle : l'indicatif après « creo que »]
Après \esp{creo que}, \esp{me parece que} et \esp{en mi opinión} (formes \textbf{affirmatives}), on emploie l'\textbf{indicatif}, car on affirme ce que l'on pense :

\esp{Creo que este texto \textbf{es} útil.} « Je pense que ce texte est utile. »

\esp{Me parece que los periodistas \textbf{tienen} una gran responsabilidad.} « Il me semble que les journalistes ont une grande responsabilité. »

En revanche, à la forme \textbf{négative}, on emploie le \textbf{subjonctif} : \esp{No creo que esta noticia \textbf{sea} verdadera.} « Je ne crois pas que cette nouvelle soit vraie. » À ton niveau, retiens surtout la forme affirmative avec l'indicatif.
\end{propriete}

\begin{exoresolu}[Identifier les parties d'un commentaire]
Énoncé. Voici trois phrases extraites d'un commentaire. Indique à quelle partie (introducción, desarrollo, conclusión) appartient chacune, puis traduis-les en français.

a) \esp{En mi opinión, los periodistas tienen una gran responsabilidad.}

b) \esp{Este texto es un artículo que trata de las redes sociales.}

c) \esp{En el segundo párrafo, el autor explica que muchos jóvenes no verifican las noticias.}
\tcblower
Solution détaillée.

a) \textbf{Conclusión} : la présence de \esp{en mi opinión} (à mon avis) marque l'opinion personnelle, qui appartient toujours à la conclusion. Traduction : « À mon avis, les journalistes ont une grande responsabilité. »

b) \textbf{Introducción} : la formule \esp{Este texto es un artículo que trata de...} sert à présenter le type et le thème du texte. Traduction : « Ce texte est un article qui traite des réseaux sociaux. »

c) \textbf{Desarrollo} : \esp{el autor explica que...} et la référence au \esp{segundo párrafo} sont des marques typiques de l'analyse. Traduction : « Dans le deuxième paragraphe, l'auteur explique que beaucoup de jeunes ne vérifient pas les nouvelles. »
\end{exoresolu}

\begin{exoresolu}[Corriger les erreurs d'un élève]
Énoncé. Voici la conclusion rédigée par un élève de Première. Elle contient trois erreurs. Repère-les et corrige-les.

\esp{En conclusión, este texto es interesante. Sin embargo, me parece que las noticias son siempre falsas. El texto dice que los jóvenes comparten bulos.}
\tcblower
Solution détaillée.

Erreur 1 : \esp{sin embargo} est mal employé : il n'y a pas d'opposition entre « le texte est intéressant » et la phrase suivante. Il faut un connecteur d'addition ou d'opinion : \esp{Además} ou directement \esp{En mi opinión}.

Erreur 2 : la généralisation \esp{las noticias son siempre falsas} (« les nouvelles sont toujours fausses ») est abusive et contredit l'esprit critique : ce sont les \esp{bulos} qui sont faux, pas toutes les nouvelles. Mieux : \esp{me parece que debemos verificar las noticias antes de compartirlas}.

Erreur 3 : la dernière phrase est une simple répétition du texte, pas une conclusion personnelle : il manque la justification avec \esp{porque}.

Version corrigée possible : \esp{En conclusión, este texto es interesante. En mi opinión, debemos verificar siempre la fuente, porque compartir un bulo puede ser peligroso.} « En conclusion, ce texte est intéressant. À mon avis, nous devons toujours vérifier la source, parce que partager une fausse information peut être dangereux. »
\end{exoresolu}

\begin{attention}[Les trois erreurs qui coûtent des points]
\begin{itemize}
\item \textbf{Traduire au lieu d'analyser} : écrire \esp{El texto dice que Juan leyó una noticia y la compartió y luego...} en racontant tout le texte n'est pas un commentaire. Sélectionne 2 ou 3 idées et commente-les.
\item \textbf{Oublier l'opinion personnelle} : la conclusion sans \esp{creo que} ou \esp{me parece que} est une conclusion ratée. Donne ton avis et justifie-le avec \esp{porque}.
\item \textbf{Ne pas citer le texte} : une affirmation sans preuve ne vaut rien. Utilise \esp{el texto dice: «...»} ou \esp{por ejemplo, en el primer párrafo...}.
\end{itemize}
\end{attention}

\begin{exercice}[À ton tour]
Rédige un mini-commentaire (8 à 10 lignes) sur un article de presse de ton choix (par exemple un article sur le football au Cameroun ou sur la vie de ton lycée). Respecte le plan : introducción con \esp{tratar de}, desarrollo con al menos una cita y los conectores \esp{primero}, \esp{además}, \esp{sin embargo}, y conclusión con \esp{en mi opinión... porque...}.
\end{exercice}

\begin{corrige}[Exercice — éléments de correction]
Un bon commentaire doit contenir :
\begin{itemize}
\item une introduction du type : \esp{Este texto es un artículo de prensa que trata del fútbol en Camerún. El tema principal es...} ;
\item un développement avec une citation courte entre guillemets, la distinction entre un fait (\esp{el hecho es que...}) et une opinion (\esp{la opinión del autor es que...}), et les connecteurs \esp{primero}, \esp{además}, \esp{sin embargo} ;
\item une conclusion : \esp{En conclusión, en mi opinión, este texto es interesante porque...}.
Vérifier aussi : accents (\esp{artículo, opinión, párrafo}), genre des noms (\esp{la noticia}, \esp{el medio}), \esp{sin embargo} = « cependant », indicatif après \esp{creo que} affirmatif, et aucune traduction mot à mot du texte.
\end{itemize}
\end{corrige}

\begin{aretenir}[L'essentiel de la section]
Le commentaire de texte se construit en trois blocs : \textbf{introducción} (\esp{Este texto es... que trata de...}), \textbf{desarrollo} (analyse avec citations, \esp{hechos} et \esp{opiniones}, connecteurs \esp{primero, además, sin embargo}) et \textbf{conclusión} (\esp{En conclusión... En mi opinión... porque...}). Après \esp{creo que} affirmatif, on emploie l'indicatif. Trois pièges à éviter : traduire le texte, oublier son opinion personnelle, et ne pas citer le texte. Avec ces outils, tu es prêt pour la production guidée de la prochaine section : \esp{escribir y opinar}.
\end{aretenir}

\coursec{Production guidée : escribir y opinar}

Dans la section précédente, nous avons appris à commenter un texte de presse : présenter le document, dégager le thème, distinguer les faits des opinions et formuler un jugement personnel. Nous allons maintenant passer de l'analyse à la \cle{production} : c'est à votre tour d'écrire et de parler. Objectif final de la leçon : savoir raconter une expérience liée à l'information, donner son avis de manière argumentée et rédiger un message responsable qui combat un \esp{bulo}. Trois tâches progressives vous y conduisent : un paragraphe écrit, un débat oral et un message pour WhatsApp.

\courssub{La démarche de production : du plan à la correction}

Avant toute rédaction, le bon hispanophone suit un chemin précis. Observez la frise ci-dessous : elle résume les cinq étapes que vous devez suivre pour chaque tâche écrite de cette leçon.

\begin{popfigure}
\begin{tikzpicture}[node distance=0.35cm]
\node[draw=popblue, fill=popblueL, rounded corners, minimum width=2.5cm, minimum height=1cm, align=center] (a) {\textbf{1. Plan}\\ \esp{ideas}};
\node[draw=popgreen, fill=popgreenL, rounded corners, minimum width=2.5cm, minimum height=1cm, align=center, right=of a] (b) {\textbf{2. Brouillon}\\ \esp{borrador}};
\node[draw=poporange, fill=poporangeL, rounded corners, minimum width=2.5cm, minimum height=1cm, align=center, right=of b] (c) {\textbf{3. Rédaction}\\ \esp{redacción}};
\node[draw=poppurple, fill=poppurpleL, rounded corners, minimum width=2.5cm, minimum height=1cm, align=center, right=of c] (d) {\textbf{4. Relecture}\\ \esp{revisión}};
\node[draw=poppink, fill=poppinkL, rounded corners, minimum width=2.5cm, minimum height=1cm, align=center, right=of d] (e) {\textbf{5. Correction}\\ \esp{versión final}};
\draw[-{Stealth}, thick, popdark] (a) -- (b);
\draw[-{Stealth}, thick, popdark] (b) -- (c);
\draw[-{Stealth}, thick, popdark] (c) -- (d);
\draw[-{Stealth}, thick, popdark] (d) -- (e);
\draw[-{Stealth}, thick, popmuted, dashed] (d.south) to[bend right=35] (c.south);
\end{tikzpicture}
\legende{Les cinq étapes de la production écrite ; la flèche en pointillés rappelle que la relecture ramène souvent à la rédaction.}
\end{popfigure}

\begin{methode}[Écrire un paragraphe d'opinion en quatre étapes]
\begin{enumerate}
\item \textbf{Présenter le fait} : racontez l'événement au \esp{pretérito indefinido} (action ponctuelle et terminée). Exemple : \esp{Ayer compartí una noticia en WhatsApp.} « Hier, j'ai partagé une nouvelle sur WhatsApp. »
\item \textbf{Donner son opinion} : utilisez \esp{creo que}, \esp{en mi opinión} ou \esp{me parece que}, suivis de l'indicatif. Exemple : \esp{En mi opinión, fue un error.} « À mon avis, ce fut une erreur. »
\item \textbf{Justifier} : introduisez la raison avec \esp{porque}. Exemple : \esp{...porque no verifiqué la fuente.} « ...parce que je n'ai pas vérifié la source. »
\item \textbf{Conclure} : terminez par une leçon ou un conseil, souvent avec \esp{ahora} ou \esp{hay que}. Exemple : \esp{Ahora sé que hay que comprobar siempre la información.} « Maintenant, je sais qu'il faut toujours vérifier l'information. »
\end{enumerate}
\end{methode}

\courssub{Tâche 1 (écrite) : « Una noticia que compartí »}

\begin{experience}[Consigne de la tâche 1]
Rédigez un court paragraphe de 5 à 6 lignes intitulé \esp{« Una noticia que compartí »}. Contraintes obligatoires :
\begin{itemize}
\item employer au moins \textbf{5 mots du lexique} de la leçon : \esp{prensa, noticia, periodista, titular, redes sociales, bulo, fuente, hecho, opinión} ;
\item employer au moins \textbf{2 expressions d'opinion} : \esp{creo que}, \esp{en mi opinión}, \esp{me parece que} ;
\item raconter l'expérience au \esp{pretérito indefinido} et donner votre avis au présent.
\end{itemize}
\end{experience}

Avant d'écrire, lisez ce modèle rédigé par un élève de Première de Yaoundé, puis observez comment les contraintes sont respectées.

\begin{exemplebox}[Modèle de paragraphe]
\esp{El mes pasado compartí en las redes sociales una noticia con un titular muy llamativo: decía que el gobierno cerraba todos los mercados de Douala. No verifiqué la fuente y era un bulo. Un periodista de la prensa local explicó el hecho verdadero en la radio. En mi opinión, compartir sin comprobar es peligroso. Creo que hay que distinguir siempre el hecho de la opinión, y me parece que debemos leer la prensa seria antes de creer cualquier noticia.}\par
\emph{Traduction : « Le mois dernier, j'ai partagé sur les réseaux sociaux une nouvelle au titre très accrocheur : elle disait que le gouvernement fermait tous les marchés de Douala. Je n'ai pas vérifié la source et c'était une fausse information. Un journaliste de la presse locale a expliqué le fait véridique à la radio. À mon avis, partager sans vérifier est dangereux. Je crois qu'il faut toujours distinguer le fait de l'opinion, et il me semble que nous devons lire la presse sérieuse avant de croire n'importe quelle nouvelle. »}
\end{exemplebox}

Vérifions ensemble le respect des contraintes : les 9 mots du lexique sont présents (\esp{redes sociales, noticia, titular, fuente, bulo, periodista, prensa, hecho, opinión}), ainsi que 3 expressions d'opinion (\esp{en mi opinión, creo que, me parece que}). Le récit des actions ponctuelles est au \esp{pretérito indefinido} : \esp{compartí, verifiqué, explicó}. Remarquez que \esp{decía}, \esp{cerraba} et \esp{era} sont au \esp{pretérito imperfecto} : ce ne sont pas des actions accomplies, mais la \textbf{description} du contenu de la nouvelle et de sa nature — l'imparfait peint le décor, l'indefinido fait avancer le récit.

\begin{attention}[Compartí ou compartía ?]
Au passé, \esp{compartí} est le \esp{pretérito indefinido} : il exprime une action \textbf{ponctuelle et terminée} (« j'ai partagé une fois »). Ne le confondez pas avec \esp{compartía} (\esp{pretérito imperfecto}), qui exprime une \textbf{habitude} ou une action en cours : \esp{Antes compartía noticias todos los días.} « Avant, je partageais des nouvelles tous les jours. » Erreur fréquente du francophone : écrire \esp{*compartía una noticia ayer*} pour raconter un événement unique. Dites : \esp{Ayer compartí una noticia.} Autre piège : l'accent de \esp{compartí} est obligatoire ; \esp{*comparti*} sans accent est une faute. Même vigilance pour \esp{verifiqué} : accent final obligatoire, et \esp{qu} devant \esp{e} pour conserver le son « k ».
\end{attention}

\begin{exoresolu}[Transformer des notes en paragraphe]
Énoncé : à partir de ces notes — \esp{noticia falsa / elecciones / no verificar la fuente / un amigo periodista / aprender la lección} — rédigez 4 phrases avec deux expressions d'opinion.
\tcblower
Solution détaillée : 1) Fait au passé : \esp{La semana pasada compartí una noticia falsa sobre las elecciones.} 2) Opinion : \esp{En mi opinión, fue una irresponsabilidad, porque no verifiqué la fuente.} 3) Justification et fait supplémentaire : \esp{Un amigo periodista me explicó que era un bulo.} 4) Conclusion : \esp{Creo que aprendí la lección: ahora siempre compruebo la información antes de compartir.} Le paragraphe respecte la méthode en quatre étapes : fait, opinion, justification, conclusion. Notez \esp{sobre las elecciones} (« sur les élections » : \esp{sobre} + nom, jamais \esp{*de las elecciones*} dans ce sens) et \esp{era} à l'imparfait, car il décrit la nature de la nouvelle.
\end{exoresolu}

\courssub{Tâche 2 (orale) : débat en binômes}

\begin{experience}[¿Las redes sociales informan o desinforman?]
Par binômes, préparez un débat de 3 minutes. L'élève A défend la thèse \esp{« Las redes sociales informan »} ; l'élève B défend \esp{« Las redes sociales desinforman »}. Règles du jeu :
\begin{itemize}
\item chaque argument commence obligatoirement par \esp{creo que}, \esp{en mi opinión} ou \esp{me parece que} ;
\item chaque opinion doit être justifiée par \esp{porque} + un exemple concret (Cameroun, école, famille) ;
\item pour réagir à votre partenaire, utilisez : \esp{Estoy de acuerdo, pero...} / \esp{No estoy de acuerdo, porque...} / \esp{Tienes razón en parte, sin embargo...}
\end{itemize}
\end{experience}

Voici des amorces pour vous aider à démarrer.

\begin{center}
\begin{tabularx}{\linewidth}{|X|X|}
\hline
\rowcolor{popblueL} \textbf{Thèse : informan} & \textbf{Thèse : desinforman} \\
\hline
\esp{Creo que informan porque las noticias llegan muy rápido.} & \esp{Me parece que desinforman porque circulan muchos bulos.} \\
\hline
\esp{En mi opinión, los jóvenes conocen el mundo gracias a las redes.} & \esp{En mi opinión, nadie verifica la fuente antes de compartir.} \\
\hline
\esp{Creo que los periodistas también usan las redes para informar.} & \esp{Creo que un titular falso se comparte más que un hecho verdadero.} \\
\hline
\end{tabularx}
\end{center}

\begin{attention}[Après « creo que », l'indicatif]
En espagnol, \esp{creo que}, \esp{me parece que} et \esp{en mi opinión} sont suivis de l'\textbf{indicatif}, jamais du subjonctif : \esp{Creo que las redes \textbf{son} útiles.} « Je crois que les réseaux sont utiles. » En revanche, à la forme négative, le subjonctif apparaît : \esp{No creo que esa noticia \textbf{sea} verdadera.} « Je ne crois pas que cette nouvelle soit vraie. » C'est une différence majeure avec le français, où l'indicatif suit souvent « je crois que » sans nuance : retenez le couple \esp{creo que + indicatif} / \esp{no creo que + subjonctif}. Attention aussi à l'accord : \esp{estar de acuerdo} est invariable en genre pour l'accord du locuteur — on dit \esp{Estoy de acuerdo} que l'on soit un garçon ou une fille.
\end{attention}

\courssub{Tâche 3 : rédiger un « mensaje responsable »}

Dernière étape : la production la plus utile de la leçon. Vous recevez dans votre groupe WhatsApp un \esp{bulo} (par exemple : « les écoles ferment demain »). Votre mission : rédiger un message de 3 à 4 phrases qui \textbf{dément} la rumeur et \textbf{indique une source fiable}.

\begin{methode}[Structure du mensaje responsable]
\begin{enumerate}
\item \textbf{Alerter} : \esp{Atención: ...} ou \esp{Cuidado: ...}
\item \textbf{Démentir} : \esp{La noticia de... es un bulo.} / \esp{Es falso que...}
\item \textbf{Donner la source fiable} : \esp{Ningún medio serio lo confirma.} / \esp{La fuente oficial (ministerio, prensa seria) dice que...}
\item \textbf{Appeler à la responsabilité} : \esp{Por favor, comprueben la fuente antes de compartir.}
\end{enumerate}
\end{methode}

\begin{textebox}[Modèle de mensaje responsable]
\esp{Atención: la noticia del cierre de las escuelas es un bulo. Ningún medio serio lo confirma. El Ministerio de Educación no ha publicado ningún comunicado oficial. Por favor, comprueben la fuente antes de compartir. Informar bien es responsabilidad de todos.}
\end{textebox}

Glossaire : \esp{cierre} « fermeture » ; \esp{comunicado oficial} « communiqué officiel » ; \esp{comprueben} « vérifiez » (subjonctif présent de \esp{comprobar}, employé ici comme impératif poli pluriel, forme \esp{ustedes}) ; \esp{responsabilidad de todos} « la responsabilité de tous ». Notez l'emploi de \esp{ningún} (apocope de \esp{ninguno} devant un nom masculin singulier : \esp{ningún medio}, mais \esp{ninguna fuente} au féminin) et le pronom \esp{lo} qui reprend \esp{el cierre} (ou le contenu de la rumeur). Observez enfin le parfait composé \esp{no ha publicado} : il présente l'absence de communiqué comme un fait d'actualité récente, ce qui est le temps naturel de la presse.

\begin{exoresolu}[Corriger un mensaje mal rédigé]
Énoncé : ce message contient trois erreurs ; corrigez-le. \esp{« Atención: la noticia es un bolo. Yo compartía ayer esta información pero es falso. Compartan sin verificar. »}
\tcblower
Solution détaillée : 1) \esp{bolo} est un faux ami orthographique : le mot correct est \esp{bulo} « fausse information » (\esp{bolo} signifie « gâteau » au Mexique ou « quilles » dans d'autres pays !). 2) \esp{compartía} exprime une habitude ; pour une action ponctuelle datée (\esp{ayer}), il faut l'indefinido : \esp{compartí}. 3) Deux problèmes dans la fin : \esp{es falso} doit s'accorder avec \esp{la información} (féminin) : \esp{es falsa} ; et le conseil est contradictoire — on ne dit pas \esp{compartan sin verificar} mais \esp{no compartan sin verificar} ou \esp{comprueben la fuente antes de compartir}. Message corrigé : \esp{« Atención: la noticia es un bulo. Ayer compartí esta información, pero es falsa. Comprueben la fuente antes de compartir. »}
\end{exoresolu}

\courssub{Relecture croisée et grille d'auto-évaluation}

Une fois votre paragraphe et votre message rédigés, échangez votre cahier avec votre voisin. Chacun évalue la production de l'autre à l'aide de la grille ci-dessous, puis la classe corrige collectivement au tableau une production choisie par le professeur : on souligne les mots du lexique, on encadre les expressions d'opinion et on corrige ensemble les fautes de conjugaison.

\begin{center}
\begin{tabularx}{\linewidth}{|X|X|c|}
\hline
\rowcolor{popblueL} \textbf{Critère} & \textbf{Qué comprobar} & \textbf{Puntos} \\
\hline
Lexique des médias & au moins 5 mots : \esp{noticia, bulo, fuente, titular, prensa...} & 2 pts \\
\hline
Fait / opinion distingués & le récit (passé) est séparé du jugement (présent) & 2 pts \\
\hline
Expression de l'opinion & au moins 2 expressions : \esp{creo que, en mi opinión, me parece que} & 2 pts \\
\hline
Correction grammaticale & \esp{compartí} avec accent, accords, ponctuation ¿ ¡ & 2 pts \\
\hline
\end{tabularx}
\end{center}

\begin{aretenir}[Les réflexes du producteur responsable]
\begin{itemize}
\item Raconter : \esp{pretérito indefinido} pour l'action unique (\esp{compartí, verifiqué, fue}) ; \esp{pretérito imperfecto} pour la description (\esp{decía, era}).
\item Opiner : \esp{creo que / en mi opinión / me parece que} + indicatif ; \esp{no creo que} + subjonctif.
\item Justifier : \esp{porque} ; conclure : \esp{ahora, hay que}.
\item Démentir : \esp{es un bulo / ningún medio serio lo confirma / comprueben la fuente}.
\end{itemize}
\end{aretenir}

\courssub{Bilan de la leçon}

Pour clôturer la séquence, chaque élève formule à voix haute, en une phrase en espagnol, ce qu'il a appris aujourd'hui. Voici trois exemples possibles :

\begin{exemplebox}[Frases de balance]
\esp{Hoy aprendí a distinguir un hecho de una opinión.} « Aujourd'hui, j'ai appris à distinguer un fait d'une opinion. »\par
\esp{Ahora sé que no debo compartir una noticia sin verificar la fuente.} « Maintenant, je sais que je ne dois pas partager une nouvelle sans vérifier la source. »\par
\esp{En mi opinión, informar bien es una responsabilidad de todos.} « À mon avis, bien informer est la responsabilité de tous. »
\end{exemplebox}

\begin{exercice}[Pour aller plus loin]
À la maison, rédigez un second \esp{mensaje responsable} (4 phrases) qui dément ce \esp{bulo} : \esp{« Mañana no hay clases en todos los liceos de Yaoundé »}. Respectez la structure en quatre étapes (alerter, démentir, source fiable, appel à la responsabilité) et soulignez vos deux expressions d'opinion.
\end{exercice}

\begin{corrige}[Exercice « Pour aller plus loin »]
Production possible : \esp{« Atención: la noticia de que mañana no hay clases en los liceos de Yaoundé es un bulo. Ningún medio serio ni el Ministerio lo confirman. En mi opinión, debemos consultar siempre la página oficial del establecimiento. Creo que compartir sin verificar es peligroso: comprueben la fuente antes de compartir. »} Barème : 2 pts de lexique (\esp{noticia, bulo, medio, fuente}), 2 pts pour les deux expressions d'opinion, 2 pts pour la distinction fait/opinion, 2 pts pour la correction grammaticale. Remarque : après \esp{ningún medio serio ni el Ministerio}, le verbe au pluriel \esp{confirman} est correct car le sujet coordonne deux éléments.
\end{corrige}

\newpage

\coursec{Activité d'intégration}

\courssub{Objectifs de l'activité}

À la fin de cette activité d'intégration, l'élève sera capable de :

\begin{itemize}
\item mobiliser le \cle{lexique des médias} (\esp{prensa}, \esp{noticia}, \esp{periodista}, \esp{titular}, \esp{redes sociales}, \esp{bulo}, \esp{fuente}, \esp{hecho}, \esp{opinión}) dans une situation de communication réelle ;
\item appliquer la \cle{grille de vérification en 4 étapes} (fuente, confirmación, fecha, imagen) pour analyser un message reçu sur les réseaux sociaux ;
\item distinguer un \cle{hecho} (fait vérifiable) d'une \cle{opinión} (jugement personnel) dans un texte court ;
\item rédiger un \cle{mensaje de verificación} clair et poli qui dément une fausse information ;
\item exprimer son opinion avec les structures \esp{creo que}, \esp{en mi opinión}, \esp{me parece que} suivies de l'indicatif ;
\item présenter oralement un verdict argumenté en une minute.
\end{itemize}

\courssub{Situation}

Tu es membre du club d'espagnol de ton lycée à Yaoundé. Ce matin, dans le groupe WhatsApp de la classe de Première A, trois messages circulent et provoquent de vives discussions. Certains camarades les croient, d'autres hésitent, et le délégué demande : \esp{¿Es verdad o es un bulo?} Le professeur d'espagnol propose alors un jeu : chaque groupe de quatre élèves devient une équipe de \emph{caza-bulos} (chasseurs de fausses informations). Votre mission : analyser les trois messages, identifier le vrai et les deux faux, puis rédiger un message de vérification pour rétablir la vérité dans le groupe WhatsApp.

\begin{textebox}[Los tres mensajes del grupo WhatsApp]
\textbf{Mensaje 1} — reenviado 47 veces\par
\esp{¡¡URGENTE!! El Ministerio de Educación cancela todos los exámenes oficiales de este año en Camerún por falta de dinero. Comparte este mensaje con todos tus amigos antes de que lo borren. ¡No estudies más, ya no hay examen!}\par\medskip
\textbf{Mensaje 2} — publicado por la tía de un alumno\par
\esp{Mi vecino, que es médico, dice que beber agua con limón cura todas las enfermedades, incluso la malaria. Es un secreto que las farmacias no quieren que sepas. Yo lo hago cada mañana y nunca estoy enferma. ¡Es la verdad!}\par\medskip
\textbf{Mensaje 3} — publicado por el club de deportes del liceo\par
\esp{El equipo de fútbol del liceo juega el sábado a las 15 h contra el Lycée de Nkol-Eton en el estadio del barrio. La entrada es gratuita. El entrenador confirma el partido en la página oficial del liceo.}
\end{textebox}

\begin{definition}[Glossaire du document]
\esp{reenviado} : transféré ; \esp{cancela} : annule ; \esp{por falta de dinero} : faute d'argent ; \esp{antes de que lo borren} : avant qu'on ne l'efface ; \esp{cura} : guérit ; \esp{las farmacias} : les pharmacies ; \esp{la entrada es gratuita} : l'entrée est gratuite ; \esp{el entrenador} : l'entraîneur.
\end{definition}

\courssub{Consignes}

Travaillez par groupes de quatre. Répartissez les rôles : un \esp{secretario} (il écrit), un \esp{portavoz} (il présente le verdict), deux \esp{verificadores} (ils appliquent la grille).

\begin{enumerate}
\item \textbf{Lecture et compréhension.} Lisez les trois messages à voix haute. Soulignez les mots du lexique de la leçon que vous reconnaissez et relevez les mots nouveaux dans le glossaire.
\item \textbf{Analyse avec la grille en 4 étapes.} Pour chaque message, remplissez la grille : ¿Cuál es la \esp{fuente}? ¿Hay \esp{confirmación} en un medio serio? ¿Cuál es la \esp{fecha}? ¿La \esp{imagen} o el texto parecen manipulados? Notez vos réponses dans un tableau.
\item \textbf{Hechos u opiniones.} Dans chaque message, classez les phrases en deux colonnes : \esp{hechos} (vérifiables) et \esp{opiniones} (jugements personnels). Attention aux signaux : majuscules exagérées, points d'exclamation, \esp{¡urgente!}, \esp{comparte}, \esp{un secreto}.
\item \textbf{Veredicto.} Désignez le message vrai et les deux \esp{bulos}. Justifiez chaque décision en une phrase avec \esp{porque}.
\item \textbf{Production écrite.} Rédigez un \esp{mensaje de verificación} de 5 lignes environ, prêt à être publié dans le groupe WhatsApp. Il doit : démentir le bulo le plus dangereux, donner la bonne information, utiliser \textbf{au moins 6 mots du lexique} de la leçon et \textbf{2 expressions d'opinion} (\esp{creo que}, \esp{en mi opinión}, \esp{me parece que}).
\item \textbf{Présentation orale.} Le \esp{portavoz} présente le verdict du groupe en une minute. La classe écoute et pose une question.
\item \textbf{Vote final.} La classe vote pour le meilleur \emph{caza-bulos} selon trois critères : exactitude de l'analyse, richesse du lexique, correction de l'espagnol.
\end{enumerate}

\courssub{Ressources à mobiliser}

\begin{aretenir}[La grille de vérification en 4 étapes]
\begin{enumerate}
\item \textbf{Fuente} : ¿quién escribe? ¿es un periodista, un medio conocido, una persona anónima?
\item \textbf{Confirmación} : ¿otros medios serios dan la misma noticia?
\item \textbf{Fecha} : ¿el mensaje es reciente o antiguo?
\item \textbf{Imagen} : ¿la foto corresponde al texto o está manipulada?
\end{enumerate}
Un message sans fuente clara, sin confirmación y con muchas exclamaciones es probablemente un \esp{bulo}.
\end{aretenir}

\begin{propriete}[Expresar la opinión : indicatif après « creo que »]
Après \esp{creo que}, \esp{me parece que}, \esp{en mi opinión}, on emploie l'\textbf{indicatif} quand la phrase est affirmative : \esp{Creo que este mensaje \textbf{es} un bulo.} « Je pense que ce message est une fausse information. » Au négatif, l'espagnol exige le subjonctif : \esp{No creo que sea verdad.} « Je ne crois pas que ce soit vrai. » À la Première, retiens surtout la forme affirmative.
\end{propriete}

\begin{attention}[Pièges fréquents des francophones]
\begin{itemize}
\item \esp{estar informado} signifie « être informé », pas « être en forme ».
\item On dit \esp{una noticia falsa} ou \esp{un bulo}, jamais \esp{una noticia falsa nueva}.
\item \esp{actualmente} veut dire « actuellement », mais \esp{de verdad} = « vraiment » ; ne confonds pas \esp{verdad} (vérité) et \esp{verde} (vert) !
\item En français on dit « je pense que c'est faux » ; en espagnol : \esp{me parece que es falso} — n'oublie pas le \esp{que}.
\end{itemize}
\end{attention}

\courssub{Proposition de production et corrigé commenté}

\begin{exoresolu}[Mensaje de verificación — modèle corrigé]
Énoncé : rédigez le message de vérification du groupe (5 lignes, 6 mots du lexique, 2 expressions d'opinion) pour démentir le bulo sur les examens.
\tcblower
\textbf{Mensaje de verificación (modèle) :}\par
\esp{Atención, compañeros: el mensaje sobre los exámenes es un \textbf{bulo}. No tiene \textbf{fuente} seria y ningún \textbf{periodista} de la \textbf{prensa} confirma esta \textbf{noticia}. El Ministerio no ha publicado ningún \textbf{titular} sobre este tema en su página oficial. En mi opinión, debemos verificar siempre las \textbf{redes sociales} antes de compartir. Creo que es mejor estudiar tranquilos: los exámenes siguen en el calendario oficial.}\par\medskip
« Attention, camarades : le message sur les examens est une fausse information. Il n'a pas de source sérieuse et aucun journaliste de la presse ne confirme cette nouvelle. Le Ministère n'a publié aucun titre sur ce sujet sur sa page officielle. À mon avis, nous devons toujours vérifier les réseaux sociaux avant de partager. Je pense qu'il vaut mieux étudier tranquillement : les examens restent au calendrier officiel. »\par\medskip
\textbf{Commentaire des choix de langue :}
\begin{itemize}
\item \textbf{Lexique} : 7 mots de la leçon sont employés (\esp{bulo}, \esp{fuente}, \esp{periodista}, \esp{prensa}, \esp{noticia}, \esp{titular}, \esp{redes sociales}) : la consigne des 6 mots est dépassée.
\item \textbf{Opinion} : deux structures correctes à l'indicatif — \esp{En mi opinión, debemos verificar...} et \esp{Creo que es mejor...}.
\item \textbf{Grammaire} : \esp{no ha publicado} (passé composé avec \esp{haber}, correct pour une action récente) ; \esp{antes de compartir} (infinitif après \esp{antes de}, sans subjonctif car le sujet ne change pas).
\item \textbf{Stratégie} : le message commence par une alerte (\esp{Atención}), donne les preuves (pas de fuente, pas de confirmación), puis la bonne information, et finit par un conseil : c'est la structure idéale d'un démenti poli et efficace.
\end{itemize}
\end{exoresolu}

\courssub{Bilan de l'activité}

À l'issue de ce jeu de \emph{caza-bulos}, tu as appris à ne plus croire un message simplement parce qu'il circule beaucoup : un message \esp{reenviado 47 veces} n'est pas plus vrai qu'un message envoyé une seule fois. Tu sais désormais appliquer une méthode simple — fuente, confirmación, fecha, imagen — pour vérifier une \esp{noticia}, distinguer un \esp{hecho} d'une \esp{opinión}, et réagir en espagnol avec un message clair qui dément le \esp{bulo} et rétablit les faits. Ces compétences te serviront tous les jours, sur WhatsApp comme dans la lecture de la presse, et elles seront réinvesties dans l'évaluation du module : lire un article simple, repérer une fausse information et exprimer ton opinion en espagnol correct. Retiens la devise du bon lecteur : \esp{Piensa antes de compartir.} « Réfléchis avant de partager. »

\newpage

\coursec{Méthodes et exercices résolus}

\courssub{Méthode 1 : Lire et comprendre un article de presse simple}

\begin{methode}[Comprendre un artículo de prensa]
\begin{enumerate}
\item \textbf{Lire le titre et le sous-titre} (\esp{titular} et \esp{subtítulo}) : ils annoncent le thème et souvent la position du journaliste.
\item \textbf{Identifier le genre} : est-ce une \esp{noticia} (information factuelle), un \esp{artículo de opinión} (opinion), un \esp{reportaje} (reportage) ?
\item \textbf{Repérer les 5 questions-clés} : \esp{¿qué?} (quoi), \esp{¿quién?} (qui), \esp{¿dónde?} (où), \esp{¿cuándo?} (quand), \esp{¿por qué?} (pourquoi). Un bon article répond à ces questions dès le premier paragraphe.
\item \textbf{Relever les marqueurs de source} : \esp{según} (selon), \esp{declaró} (a déclaré), \esp{anunció} (a annoncé). Ils indiquent d'où vient l'information.
\item \textbf{Deviner le vocabulaire inconnu par le contexte} avant de consulter le glossaire : famille de mots, mots apparentés au français (\esp{información}, \esp{comunicación}).
\item \textbf{Résumer en une phrase} : \esp{El artículo trata de...} (L'article traite de...).
\end{enumerate}
\textbf{Structures à retenir :} \esp{El texto trata de...}, \esp{Según el periodista...}, \esp{La noticia fue publicada en...}, \esp{El tema principal es...}
\end{methode}

\begin{exoresolu}[Compréhension d'un article]
\textbf{Énoncé.} Lis le texte suivant, puis réponds en espagnol par des phrases complètes.

\esp{Un grupo de alumnos de un instituto de Yaoundé creó un club de prensa para luchar contra los bulos en las redes sociales. Cada semana, los jóvenes verifican las noticias que circulan por WhatsApp antes de compartirlas. Según su profesora de español, la iniciativa ya ha ayudado a muchas familias a distinguir los hechos de las opiniones.}

\textbf{Questions :} 1) \esp{¿Quiénes crearon el club?} 2) \esp{¿Qué hacen los jóvenes cada semana?} 3) \esp{¿Cuál es la opinión de la profesora?}
\tcblower
\textbf{Raisonnement.} On applique la méthode : question 1 = \esp{¿quién?}, question 2 = \esp{¿qué?}, question 3 = recherche du marqueur d'opinion \esp{según}.

\textbf{Réponses rédigées :}

1) \esp{Un grupo de alumnos de un instituto de Yaoundé creó el club de prensa.} — On reprend le sujet exact du texte ; attention à l'accord : \esp{un grupo de alumnos} (singulier + pluriel).

2) \esp{Cada semana, los jóvenes verifican las noticias que circulan por WhatsApp antes de compartirlas.} — \esp{compartirlas} : le pronom \esp{las} (les = \esp{las noticias}, féminin pluriel) se colle à l'infinitif.

3) \esp{Según la profesora, la iniciativa ya ha ayudado a muchas familias a distinguir los hechos de las opiniones.} — On introduit l'opinion avec \esp{según} ; le présent perfect \esp{ha ayudado} exprime un bilan d'une action récente.

\textbf{Conclusion.} Une bonne réponse de compréhension reprend les mots du texte, les organise en phrase complète et respecte les accords.
\end{exoresolu}

\courssub{Méthode 2 : Distinguer hecho y opinión}

\begin{methode}[Distinguer un fait d'une opinion]
\begin{enumerate}
\item \textbf{Un fait} (\esp{hecho}) peut être vérifié : chiffres, dates, événements, citations entre guillemets. Exemple : \esp{El club se creó en enero.} (Le club a été créé en janvier.)
\item \textbf{Une opinion} (\esp{opinión}) exprime un jugement, un sentiment, une prévision : adjectifs évaluatifs (\esp{bueno, terrible, importante}), verbes d'opinion (\esp{creer, parecer, pensar}), conditionnel de rumeur.
\item \textbf{Repérer les marqueurs d'opinion} : \esp{creo que}, \esp{en mi opinión}, \esp{me parece que}, \esp{pienso que}, \esp{es evidente que}, \esp{desgraciadamente}.
\item \textbf{Attention au mélange} : un journaliste peut glisser une opinion dans une nouvelle. Souligne les adjectifs et les verbes de jugement.
\item \textbf{Vérifier la source} (\esp{fuente}) : qui parle ? Est-ce un témoin, un expert, un anonyme ? Une information sans source est suspecte.
\end{enumerate}
\textbf{Réflexe :} demande-toi toujours \esp{¿Se puede comprobar?} (Peut-on le vérifier ?). Si oui, c'est un fait ; si non, c'est une opinion.
\end{methode}

\begin{exoresolu}[Hecho u opinión]
\textbf{Énoncé.} Indique si chaque phrase est un \esp{hecho} (H) ou una \esp{opinión} (O) et justifie ta réponse en français.

1) \esp{El periódico publicó la noticia el lunes.} 2) \esp{Creo que las redes sociales son muy peligrosas.} 3) \esp{En mi opinión, el club de prensa es una idea excelente.} 4) \esp{Más de cien alumnos participaron en la reunión.}
\tcblower
\textbf{Raisonnement.} On applique le test : « Peut-on vérifier ? »

1) \textbf{H.} \esp{El periódico publicó la noticia el lunes.} — Date précise, événement vérifiable : c'est un fait.

2) \textbf{O.} Marqueur explicite \esp{creo que} + adjectif évaluatif \esp{peligrosas} : jugement personnel, invérifiable.

3) \textbf{O.} Marqueur \esp{en mi opinión} + adjectif évaluatif \esp{excelente} : opinion clairement signalée.

4) \textbf{H.} \esp{Más de cien alumnos participaron} : chiffre et événement vérifiables, malgré l'approximation \esp{más de}.

\textbf{Conclusion.} Les marqueurs \esp{creo que} et \esp{en mi opinión} trahissent toujours une opinion ; les dates et les chiffres signalent généralement un fait.
\end{exoresolu}

\courssub{Méthode 3 : Repérer une fausse information (bulo)}

\begin{methode}[Detectar un bulo en las redes sociales]
\begin{enumerate}
\item \textbf{Lire au-delà du titre} : un \esp{titular} sensationnaliste (\esp{¡Increíble!}, \esp{¡Urgente!}) cache souvent un article vide.
\item \textbf{Vérifier la source} : le site est-il connu ? L'adresse est-elle correcte ? Un \esp{bulo} vient souvent d'un compte anonyme.
\item \textbf{Croiser l'information} : chercher la même nouvelle dans plusieurs \esp{medios} sérieux. Si un seul compte en parle, méfiance.
\item \textbf{Examiner les images} : une photo peut être ancienne ou prise dans un autre pays.
\item \textbf{Repérer les signaux} : fautes d'orthographe, MAJUSCULES partout, absence d'auteur, appel à partager vite (\esp{¡Comparte antes de que lo borren!}).
\item \textbf{Ne pas partager avant de vérifier} : \esp{Antes de compartir, verifica.}
\end{enumerate}
\end{methode}

\begin{exoresolu}[Analyse d'un message suspect]
\textbf{Énoncé.} Voici un message reçu sur WhatsApp :

\esp{¡¡URGENTE!! El Gobierno regala 50.000 francos a todos los estudiantes que compartan este mensaje con 10 grupos. ¡No es un bulo! ¡Comparte YA!}

Relève trois signaux qui montrent qu'il s'agit probablement d'un \esp{bulo} et explique-les en français.
\tcblower
\textbf{Raisonnement.} On applique la grille de détection.

\textbf{Signaux relevés :}

1) \textbf{Le ton sensationnaliste} : \esp{¡¡URGENTE!!}, \esp{¡Comparte YA!} — les majuscules et les doubles points d'exclamation cherchent à provoquer une réaction émotionnelle rapide, sans réflexion.

2) \textbf{L'absence de source} : \esp{el Gobierno} est vague (quel ministère ? quel site officiel ?). Aucune date, aucun nom, aucun lien vérifiable.

3) \textbf{La condition de partage} : \esp{compartan este mensaje con 10 grupos} — une vraie aide financière ne dépend jamais du partage d'un message. C'est le mécanisme typique de propagation d'un \esp{bulo}.

4) \textbf{La négation anticipée} : \esp{¡No es un bulo!} — paradoxalement, affirmer qu'un message n'est pas un canular est un signal classique de canular.

\textbf{Conclusion.} Ce message cumule presque tous les signaux du \esp{bulo} : il ne faut ni le croire ni le partager. \esp{Antes de compartir, verifica.}
\end{exoresolu}

\courssub{Méthode 4 : Exprimer son opinion à l'écrit}

\begin{methode}[Expresar la opinión : creo que, en mi opinión, me parece que]
\begin{enumerate}
\item \textbf{Choisir un marqueur} : \esp{creo que} (je crois que), \esp{pienso que} (je pense que), \esp{en mi opinión} (à mon avis), \esp{me parece que} (il me semble que), \esp{estoy de acuerdo con} (je suis d'accord avec), \esp{no estoy de acuerdo con} (je ne suis pas d'accord avec).
\item \textbf{Construire la phrase} : marqueur + \esp{que} + phrase à l'\textbf{indicatif} pour affirmer : \esp{Creo que la prensa es importante.} (Je crois que la presse est importante.)
\item \textbf{Justifier avec} \esp{porque} : \esp{...porque informa a los ciudadanos.} (...parce qu'elle informe les citoyens.)
\item \textbf{Nuance possible} : \esp{Es verdad que..., pero...} (Il est vrai que..., mais...).
\item \textbf{Conclure} : \esp{En conclusión}, \esp{Por eso} (C'est pourquoi).
\end{enumerate}
\textbf{Attention :} après la négation, le subjonctif s'impose : \esp{No creo que sea verdad.} (Je ne crois pas que ce soit vrai.)
\end{methode}

\begin{exoresolu}[Rédaction : dar la opinión]
\textbf{Énoncé.} Rédige un court paragraphe (5 à 6 phrases) en espagnol : \esp{¿Son útiles las redes sociales para informarse? Da tu opinión y justifícala.}
\tcblower
\textbf{Raisonnement.} Plan : marqueur d'opinion + thèse → justification avec \esp{porque} → nuance (\esp{es verdad que... pero}) → conclusion.

\textbf{Production modèle :}

\esp{En mi opinión, las redes sociales son útiles para informarse, pero hay que tener cuidado. Creo que son muy rápidas porque las noticias llegan enseguida a nuestros teléfonos. Es verdad que muchos jóvenes se informan solo por WhatsApp o Facebook, pero allí también circulan muchos bulos. Me parece que debemos verificar siempre las fuentes antes de compartir una noticia. Por eso, pienso que las redes son una buena herramienta si las usamos con espíritu crítico.}

« À mon avis, les réseaux sociaux sont utiles pour s'informer, mais il faut faire attention. Je crois qu'ils sont très rapides parce que les nouvelles arrivent aussitôt sur nos téléphones. Il est vrai que beaucoup de jeunes s'informent seulement par WhatsApp ou Facebook, mais là circulent aussi beaucoup de canulars. Il me semble que nous devons toujours vérifier les sources avant de partager une nouvelle. C'est pourquoi je pense que les réseaux sont un bon outil si nous les utilisons avec esprit critique. »

\textbf{Justifications grammaticales :} indicatif après \esp{creo que} et \esp{me parece que} (\esp{son}, \esp{debemos}) ; \esp{hay que} + infinitif pour l'obligation générale ; \esp{antes de} + infinitif ; accord de \esp{muchas} avec \esp{noticias/bulos} selon le genre.

\textbf{Conclusion.} Un paragraphe d'opinion efficace = marqueur + thèse + justification + nuance + conclusion.
\end{exoresolu}

\courssub{Méthode 5 : Traduire un texte sur les médias (thème et version)}

\begin{methode}[Traduire sans calquer le français]
\begin{enumerate}
\item \textbf{Lire toute la phrase} avant de traduire : repérer le verbe, son temps, son sujet.
\item \textbf{Traduire les idées, pas les mots} : « se tenir informé » = \esp{estar informado} ; « une fausse nouvelle » = \esp{un bulo} ou \esp{una noticia falsa} (jamais \esp{un falso amigo}... qui signifie « faux ami » !).
\item \textbf{Surveiller les temps} : passé composé français récent = \esp{pretérito perfecto} (\esp{ha publicado}) ; événement daté et achevé = \esp{pretérito indefinido} (\esp{publicó}).
\item \textbf{Vérifier les accords} : genre et nombre des adjectifs (\esp{las noticias falsas}).
\item \textbf{Placer les pronoms} : devant le verbe conjugué (\esp{la leemos}), collés à l'infinitif (\esp{antes de compartirla}).
\item \textbf{Relire} : accents, \esp{ñ}, points d'interrogation et d'exclamation ouvrants ¿ ¡.
\end{enumerate}
\end{methode}

\begin{exoresolu}[Thème : phrases sur les médias]
\textbf{Énoncé.} Traduis en espagnol :

1) Les journalistes vérifient toujours leurs sources. 2) Hier, le journal a publié une nouvelle importante. 3) Je crois que cette information est fausse. 4) Avant de partager le message, nous devons vérifier s'il est vrai.
\tcblower
\textbf{Solutions détaillées :}

1) \esp{Los periodistas verifican siempre sus fuentes.} — Présent d'habitude ; \esp{fuente} est féminin : \esp{sus fuentes}. Ne pas confondre avec \esp{fuente} = « fontaine » aussi, ici « source ».

2) \esp{Ayer, el periódico publicó una noticia importante.} — \esp{Ayer} impose le \esp{pretérito indefinido} : \esp{publicó} (avec accent sur la 3\up{e} personne). « Journal » = \esp{periódico} (et non \esp{jornal}, qui existe mais signifie « salaire journalier »).

3) \esp{Creo que esta información es falsa.} — Indicatif après \esp{creo que} ; \esp{información} est féminin singulier : \esp{esta información... falsa}. Faux ami : \esp{información} est singulier en espagnol alors que « informations » est souvent pluriel en français.

4) \esp{Antes de compartir el mensaje, debemos verificar si es verdad.} — \esp{Antes de} + infinitif ; « si » (whether) = \esp{si} sans accent ; \esp{es verdad} : on utilise \esp{ser} car être vrai est une qualité intrinsèque de l'information.

\textbf{Conclusion.} En thème, chaque choix (temps, genre, ser/estar, accents) doit être justifié mentalement avant d'écrire.
\end{exoresolu}

\begin{attention}[Les 5 erreurs qui coûtent des points au Bac]
\begin{enumerate}
\item \textbf{Oublier les accents et les signes espagnols} : \esp{opinión}, \esp{información}, \esp{periódico}, \esp{publicó} prennent un accent ; une question commence par ¿ et une exclamation par ¡. Écrire \esp{opinion} sans accent est une faute.
\item \textbf{Les faux amis du lexique des médias} : \esp{una noticia} = une nouvelle/une information (pas « une notice ») ; \esp{un bulo} = un canular ; \esp{actualmente} = actuellement/en ce moment (pas « actuellement » au sens de « en fait ») ; \esp{información} est singulier : \esp{una información}, jamais \esp{unas informaciones} pour « des informations ».
\item \textbf{Le calque du français « je suis d'accord »} : on dit \esp{estoy de acuerdo} (avec \esp{estar}), jamais \esp{soy de acuerdo}. De même, « avoir raison » = \esp{tener razón} (pas \esp{ser razón}).
\item \textbf{L'indicatif/subjonctif après les verbes d'opinion} : \esp{Creo que es verdad} (indicatif, affirmation) mais \esp{No creo que sea verdad} (subjonctif après la négation). Mélanger les deux modes est sanctionné.
\item \textbf{Le choix du passé} : avec \esp{ayer}, \esp{el lunes}, une date précise, utilise le \esp{pretérito indefinido} (\esp{publicó, creó, dijo}) ; le \esp{pretérito perfecto} (\esp{ha publicado}) s'emploie pour un passé récent lié au présent (\esp{esta semana, hoy, ya}). Ne traduis jamais mécaniquement le passé composé français.
\end{enumerate}
\end{attention}

\newpage

\coursec{Exercices}

\courssub{Je vérifie mes connaissances}

\begin{exercice}[QCM : el vocabulario de los medios]
Entoure la bonne réponse pour chaque question.
\begin{enumerate}
\item Un \esp{bulo} es : \quad a) una noticia verdadera \quad b) una noticia falsa \quad c) un titular largo
\item La persona que escribe artículos en un periódico es : \quad a) el periodista \quad b) el locutor \quad c) el lector
\item Un \esp{titular} es : \quad a) la foto de un artículo \quad b) el título de una noticia \quad c) la firma del autor
\item Una \esp{fuente} es : \quad a) el origen de una información \quad b) una red social \quad c) una opinión personal
\end{enumerate}
\end{exercice}

\begin{exercice}[Vrai ou faux ? Justifie ta réponse]
Réponds par V (verdadero) ou F (falso) et justifie chaque réponse par une courte phrase en espagnol.
\begin{enumerate}
\item \esp{Las redes sociales siempre publican noticias verdaderas.}
\item \esp{Un hecho se puede verificar ; una opinión expresa un sentimiento personal.}
\item \esp{Antes de compartir una noticia, hay que comprobar la fuente.}
\item \esp{« Me parece que » sirve para expresar un hecho objetivo.}
\end{enumerate}
\end{exercice}

\begin{exercice}[Complète avec le lexique]
Complète les phrases avec les mots suivants : \esp{periodista / fuente / noticia / titular / prensa}.
\begin{enumerate}
\item El \_\_\_\_\_\_ verificó la \_\_\_\_\_\_ antes de publicar la \_\_\_\_\_\_.
\item El \_\_\_\_\_\_ del artículo es muy llamativo : « ¡Escándalo en el mercado de Yaoundé ! »
\item La \_\_\_\_\_\_ escrita informa cada mañana a los lectores de Douala.
\end{enumerate}
\end{exercice}

\begin{exercice}[Conjugaison : l'expression de l'opinion]
Complète les phrases avec la forme correcte du verbe entre parenthèses au présent de l'indicatif.
\begin{enumerate}
\item Yo \_\_\_\_\_\_ (creer) que las redes sociales son útiles.
\item En mi opinión, los jóvenes \_\_\_\_\_\_ (leer) demasiados titulares sin verificar.
\item ¿Tú \_\_\_\_\_\_ (pensar) que este titular es un bulo ?
\item Nosotros \_\_\_\_\_\_ (preferir) la prensa escrita a las redes sociales.
\item Me parece que ellos no \_\_\_\_\_\_ (decir) la verdad.
\end{enumerate}
\end{exercice}

\courssub{Je m'entraîne}

\begin{exercice}[Traduction : thème]
Traduis les phrases suivantes en espagnol.
\begin{enumerate}
\item Je crois que les réseaux sociaux sont dangereux pour l'information.
\item Le journaliste a vérifié la source avant de publier la nouvelle.
\item À mon avis, ce titre est un mensonge (bulo).
\item Il faut distinguer les faits et les opinions quand on lit la presse.
\end{enumerate}
\end{exercice}

\begin{exercice}[Reformulation : exprimer l'opinion]
Transforme chaque phrase en commençant par \esp{Me parece que...} ou \esp{En mi opinión,...}.
\begin{enumerate}
\item \esp{Las redes sociales difunden muchos bulos.}
\item \esp{La prensa escrita es más seria que Internet.}
\item \esp{Los jóvenes cameruneses usan demasiado WhatsApp.}
\item \esp{Este periodista trabaja muy bien.}
\item \esp{Verificar las fuentes es importante.}
\end{enumerate}
\end{exercice}

\begin{exercice}[Texte à trous : un artículo]
Complète le texte avec les mots suivants : \esp{bulo / noticias / hechos / fuentes / redes sociales / opinión}.

\begin{textebox}[Un mensaje de WhatsApp]
Ayer recibí en mi teléfono un mensaje que circulaba por las \_\_\_\_\_\_ : « ¡El lunes no hay clases en todo el país ! ». Mi hermano mayor, que es estudiante de periodismo, me dijo : « Cuidado, eso puede ser un \_\_\_\_\_\_ ». Juntos buscamos la información en la página oficial del Ministerio : no había ningún comunicado. Comprendí que hay que separar los \_\_\_\_\_\_ de la simple \_\_\_\_\_\_ de la gente, y comprobar siempre las \_\_\_\_\_\_ antes de creer las \_\_\_\_\_\_ que recibimos.
\end{textebox}
\end{exercice}

\begin{exercice}[Appariements : palabras y definiciones]
Associe chaque mot (1-6) à sa définition (a-f).
\begin{enumerate}
\item \esp{el titular}
\item \esp{el bulo}
\item \esp{la fuente}
\item \esp{el hecho}
\item \esp{la opinión}
\item \esp{el periodista}
\end{enumerate}
\begin{itemize}
\item[a)] \esp{Lo que una persona piensa o siente.}
\item[b)] \esp{El profesional que busca y escribe noticias.}
\item[c)] \esp{El título de un artículo de prensa.}
\item[d)] \esp{Una información falsa que circula por Internet.}
\item[e)] \esp{Algo real, que se puede comprobar.}
\item[f)] \esp{El origen de donde viene una información.}
\end{itemize}
\end{exercice}

\courssub{Je me prépare au Bac}

\begin{exercice}[Compréhension écrite type Bac]
Lis le texte suivant, puis réponds aux questions \textbf{en espagnol}, avec des phrases complètes.

\begin{textebox}[Los jóvenes y la información en la era digital]
Hoy en día, muchos jóvenes africanos se informan principalmente a través de las redes sociales como WhatsApp, Facebook o TikTok. En las ciudades de Camerún, como Yaoundé o Douala, es raro ver a un adolescente leer un periódico de papel : la mayoría prefiere mirar su teléfono móvil.

Esta situación tiene ventajas : las noticias llegan muy rápido y todo el mundo puede compartir información. Sin embargo, también tiene peligros. Según los profesores de secundaria, muchos estudiantes creen todo lo que leen sin verificar las fuentes. Los bulos circulan de grupo en grupo y, en pocos minutos, miles de personas comparten una noticia falsa.

« Antes de compartir, hay que pensar », explica un periodista camerunés. « Un buen lector debe hacerse tres preguntas : ¿quién escribe esta noticia ?, ¿dónde se publicó ?, ¿otros medios serios hablan de lo mismo ? ».

Algunas escuelas ya organizan talleres para enseñar a los alumnos a distinguir los hechos de las opiniones. Para muchos educadores, formar lectores críticos es hoy tan importante como enseñar a leer y escribir. La información es un poder : aprender a usarla con responsabilidad es un deber de todos los ciudadanos.
\end{textebox}

\textbf{Glossaire :} \esp{sin embargo} = cependant ; \esp{talleres} = ateliers ; \esp{deber} = devoir.

\textbf{Questions de comprensión :}
\begin{enumerate}
\item ¿Cómo se informan los jóvenes cameruneses según el texto ?
\item ¿Cuáles son las ventajas de las redes sociales ?
\item ¿Qué hacen muchos estudiantes con las noticias que leen ?
\item ¿Qué tres preguntas propone el periodista ?
\item ¿Qué hacen algunas escuelas para ayudar a los alumnos ?
\end{enumerate}

\textbf{Vrai ou faux :} réponds V ou F et justifie avec une phrase du texte.
\begin{enumerate}
\item \esp{La mayoría de los adolescentes lee periódicos de papel.}
\item \esp{Los bulos se difunden muy rápidamente.}
\end{enumerate}

\textbf{Questions de langue :}
\begin{enumerate}
\item Trouve dans le texte deux mots de la famille de \esp{información}.
\item Mets à la forme négative : \esp{Muchos estudiantes verifican las fuentes.}
\end{enumerate}
\end{exercice}

\begin{exercice}[Thème et version type Bac]
\textbf{A. Thème.} Traduis en espagnol.
\begin{enumerate}
\item Les réseaux sociaux transmettent les nouvelles très rapidement.
\item À mon avis, il faut toujours vérifier la source d'une information.
\item Le journaliste a publié un article sur les mensonges (bulos) d'Internet.
\item Beaucoup de jeunes croient tout ce qu'ils lisent sur leur téléphone.
\end{enumerate}

\textbf{B. Version.} Traduis en français.

\begin{textebox}[Un consejo útil]
Mi abuelo lee el periódico todas las mañanas en el patio de la casa. Ayer me dijo : « No creas todas las noticias que recibes por WhatsApp. Un hecho se puede comprobar ; una opinión es solo lo que piensa una persona. Si aprendes a distinguirlos, serás un ciudadano libre ». Creo que mi abuelo tiene razón : pensar antes de compartir es una forma de respetar la verdad.
\end{textebox}
\end{exercice}

\begin{exercice}[Expression écrite type Bac]
\textbf{Sujet :} \esp{¿Crees que las redes sociales son peligrosas para la información ?}

Rédige un texte de \textbf{80 à 100 mots} en espagnol. Tu donneras ton avis personnel et tu le justifieras avec au moins deux arguments et un exemple.

\textbf{Consignes obligatoires :}
\begin{itemize}
\item Utilise au moins une fois chacune de ces expressions : \esp{creo que}, \esp{en mi opinión}, \esp{me parece que}.
\item Utilise au moins quatre mots du lexique de la leçon : \esp{prensa, noticia, periodista, titular, bulo, fuente, hecho, opinión}.
\item Structure ton texte : introduction (ton avis), développement (arguments et exemple), conclusion (une phrase de bilan).
\end{itemize}
\end{exercice}

\newpage

\coursec{Corrigés des exercices}

\begin{corrige}[Exercice 1 — QCM : el vocabulario de los medios]
\begin{enumerate}
\item \textbf{b) una noticia falsa.} Un \esp{bulo} est une fausse nouvelle qui circule, surtout sur Internet et les réseaux sociaux.
\item \textbf{a) el periodista.} C'est le professionnel qui recherche et écrit les articles. Le \esp{locutor} parle à la radio ou à la télévision ; le \esp{lector} lit.
\item \textbf{b) el título de una noticia.} Le \esp{titular} est le titre d'un article, souvent écrit en gros caractères pour attirer le lecteur.
\item \textbf{a) el origen de una información.} La \esp{fuente} est la source : la personne, le document ou le site d'où vient l'information.
\end{enumerate}
\textbf{Point de vigilance :} ne pas confondre \esp{noticia} (une nouvelle, une information) et \esp{nota} (une note scolaire) : faux ami fréquent.
\end{corrige}

\begin{corrige}[Exercice 2 — Vrai ou faux ?]
\begin{enumerate}
\item \textbf{F.} \esp{Las redes sociales también difunden bulos y noticias falsas.} (Les réseaux sociaux diffusent aussi de fausses nouvelles.)
\item \textbf{V.} \esp{Un hecho es real y se puede comprobar ; una opinión es lo que piensa una persona.} C'est la définition même de la distinction fait/opinion.
\item \textbf{V.} \esp{Hay que comprobar la fuente para saber si la noticia es verdadera.} Vérifier la source est la première règle de l'esprit critique.
\item \textbf{F.} \esp{« Me parece que » sirve para expresar una opinión personal, no un hecho.} Les verbes d'opinion (\esp{creer, parecer, pensar}) introduisent un point de vue subjectif.
\end{enumerate}
\end{corrige}

\begin{corrige}[Exercice 3 — Complète avec le lexique]
\begin{enumerate}
\item El \textbf{periodista} verificó la \textbf{fuente} antes de publicar la \textbf{noticia}.\\
« Le journaliste a vérifié la source avant de publier la nouvelle. »
\item El \textbf{titular} del artículo es muy llamativo : « ¡Escándalo en el mercado de Yaoundé ! »\\
« Le titre de l'article est très accrocheur. »
\item La \textbf{prensa} escrita informa cada mañana a los lectores de Douala.\\
« La presse écrite informe chaque matin les lecteurs de Douala. »
\end{enumerate}
\textbf{Justification :} seul le \esp{periodista} peut être le sujet humain de \esp{verificó} ; la \esp{prensa escrita} est l'expression consacrée pour « la presse écrite » (journaux et magazines).
\end{corrige}

\begin{corrige}[Exercice 4 — Conjugaison : l'expression de l'opinion]
\begin{enumerate}
\item Yo \textbf{creo} que las redes sociales son útiles. — \esp{creer} est régulier au présent : creo, crees, cree...
\item En mi opinión, los jóvenes \textbf{leen} demasiados titulares sin verificar. — \esp{leer}, 3\up{e} personne du pluriel : leen.
\item ¿Tú \textbf{piensas} que este titular es un bulo ? — \esp{pensar} est un verbe à changement radical \textbf{e $\rightarrow$ ie} : pienso, piensas, piensa, pensamos, pensáis, piensan.
\item Nosotros \textbf{preferimos} la prensa escrita a las redes sociales. — \esp{preferir} : changement \textbf{e $\rightarrow$ ie}, mais \emph{pas} à la 1\iere{} et 2\up{e} personnes du pluriel : preferimos, preferís.
\item Me parece que ellos no \textbf{dicen} la verdad. — \esp{decir} : changement \textbf{e $\rightarrow$ i} : digo, dices, dice, decimos, decís, dicen.
\end{enumerate}
\textbf{Point de vigilance :} les francophones oublient souvent le changement de voyelle : on n'écrit jamais \esp{*pensas} ni \esp{*decen}.
\end{corrige}

\begin{corrige}[Exercice 5 — Traduction : thème]
\begin{enumerate}
\item \esp{Creo que las redes sociales son peligrosas para la información.}
\item \esp{El periodista verificó (ha verificado) la fuente antes de publicar la noticia.}\\
Attention : « avant de + infinitif » se traduit par \esp{antes de + infinitivo} ; le passé composé français se rend ici par le \esp{pretérito perfecto} (\esp{ha verificado}) ou l'\esp{indefinido} (\esp{verificó}).
\item \esp{En mi opinión, este titular es un bulo.}
\item \esp{Hay que distinguir los hechos y las opiniones (de las opiniones) cuando se lee la prensa.}\\
« Il faut » impersonnel = \esp{hay que + infinitivo}.
\end{enumerate}
\textbf{Points de vigilance :} ne pas traduire « dangereux » par \esp{*dangerosos} (le mot espagnol est \esp{peligrosos}) ; « l'information » au sens des médias = \esp{la información}, toujours singulier.
\end{corrige}

\begin{corrige}[Exercice 6 — Reformulation : exprimer l'opinion]
Réponses possibles (d'autres formulations sont acceptables) :
\begin{enumerate}
\item \esp{Me parece que las redes sociales difunden muchos bulos.}
\item \esp{En mi opinión, la prensa escrita es más seria que Internet.}
\item \esp{Me parece que los jóvenes cameruneses usan demasiado WhatsApp.}
\item \esp{En mi opinión, este periodista trabaja muy bien.}
\item \esp{Me parece que verificar las fuentes es importante.}
\end{enumerate}
\textbf{Règle :} après \esp{me parece que} et \esp{en mi opinión}, le verbe reste à l'\textbf{indicatif} car on affirme son point de vue. \textbf{Attention :} avec \esp{en mi opinión}, la virgule est obligatoire après l'expression ; avec \esp{me parece que}, on n'oublie pas le \esp{que}.
\end{corrige}

\begin{corrige}[Exercice 7 — Texte à trous : un artículo]
Texte complété :

\esp{Ayer recibí en mi teléfono un mensaje que circulaba por las \textbf{redes sociales} : « ¡El lunes no hay clases en todo el país ! ». Mi hermano mayor, que es estudiante de periodismo, me dijo : « Cuidado, eso puede ser un \textbf{bulo} ». Juntos buscamos la información en la página oficial del Ministerio : no había ningún comunicado. Comprendí que hay que separar los \textbf{hechos} de la simple \textbf{opinión} de la gente, y comprobar siempre las \textbf{fuentes} antes de creer las \textbf{noticias} que recibimos.}

\textbf{Traduction :} « Hier, j'ai reçu sur mon téléphone un message qui circulait sur les réseaux sociaux : "Lundi, pas de cours dans tout le pays !". Mon frère aîné, qui est étudiant en journalisme, m'a dit : "Attention, c'est peut-être une fausse nouvelle". Ensemble, nous avons cherché l'information sur la page officielle du Ministère : il n'y avait aucun communiqué. J'ai compris qu'il faut séparer les faits de la simple opinion des gens, et toujours vérifier les sources avant de croire les nouvelles que nous recevons. »

\textbf{Justification :} \esp{un bulo} est masculin singulier (article \esp{un}) ; \esp{los hechos} s'oppose logiquement à \esp{la opinión} ; on vérifie les \esp{fuentes} et l'on croit (ou non) les \esp{noticias}.
\end{corrige}

\begin{corrige}[Exercice 8 — Appariements : palabras y definiciones]
\begin{center}
\begin{tabular}{|l|l|}
\hline
\rowcolor{popblueL} \textbf{Palabra} & \textbf{Definición} \\
\hline
1. el titular & c) El título de un artículo de prensa. \\
\hline
2. el bulo & d) Una información falsa que circula por Internet. \\
\hline
3. la fuente & f) El origen de donde viene una información. \\
\hline
4. el hecho & e) Algo real, que se puede comprobar. \\
\hline
5. la opinión & a) Lo que una persona piensa o siente. \\
\hline
6. el periodista & b) El profesional que busca y escribe noticias. \\
\hline
\end{tabular}
\end{center}
\textbf{Astuce de mémorisation :} le couple \esp{hecho / opinión} fonctionne comme le couple français « fait / opinion » : le fait se vérifie, l'opinion se discute.
\end{corrige}

\begin{corrige}[Exercice 9 — Compréhension écrite type Bac]
\textbf{Barème indicatif (10 points) :} questions de compréhension : 5 $\times$ 1 pt ; vrai/faux : 2 $\times$ 1,5 pt (0,5 pour V/F + 1 pour la justification) ; questions de langue : 2 $\times$ 1 pt.

\textbf{Questions de comprensión — réponses modèles :}
\begin{enumerate}
\item \esp{Los jóvenes cameruneses se informan principalmente a través de las redes sociales como WhatsApp, Facebook o TikTok.}
\item \esp{Las ventajas son que las noticias llegan muy rápido y que todo el mundo puede compartir información.}
\item \esp{Muchos estudiantes creen todo lo que leen sin verificar las fuentes.}
\item \esp{El periodista propone tres preguntas : ¿quién escribe esta noticia ?, ¿dónde se publicó ? y ¿otros medios serios hablan de lo mismo ?}
\item \esp{Algunas escuelas organizan talleres para enseñar a los alumnos a distinguir los hechos de las opiniones.}
\end{enumerate}

\textbf{Vrai ou faux :}
\begin{enumerate}
\item \textbf{F.} \esp{« Es raro ver a un adolescente leer un periódico de papel : la mayoría prefiere mirar su teléfono móvil. »}
\item \textbf{V.} \esp{« En pocos minutos, miles de personas comparten una noticia falsa. »}
\end{enumerate}

\textbf{Questions de langue :}
\begin{enumerate}
\item Deux mots de la famille de \esp{información} : \esp{informan} (verbe) et \esp{informarse} (verbe pronominal). On accepte aussi \esp{información} elle-même reprise avec \esp{informan}.
\item Forme négative : \esp{Muchos estudiantes \textbf{no} verifican las fuentes.} — En espagnol, la négation se fait avec \esp{no} placé \emph{devant} le verbe, sans équivalent de « pas ».
\end{enumerate}

\textbf{Points de vigilance :} répondre par des phrases complètes (sujet + verbe) ; recopier la justification \emph{exactement} telle qu'elle figure dans le texte, entre guillemets ; ne pas traduire les questions en français.
\end{corrige}

\begin{corrige}[Exercice 10 — Thème et version type Bac]
\textbf{Barème indicatif (10 points) :} thème : 4 $\times$ 1,5 pt ; version : 4 pts (respect du sens, correction grammaticale, qualité du français).

\textbf{A. Thème — corrigé :}
\begin{enumerate}
\item \esp{Las redes sociales transmiten las noticias muy rápidamente.}
\item \esp{En mi opinión, hay que verificar siempre la fuente de una información.}
\item \esp{El periodista ha publicado (publicó) un artículo sobre los bulos de Internet.}
\item \esp{Muchos jóvenes creen todo lo que leen en su teléfono.}
\end{enumerate}
\textbf{Points de vigilance :} « il faut » = \esp{hay que} ; « tout ce qu'ils lisent » = \esp{todo lo que leen} (\esp{lo que} neutre, jamais \esp{*todo que}) ; « sur leur téléphone » = \esp{en su teléfono}.

\textbf{B. Version — traduction modèle :}

« Mon grand-père lit le journal tous les matins dans la cour de la maison. Hier, il m'a dit : "Ne crois pas toutes les nouvelles que tu reçois par WhatsApp. Un fait peut être vérifié ; une opinion est seulement ce que pense une personne. Si tu apprends à les distinguer, tu seras un citoyen libre." Je crois que mon grand-père a raison : réfléchir avant de partager est une façon de respecter la vérité. »

\textbf{Points de vigilance :} \esp{No creas} est un impératif négatif (subjonctif) : « ne crois pas » ; \esp{tener razón} = « avoir raison » (ne jamais traduire par « avoir raison » avec \esp{*ser razón}) ; \esp{el patio} = ici « la cour » et non « le patio ».
\end{corrige}

\begin{corrige}[Exercice 11 — Expression écrite type Bac]
\textbf{Barème indicatif (10 points) :} respect de la consigne et de la longueur (2 pts) ; structure introduction/développement/conclusion (2 pts) ; emploi des trois expressions d'opinion et de quatre mots du lexique (3 pts) ; correction grammaticale et richesse du vocabulaire (3 pts).

\textbf{Proposition de rédaction modèle (environ 95 mots) :}

\begin{textebox}[¿Son peligrosas las redes sociales ?]
En mi opinión, las redes sociales son útiles pero también peligrosas para la información.

Primero, creo que difunden las noticias muy rápidamente, y eso es una ventaja. Sin embargo, me parece que muchos usuarios comparten bulos sin verificar la fuente. Por ejemplo, el mes pasado circuló en WhatsApp un titular falso sobre el cierre de los mercados de Douala, y muchas personas lo creyeron. Además, pocos jóvenes distinguen un hecho de una simple opinión.

En conclusión, las redes sociales no son malas, pero hay que usarlas con espíritu crítico, como hace un buen periodista con la prensa.
\end{textebox}

\textbf{Traduction :} « À mon avis, les réseaux sociaux sont utiles mais aussi dangereux pour l'information. D'abord, je crois qu'ils transmettent les nouvelles très rapidement, et c'est un avantage. Cependant, il me semble que beaucoup d'utilisateurs partagent de fausses nouvelles sans vérifier la source. Par exemple, le mois dernier, un faux titre sur la fermeture des marchés de Douala a circulé sur WhatsApp, et beaucoup de personnes l'ont cru. De plus, peu de jeunes distinguent un fait d'une simple opinion. En conclusion, les réseaux sociaux ne sont pas mauvais, mais il faut les utiliser avec esprit critique, comme le fait un bon journaliste avec la presse. »

\textbf{Points de vigilance :}
\begin{itemize}
\item Les trois expressions d'opinion exigées sont bien présentes : \esp{en mi opinión}, \esp{creo que}, \esp{me parece que}, chacune suivie de l'indicatif.
\item Le lexique de la leçon est mobilisé : \esp{noticias, bulos, fuente, titular, hecho, opinión, periodista, prensa}.
\item Éviter les calques du français : « les réseaux sociaux sont dangereux » = \esp{las redes sociales son peligrosas} (accord féminin pluriel) ; « partager » = \esp{compartir}, jamais \esp{*partir}.
\item Compter ses mots : un texte trop court (moins de 80 mots) ou trop long est pénalisé.
\end{itemize}
\end{corrige}

\newpage

\coursec{Fiche bilan}

\begin{popfigure}
\begin{tikzpicture}[mindmap, grow cyclic, every node/.style={concept, align=center, font=\small}, root concept/.append style={concept color=popblue, font=\bfseries\small, minimum size=3.2cm}, level 1 concept/.append style={level distance=4.2cm, sibling angle=51.4}, text=white]
\node[root concept]{Los medios y la información responsable}
  child[concept color=poporange]{ node{Lexique : prensa, noticia, titular, periodista} }
  child[concept color=popgreen]{ node{Redes sociales : bulo, fuente, viralizar} }
  child[concept color=poppurple]{ node{Hecho vs opinión : datos vérifiables / juicios} }
  child[concept color=poppink]{ node{Opinar : creo que, en mi opinión, me parece que} }
  child[concept color=popgold]{ node{Espíritu crítico : verificar, comparar fuentes} }
  child[concept color=popblue!70]{ node{Méthode : comentario de texto} }
  child[concept color=popgreen!70]{ node{Producción : escribir y opinar} };
\end{tikzpicture}
\legende{Carte mentale de la leçon : les médias et l'information responsable}
\end{popfigure}

\begin{definition}[Lexique clé : los medios de comunicación (à connaître par cœur)]
\begin{itemize}
\item \esp{la prensa} : la presse ; \esp{el periódico / el diario} : le journal ; \esp{la revista} : le magazine.
\item \esp{la noticia} : la nouvelle, l'information ; \esp{el titular} : le titre (d'un article) ; \esp{el artículo} : l'article.
\item \esp{el/la periodista} : le/la journaliste ; \esp{la redacción} : la rédaction ; \esp{informar} : informer ; \esp{el reportaje} : le reportage.
\item \esp{las redes sociales} : les réseaux sociaux ; \esp{compartir} : partager ; \esp{viralizar} : rendre viral ; \esp{publicar} : publier.
\item \esp{el bulo} : la fausse information, le canard (fake news) ; \esp{la fuente} : la source ; \esp{verificar / comprobar} : vérifier.
\item \esp{el hecho} : le fait ; \esp{la opinión} : l'opinion ; \esp{verdadero / falso} : vrai / faux.
\end{itemize}
\end{definition}

\begin{attention}[Faux amis et pièges de traduction]
\begin{itemize}
\item \esp{la noticia} signifie « la nouvelle, l'information », et non « la notice » (une notice = \esp{las instrucciones}).
\item \esp{el bulo} signifie « la fausse information », et non « le bulletin » (un bulletin scolaire = \esp{el boletín de notas}).
\item \esp{actual} signifie « actuel » ; « actuel » au sens de vrai se dit \esp{real, verdadero}.
\item \esp{informarse} = s'informer ; attention à la construction pronominale : \esp{Me informo en Internet.} « Je m'informe sur Internet. »
\end{itemize}
\end{attention}

\begin{propriete}[Hecho u opinión : les indices à repérer]
Un \cle{hecho} (fait) est une information vérifiable : chiffres, dates, lieux, citations exactes. Une \cle{opinión} est un jugement personnel, marqué par des verbes d'opinion et des adjectifs évaluatifs.
\begin{center}
\begin{tabularx}{\linewidth}{|X|X|X|}\hline
\rowcolor{popblueL} Hecho (fait) & Opinión (opinion) & Remarque \\ \hline
\esp{El partido terminó 2-0.} & \esp{Fue el mejor partido del año.} & Le fait se vérifie ; l'opinion se discute. \\ \hline
Chiffres, dates, lieux, citations & \esp{creo que, me parece, es evidente, mejor/peor} & Adjectifs évaluatifs = opinion. \\ \hline
Verbes objectifs : \esp{decir, anunciar, ocurrir} & Verbes d'opinion : \esp{creer, pensar, opinar} & Repérer le verbe introducteur. \\ \hline
\end{tabularx}
\end{center}
\end{propriete}

\begin{propriete}[Grammaire : exprimer son opinion]
\begin{itemize}
\item \esp{Creo que + indicatif} : \esp{Creo que la noticia es falsa.} « Je pense que la nouvelle est fausse. »
\item \esp{No creo que + subjonctif} : \esp{No creo que sea verdad.} « Je ne pense pas que ce soit vrai. » (La négation du verbe d'opinion entraîne le subjonctif dans la subordonnée.)
\item \esp{En mi opinión, ...} / \esp{Me parece que + indicatif} : \esp{Me parece que el titular es exagerado.} « Il me semble que le titre est exagéré. »
\item Variantes pour nuancer : \esp{Pienso que...} « je pense que... » ; \esp{Opino que...} « j'estime que... » ; \esp{Estoy de acuerdo con...} « je suis d'accord avec... » ; \esp{No estoy de acuerdo con...} « je ne suis pas d'accord avec... »
\end{itemize}
\end{propriete}

\begin{methode}[Vérifier une information : la méthode anti-bulo en 5 étapes]
\begin{enumerate}
\item Lire au-delà du \esp{titular} : le titre peut être exagéré ou trompeur.
\item Identifier la \esp{fuente} : qui publie ? un média connu ou un compte anonyme ?
\item Comparer : chercher la même \esp{noticia} dans d'autres médias sérieux.
\item Distinguer \esp{hecho} et \esp{opinión} dans le texte.
\item Vérifier la date et les images avant de \esp{compartir}.
\end{enumerate}
\end{methode}

\begin{methode}[Commentaire de texte journalistique : la démarche en 4 étapes]
\begin{enumerate}
\item Présenter : type de texte, source, thème. Modèle : \esp{El texto es un artículo de prensa publicado en... Trata de...} « Le texte est un article de presse publié dans... Il traite de... »
\item Résumer les idées principales en 2-3 phrases, sans recopier le texte.
\item Analyser : faits rapportés, opinions exprimées, ton du journaliste (neutre, critique, alarmiste), lexique des médias.
\item Donner son avis argumenté : \esp{En mi opinión, este artículo es interesante porque...} et conclure.
\end{enumerate}
\end{methode}

\begin{exemplebox}[Modèles de phrases pour donner son avis]
\begin{itemize}
\item \esp{En mi opinión, las redes sociales informan muy rápido, pero a veces difunden bulos.} « À mon avis, les réseaux sociaux informent très vite, mais diffusent parfois de fausses informations. »
\item \esp{Creo que es importante verificar las fuentes antes de compartir una noticia.} « Je pense qu'il est important de vérifier les sources avant de partager une nouvelle. »
\item \esp{No creo que todos los titulares sean verdaderos.} « Je ne pense pas que tous les titres soient vrais. » (subjonctif après \esp{no creo que})
\item \esp{Sin embargo, la prensa seria sigue siendo necesaria. Además, los periodistas verifican los hechos.} « Cependant, la presse sérieuse reste nécessaire. De plus, les journalistes vérifient les faits. »
\end{itemize}
\end{exemplebox}

\begin{aretenir}[Checklist avant le Bac]
\begin{itemize}
\item[$\square$] Je connais le lexique des médias : \esp{prensa, noticia, titular, periodista, fuente, bulo}.
\item[$\square$] Je sais distinguer un \esp{hecho} d'une \esp{opinión} dans un article.
\item[$\square$] Je repère les marqueurs d'opinion : \esp{creo que, me parece que, en mi opinión}.
\item[$\square$] Je sais que \esp{no creo que} est suivi du subjonctif : \esp{No creo que sea verdad.}
\item[$\square$] Je sais vérifier une information : source, date, comparaison des médias.
\item[$\square$] Je ne confonds pas \esp{noticia} (une information) et « notice ».
\item[$\square$] Je sais que \esp{bulo} signifie « fausse information » et non « bulletin ».
\item[$\square$] Je sais présenter un article : type, source, thème, idées principales.
\item[$\square$] Je peux donner mon avis en 3-4 phrases avec des connecteurs (\esp{sin embargo, además, por eso}).
\item[$\square$] Je vérifie les accents : \esp{opinión, información, artículo, crítico}.
\end{itemize}
\end{aretenir}

\findecours

\end{document}
